\documentclass[11pt]{article}

    \usepackage[breakable]{tcolorbox}
    \usepackage{parskip} % Stop auto-indenting (to mimic markdown behaviour)
    

    % Basic figure setup, for now with no caption control since it's done
    % automatically by Pandoc (which extracts ![](path) syntax from Markdown).
    \usepackage{graphicx}
    % Keep aspect ratio if custom image width or height is specified
    \setkeys{Gin}{keepaspectratio}
    % Maintain compatibility with old templates. Remove in nbconvert 6.0
    \let\Oldincludegraphics\includegraphics
    % Ensure that by default, figures have no caption (until we provide a
    % proper Figure object with a Caption API and a way to capture that
    % in the conversion process - todo).
    \usepackage{caption}
    \DeclareCaptionFormat{nocaption}{}
    \captionsetup{format=nocaption,aboveskip=0pt,belowskip=0pt}

    \usepackage{float}
    \floatplacement{figure}{H} % forces figures to be placed at the correct location
    \usepackage{xcolor} % Allow colors to be defined
    \usepackage{enumerate} % Needed for markdown enumerations to work
    \usepackage{geometry} % Used to adjust the document margins
    \usepackage{amsmath} % Equations
    \usepackage{amssymb} % Equations
    \usepackage{textcomp} % defines textquotesingle
    % Hack from http://tex.stackexchange.com/a/47451/13684:
    \AtBeginDocument{%
        \def\PYZsq{\textquotesingle}% Upright quotes in Pygmentized code
    }
    \usepackage{upquote} % Upright quotes for verbatim code
    \usepackage{eurosym} % defines \euro

    \usepackage{iftex}
    \ifPDFTeX
        \usepackage[T1]{fontenc}
        \IfFileExists{alphabeta.sty}{
              \usepackage{alphabeta}
          }{
              \usepackage[mathletters]{ucs}
              \usepackage[utf8x]{inputenc}
          }
    \else
        \usepackage{fontspec}
        \usepackage{unicode-math}
    \fi

    \usepackage{fancyvrb} % verbatim replacement that allows latex
    \usepackage{grffile} % extends the file name processing of package graphics
                         % to support a larger range
    \makeatletter % fix for old versions of grffile with XeLaTeX
    \@ifpackagelater{grffile}{2019/11/01}
    {
      % Do nothing on new versions
    }
    {
      \def\Gread@@xetex#1{%
        \IfFileExists{"\Gin@base".bb}%
        {\Gread@eps{\Gin@base.bb}}%
        {\Gread@@xetex@aux#1}%
      }
    }
    \makeatother
    \usepackage[Export]{adjustbox} % Used to constrain images to a maximum size
    \adjustboxset{max size={0.9\linewidth}{0.9\paperheight}}

    % The hyperref package gives us a pdf with properly built
    % internal navigation ('pdf bookmarks' for the table of contents,
    % internal cross-reference links, web links for URLs, etc.)
    \usepackage{hyperref}
    % The default LaTeX title has an obnoxious amount of whitespace. By default,
    % titling removes some of it. It also provides customization options.
    \usepackage{titling}
    \usepackage{longtable} % longtable support required by pandoc >1.10
    \usepackage{booktabs}  % table support for pandoc > 1.12.2
    \usepackage{array}     % table support for pandoc >= 2.11.3
    \usepackage{calc}      % table minipage width calculation for pandoc >= 2.11.1
    \usepackage[inline]{enumitem} % IRkernel/repr support (it uses the enumerate* environment)
    \usepackage[normalem]{ulem} % ulem is needed to support strikethroughs (\sout)
                                % normalem makes italics be italics, not underlines
    \usepackage{soul}      % strikethrough (\st) support for pandoc >= 3.0.0
    \usepackage{mathrsfs}
    

    
    % Colors for the hyperref package
    \definecolor{urlcolor}{rgb}{0,.145,.698}
    \definecolor{linkcolor}{rgb}{.71,0.21,0.01}
    \definecolor{citecolor}{rgb}{.12,.54,.11}

    % ANSI colors
    \definecolor{ansi-black}{HTML}{3E424D}
    \definecolor{ansi-black-intense}{HTML}{282C36}
    \definecolor{ansi-red}{HTML}{E75C58}
    \definecolor{ansi-red-intense}{HTML}{B22B31}
    \definecolor{ansi-green}{HTML}{00A250}
    \definecolor{ansi-green-intense}{HTML}{007427}
    \definecolor{ansi-yellow}{HTML}{DDB62B}
    \definecolor{ansi-yellow-intense}{HTML}{B27D12}
    \definecolor{ansi-blue}{HTML}{208FFB}
    \definecolor{ansi-blue-intense}{HTML}{0065CA}
    \definecolor{ansi-magenta}{HTML}{D160C4}
    \definecolor{ansi-magenta-intense}{HTML}{A03196}
    \definecolor{ansi-cyan}{HTML}{60C6C8}
    \definecolor{ansi-cyan-intense}{HTML}{258F8F}
    \definecolor{ansi-white}{HTML}{C5C1B4}
    \definecolor{ansi-white-intense}{HTML}{A1A6B2}
    \definecolor{ansi-default-inverse-fg}{HTML}{FFFFFF}
    \definecolor{ansi-default-inverse-bg}{HTML}{000000}

    % common color for the border for error outputs.
    \definecolor{outerrorbackground}{HTML}{FFDFDF}

    % commands and environments needed by pandoc snippets
    % extracted from the output of `pandoc -s`
    \providecommand{\tightlist}{%
      \setlength{\itemsep}{0pt}\setlength{\parskip}{0pt}}
    \DefineVerbatimEnvironment{Highlighting}{Verbatim}{commandchars=\\\{\}}
    % Add ',fontsize=\small' for more characters per line
    \newenvironment{Shaded}{}{}
    \newcommand{\KeywordTok}[1]{\textcolor[rgb]{0.00,0.44,0.13}{\textbf{{#1}}}}
    \newcommand{\DataTypeTok}[1]{\textcolor[rgb]{0.56,0.13,0.00}{{#1}}}
    \newcommand{\DecValTok}[1]{\textcolor[rgb]{0.25,0.63,0.44}{{#1}}}
    \newcommand{\BaseNTok}[1]{\textcolor[rgb]{0.25,0.63,0.44}{{#1}}}
    \newcommand{\FloatTok}[1]{\textcolor[rgb]{0.25,0.63,0.44}{{#1}}}
    \newcommand{\CharTok}[1]{\textcolor[rgb]{0.25,0.44,0.63}{{#1}}}
    \newcommand{\StringTok}[1]{\textcolor[rgb]{0.25,0.44,0.63}{{#1}}}
    \newcommand{\CommentTok}[1]{\textcolor[rgb]{0.38,0.63,0.69}{\textit{{#1}}}}
    \newcommand{\OtherTok}[1]{\textcolor[rgb]{0.00,0.44,0.13}{{#1}}}
    \newcommand{\AlertTok}[1]{\textcolor[rgb]{1.00,0.00,0.00}{\textbf{{#1}}}}
    \newcommand{\FunctionTok}[1]{\textcolor[rgb]{0.02,0.16,0.49}{{#1}}}
    \newcommand{\RegionMarkerTok}[1]{{#1}}
    \newcommand{\ErrorTok}[1]{\textcolor[rgb]{1.00,0.00,0.00}{\textbf{{#1}}}}
    \newcommand{\NormalTok}[1]{{#1}}

    % Additional commands for more recent versions of Pandoc
    \newcommand{\ConstantTok}[1]{\textcolor[rgb]{0.53,0.00,0.00}{{#1}}}
    \newcommand{\SpecialCharTok}[1]{\textcolor[rgb]{0.25,0.44,0.63}{{#1}}}
    \newcommand{\VerbatimStringTok}[1]{\textcolor[rgb]{0.25,0.44,0.63}{{#1}}}
    \newcommand{\SpecialStringTok}[1]{\textcolor[rgb]{0.73,0.40,0.53}{{#1}}}
    \newcommand{\ImportTok}[1]{{#1}}
    \newcommand{\DocumentationTok}[1]{\textcolor[rgb]{0.73,0.13,0.13}{\textit{{#1}}}}
    \newcommand{\AnnotationTok}[1]{\textcolor[rgb]{0.38,0.63,0.69}{\textbf{\textit{{#1}}}}}
    \newcommand{\CommentVarTok}[1]{\textcolor[rgb]{0.38,0.63,0.69}{\textbf{\textit{{#1}}}}}
    \newcommand{\VariableTok}[1]{\textcolor[rgb]{0.10,0.09,0.49}{{#1}}}
    \newcommand{\ControlFlowTok}[1]{\textcolor[rgb]{0.00,0.44,0.13}{\textbf{{#1}}}}
    \newcommand{\OperatorTok}[1]{\textcolor[rgb]{0.40,0.40,0.40}{{#1}}}
    \newcommand{\BuiltInTok}[1]{{#1}}
    \newcommand{\ExtensionTok}[1]{{#1}}
    \newcommand{\PreprocessorTok}[1]{\textcolor[rgb]{0.74,0.48,0.00}{{#1}}}
    \newcommand{\AttributeTok}[1]{\textcolor[rgb]{0.49,0.56,0.16}{{#1}}}
    \newcommand{\InformationTok}[1]{\textcolor[rgb]{0.38,0.63,0.69}{\textbf{\textit{{#1}}}}}
    \newcommand{\WarningTok}[1]{\textcolor[rgb]{0.38,0.63,0.69}{\textbf{\textit{{#1}}}}}
    \makeatletter
    \newsavebox\pandoc@box
    \newcommand*\pandocbounded[1]{%
      \sbox\pandoc@box{#1}%
      % scaling factors for width and height
      \Gscale@div\@tempa\textheight{\dimexpr\ht\pandoc@box+\dp\pandoc@box\relax}%
      \Gscale@div\@tempb\linewidth{\wd\pandoc@box}%
      % select the smaller of both
      \ifdim\@tempb\p@<\@tempa\p@
        \let\@tempa\@tempb
      \fi
      % scaling accordingly (\@tempa < 1)
      \ifdim\@tempa\p@<\p@
        \scalebox{\@tempa}{\usebox\pandoc@box}%
      % scaling not needed, use as it is
      \else
        \usebox{\pandoc@box}%
      \fi
    }
    \makeatother

    % Define a nice break command that doesn't care if a line doesn't already
    % exist.
    \def\br{\hspace*{\fill} \\* }
    % Math Jax compatibility definitions
    \def\gt{>}
    \def\lt{<}
    \let\Oldtex\TeX
    \let\Oldlatex\LaTeX
    \renewcommand{\TeX}{\textrm{\Oldtex}}
    \renewcommand{\LaTeX}{\textrm{\Oldlatex}}
    % Document parameters
    % Document title
    \title{Tarea\_1\_Alvear\_Benjamin}
    
    
    
    
    
    
    
% Pygments definitions
\makeatletter
\def\PY@reset{\let\PY@it=\relax \let\PY@bf=\relax%
    \let\PY@ul=\relax \let\PY@tc=\relax%
    \let\PY@bc=\relax \let\PY@ff=\relax}
\def\PY@tok#1{\csname PY@tok@#1\endcsname}
\def\PY@toks#1+{\ifx\relax#1\empty\else%
    \PY@tok{#1}\expandafter\PY@toks\fi}
\def\PY@do#1{\PY@bc{\PY@tc{\PY@ul{%
    \PY@it{\PY@bf{\PY@ff{#1}}}}}}}
\def\PY#1#2{\PY@reset\PY@toks#1+\relax+\PY@do{#2}}

\@namedef{PY@tok@w}{\def\PY@tc##1{\textcolor[rgb]{0.73,0.73,0.73}{##1}}}
\@namedef{PY@tok@c}{\let\PY@it=\textit\def\PY@tc##1{\textcolor[rgb]{0.24,0.48,0.48}{##1}}}
\@namedef{PY@tok@cp}{\def\PY@tc##1{\textcolor[rgb]{0.61,0.40,0.00}{##1}}}
\@namedef{PY@tok@k}{\let\PY@bf=\textbf\def\PY@tc##1{\textcolor[rgb]{0.00,0.50,0.00}{##1}}}
\@namedef{PY@tok@kp}{\def\PY@tc##1{\textcolor[rgb]{0.00,0.50,0.00}{##1}}}
\@namedef{PY@tok@kt}{\def\PY@tc##1{\textcolor[rgb]{0.69,0.00,0.25}{##1}}}
\@namedef{PY@tok@o}{\def\PY@tc##1{\textcolor[rgb]{0.40,0.40,0.40}{##1}}}
\@namedef{PY@tok@ow}{\let\PY@bf=\textbf\def\PY@tc##1{\textcolor[rgb]{0.67,0.13,1.00}{##1}}}
\@namedef{PY@tok@nb}{\def\PY@tc##1{\textcolor[rgb]{0.00,0.50,0.00}{##1}}}
\@namedef{PY@tok@nf}{\def\PY@tc##1{\textcolor[rgb]{0.00,0.00,1.00}{##1}}}
\@namedef{PY@tok@nc}{\let\PY@bf=\textbf\def\PY@tc##1{\textcolor[rgb]{0.00,0.00,1.00}{##1}}}
\@namedef{PY@tok@nn}{\let\PY@bf=\textbf\def\PY@tc##1{\textcolor[rgb]{0.00,0.00,1.00}{##1}}}
\@namedef{PY@tok@ne}{\let\PY@bf=\textbf\def\PY@tc##1{\textcolor[rgb]{0.80,0.25,0.22}{##1}}}
\@namedef{PY@tok@nv}{\def\PY@tc##1{\textcolor[rgb]{0.10,0.09,0.49}{##1}}}
\@namedef{PY@tok@no}{\def\PY@tc##1{\textcolor[rgb]{0.53,0.00,0.00}{##1}}}
\@namedef{PY@tok@nl}{\def\PY@tc##1{\textcolor[rgb]{0.46,0.46,0.00}{##1}}}
\@namedef{PY@tok@ni}{\let\PY@bf=\textbf\def\PY@tc##1{\textcolor[rgb]{0.44,0.44,0.44}{##1}}}
\@namedef{PY@tok@na}{\def\PY@tc##1{\textcolor[rgb]{0.41,0.47,0.13}{##1}}}
\@namedef{PY@tok@nt}{\let\PY@bf=\textbf\def\PY@tc##1{\textcolor[rgb]{0.00,0.50,0.00}{##1}}}
\@namedef{PY@tok@nd}{\def\PY@tc##1{\textcolor[rgb]{0.67,0.13,1.00}{##1}}}
\@namedef{PY@tok@s}{\def\PY@tc##1{\textcolor[rgb]{0.73,0.13,0.13}{##1}}}
\@namedef{PY@tok@sd}{\let\PY@it=\textit\def\PY@tc##1{\textcolor[rgb]{0.73,0.13,0.13}{##1}}}
\@namedef{PY@tok@si}{\let\PY@bf=\textbf\def\PY@tc##1{\textcolor[rgb]{0.64,0.35,0.47}{##1}}}
\@namedef{PY@tok@se}{\let\PY@bf=\textbf\def\PY@tc##1{\textcolor[rgb]{0.67,0.36,0.12}{##1}}}
\@namedef{PY@tok@sr}{\def\PY@tc##1{\textcolor[rgb]{0.64,0.35,0.47}{##1}}}
\@namedef{PY@tok@ss}{\def\PY@tc##1{\textcolor[rgb]{0.10,0.09,0.49}{##1}}}
\@namedef{PY@tok@sx}{\def\PY@tc##1{\textcolor[rgb]{0.00,0.50,0.00}{##1}}}
\@namedef{PY@tok@m}{\def\PY@tc##1{\textcolor[rgb]{0.40,0.40,0.40}{##1}}}
\@namedef{PY@tok@gh}{\let\PY@bf=\textbf\def\PY@tc##1{\textcolor[rgb]{0.00,0.00,0.50}{##1}}}
\@namedef{PY@tok@gu}{\let\PY@bf=\textbf\def\PY@tc##1{\textcolor[rgb]{0.50,0.00,0.50}{##1}}}
\@namedef{PY@tok@gd}{\def\PY@tc##1{\textcolor[rgb]{0.63,0.00,0.00}{##1}}}
\@namedef{PY@tok@gi}{\def\PY@tc##1{\textcolor[rgb]{0.00,0.52,0.00}{##1}}}
\@namedef{PY@tok@gr}{\def\PY@tc##1{\textcolor[rgb]{0.89,0.00,0.00}{##1}}}
\@namedef{PY@tok@ge}{\let\PY@it=\textit}
\@namedef{PY@tok@gs}{\let\PY@bf=\textbf}
\@namedef{PY@tok@ges}{\let\PY@bf=\textbf\let\PY@it=\textit}
\@namedef{PY@tok@gp}{\let\PY@bf=\textbf\def\PY@tc##1{\textcolor[rgb]{0.00,0.00,0.50}{##1}}}
\@namedef{PY@tok@go}{\def\PY@tc##1{\textcolor[rgb]{0.44,0.44,0.44}{##1}}}
\@namedef{PY@tok@gt}{\def\PY@tc##1{\textcolor[rgb]{0.00,0.27,0.87}{##1}}}
\@namedef{PY@tok@err}{\def\PY@bc##1{{\setlength{\fboxsep}{\string -\fboxrule}\fcolorbox[rgb]{1.00,0.00,0.00}{1,1,1}{\strut ##1}}}}
\@namedef{PY@tok@kc}{\let\PY@bf=\textbf\def\PY@tc##1{\textcolor[rgb]{0.00,0.50,0.00}{##1}}}
\@namedef{PY@tok@kd}{\let\PY@bf=\textbf\def\PY@tc##1{\textcolor[rgb]{0.00,0.50,0.00}{##1}}}
\@namedef{PY@tok@kn}{\let\PY@bf=\textbf\def\PY@tc##1{\textcolor[rgb]{0.00,0.50,0.00}{##1}}}
\@namedef{PY@tok@kr}{\let\PY@bf=\textbf\def\PY@tc##1{\textcolor[rgb]{0.00,0.50,0.00}{##1}}}
\@namedef{PY@tok@bp}{\def\PY@tc##1{\textcolor[rgb]{0.00,0.50,0.00}{##1}}}
\@namedef{PY@tok@fm}{\def\PY@tc##1{\textcolor[rgb]{0.00,0.00,1.00}{##1}}}
\@namedef{PY@tok@vc}{\def\PY@tc##1{\textcolor[rgb]{0.10,0.09,0.49}{##1}}}
\@namedef{PY@tok@vg}{\def\PY@tc##1{\textcolor[rgb]{0.10,0.09,0.49}{##1}}}
\@namedef{PY@tok@vi}{\def\PY@tc##1{\textcolor[rgb]{0.10,0.09,0.49}{##1}}}
\@namedef{PY@tok@vm}{\def\PY@tc##1{\textcolor[rgb]{0.10,0.09,0.49}{##1}}}
\@namedef{PY@tok@sa}{\def\PY@tc##1{\textcolor[rgb]{0.73,0.13,0.13}{##1}}}
\@namedef{PY@tok@sb}{\def\PY@tc##1{\textcolor[rgb]{0.73,0.13,0.13}{##1}}}
\@namedef{PY@tok@sc}{\def\PY@tc##1{\textcolor[rgb]{0.73,0.13,0.13}{##1}}}
\@namedef{PY@tok@dl}{\def\PY@tc##1{\textcolor[rgb]{0.73,0.13,0.13}{##1}}}
\@namedef{PY@tok@s2}{\def\PY@tc##1{\textcolor[rgb]{0.73,0.13,0.13}{##1}}}
\@namedef{PY@tok@sh}{\def\PY@tc##1{\textcolor[rgb]{0.73,0.13,0.13}{##1}}}
\@namedef{PY@tok@s1}{\def\PY@tc##1{\textcolor[rgb]{0.73,0.13,0.13}{##1}}}
\@namedef{PY@tok@mb}{\def\PY@tc##1{\textcolor[rgb]{0.40,0.40,0.40}{##1}}}
\@namedef{PY@tok@mf}{\def\PY@tc##1{\textcolor[rgb]{0.40,0.40,0.40}{##1}}}
\@namedef{PY@tok@mh}{\def\PY@tc##1{\textcolor[rgb]{0.40,0.40,0.40}{##1}}}
\@namedef{PY@tok@mi}{\def\PY@tc##1{\textcolor[rgb]{0.40,0.40,0.40}{##1}}}
\@namedef{PY@tok@il}{\def\PY@tc##1{\textcolor[rgb]{0.40,0.40,0.40}{##1}}}
\@namedef{PY@tok@mo}{\def\PY@tc##1{\textcolor[rgb]{0.40,0.40,0.40}{##1}}}
\@namedef{PY@tok@ch}{\let\PY@it=\textit\def\PY@tc##1{\textcolor[rgb]{0.24,0.48,0.48}{##1}}}
\@namedef{PY@tok@cm}{\let\PY@it=\textit\def\PY@tc##1{\textcolor[rgb]{0.24,0.48,0.48}{##1}}}
\@namedef{PY@tok@cpf}{\let\PY@it=\textit\def\PY@tc##1{\textcolor[rgb]{0.24,0.48,0.48}{##1}}}
\@namedef{PY@tok@c1}{\let\PY@it=\textit\def\PY@tc##1{\textcolor[rgb]{0.24,0.48,0.48}{##1}}}
\@namedef{PY@tok@cs}{\let\PY@it=\textit\def\PY@tc##1{\textcolor[rgb]{0.24,0.48,0.48}{##1}}}

\def\PYZbs{\char`\\}
\def\PYZus{\char`\_}
\def\PYZob{\char`\{}
\def\PYZcb{\char`\}}
\def\PYZca{\char`\^}
\def\PYZam{\char`\&}
\def\PYZlt{\char`\<}
\def\PYZgt{\char`\>}
\def\PYZsh{\char`\#}
\def\PYZpc{\char`\%}
\def\PYZdl{\char`\$}
\def\PYZhy{\char`\-}
\def\PYZsq{\char`\'}
\def\PYZdq{\char`\"}
\def\PYZti{\char`\~}
% for compatibility with earlier versions
\def\PYZat{@}
\def\PYZlb{[}
\def\PYZrb{]}
\makeatother


    % For linebreaks inside Verbatim environment from package fancyvrb.
    \makeatletter
        \newbox\Wrappedcontinuationbox
        \newbox\Wrappedvisiblespacebox
        \newcommand*\Wrappedvisiblespace {\textcolor{red}{\textvisiblespace}}
        \newcommand*\Wrappedcontinuationsymbol {\textcolor{red}{\llap{\tiny$\m@th\hookrightarrow$}}}
        \newcommand*\Wrappedcontinuationindent {3ex }
        \newcommand*\Wrappedafterbreak {\kern\Wrappedcontinuationindent\copy\Wrappedcontinuationbox}
        % Take advantage of the already applied Pygments mark-up to insert
        % potential linebreaks for TeX processing.
        %        {, <, #, %, $, ' and ": go to next line.
        %        _, }, ^, &, >, - and ~: stay at end of broken line.
        % Use of \textquotesingle for straight quote.
        \newcommand*\Wrappedbreaksatspecials {%
            \def\PYGZus{\discretionary{\char`\_}{\Wrappedafterbreak}{\char`\_}}%
            \def\PYGZob{\discretionary{}{\Wrappedafterbreak\char`\{}{\char`\{}}%
            \def\PYGZcb{\discretionary{\char`\}}{\Wrappedafterbreak}{\char`\}}}%
            \def\PYGZca{\discretionary{\char`\^}{\Wrappedafterbreak}{\char`\^}}%
            \def\PYGZam{\discretionary{\char`\&}{\Wrappedafterbreak}{\char`\&}}%
            \def\PYGZlt{\discretionary{}{\Wrappedafterbreak\char`\<}{\char`\<}}%
            \def\PYGZgt{\discretionary{\char`\>}{\Wrappedafterbreak}{\char`\>}}%
            \def\PYGZsh{\discretionary{}{\Wrappedafterbreak\char`\#}{\char`\#}}%
            \def\PYGZpc{\discretionary{}{\Wrappedafterbreak\char`\%}{\char`\%}}%
            \def\PYGZdl{\discretionary{}{\Wrappedafterbreak\char`\$}{\char`\$}}%
            \def\PYGZhy{\discretionary{\char`\-}{\Wrappedafterbreak}{\char`\-}}%
            \def\PYGZsq{\discretionary{}{\Wrappedafterbreak\textquotesingle}{\textquotesingle}}%
            \def\PYGZdq{\discretionary{}{\Wrappedafterbreak\char`\"}{\char`\"}}%
            \def\PYGZti{\discretionary{\char`\~}{\Wrappedafterbreak}{\char`\~}}%
        }
        % Some characters . , ; ? ! / are not pygmentized.
        % This macro makes them "active" and they will insert potential linebreaks
        \newcommand*\Wrappedbreaksatpunct {%
            \lccode`\~`\.\lowercase{\def~}{\discretionary{\hbox{\char`\.}}{\Wrappedafterbreak}{\hbox{\char`\.}}}%
            \lccode`\~`\,\lowercase{\def~}{\discretionary{\hbox{\char`\,}}{\Wrappedafterbreak}{\hbox{\char`\,}}}%
            \lccode`\~`\;\lowercase{\def~}{\discretionary{\hbox{\char`\;}}{\Wrappedafterbreak}{\hbox{\char`\;}}}%
            \lccode`\~`\:\lowercase{\def~}{\discretionary{\hbox{\char`\:}}{\Wrappedafterbreak}{\hbox{\char`\:}}}%
            \lccode`\~`\?\lowercase{\def~}{\discretionary{\hbox{\char`\?}}{\Wrappedafterbreak}{\hbox{\char`\?}}}%
            \lccode`\~`\!\lowercase{\def~}{\discretionary{\hbox{\char`\!}}{\Wrappedafterbreak}{\hbox{\char`\!}}}%
            \lccode`\~`\/\lowercase{\def~}{\discretionary{\hbox{\char`\/}}{\Wrappedafterbreak}{\hbox{\char`\/}}}%
            \catcode`\.\active
            \catcode`\,\active
            \catcode`\;\active
            \catcode`\:\active
            \catcode`\?\active
            \catcode`\!\active
            \catcode`\/\active
            \lccode`\~`\~
        }
    \makeatother

    \let\OriginalVerbatim=\Verbatim
    \makeatletter
    \renewcommand{\Verbatim}[1][1]{%
        %\parskip\z@skip
        \sbox\Wrappedcontinuationbox {\Wrappedcontinuationsymbol}%
        \sbox\Wrappedvisiblespacebox {\FV@SetupFont\Wrappedvisiblespace}%
        \def\FancyVerbFormatLine ##1{\hsize\linewidth
            \vtop{\raggedright\hyphenpenalty\z@\exhyphenpenalty\z@
                \doublehyphendemerits\z@\finalhyphendemerits\z@
                \strut ##1\strut}%
        }%
        % If the linebreak is at a space, the latter will be displayed as visible
        % space at end of first line, and a continuation symbol starts next line.
        % Stretch/shrink are however usually zero for typewriter font.
        \def\FV@Space {%
            \nobreak\hskip\z@ plus\fontdimen3\font minus\fontdimen4\font
            \discretionary{\copy\Wrappedvisiblespacebox}{\Wrappedafterbreak}
            {\kern\fontdimen2\font}%
        }%

        % Allow breaks at special characters using \PYG... macros.
        \Wrappedbreaksatspecials
        % Breaks at punctuation characters . , ; ? ! and / need catcode=\active
        \OriginalVerbatim[#1,codes*=\Wrappedbreaksatpunct]%
    }
    \makeatother

    % Exact colors from NB
    \definecolor{incolor}{HTML}{303F9F}
    \definecolor{outcolor}{HTML}{D84315}
    \definecolor{cellborder}{HTML}{CFCFCF}
    \definecolor{cellbackground}{HTML}{F7F7F7}

    % prompt
    \makeatletter
    \newcommand{\boxspacing}{\kern\kvtcb@left@rule\kern\kvtcb@boxsep}
    \makeatother
    \newcommand{\prompt}[4]{
        {\ttfamily\llap{{\color{#2}[#3]:\hspace{3pt}#4}}\vspace{-\baselineskip}}
    }
    

    
    % Prevent overflowing lines due to hard-to-break entities
    \sloppy
    % Setup hyperref package
    \hypersetup{
      breaklinks=true,  % so long urls are correctly broken across lines
      colorlinks=true,
      urlcolor=urlcolor,
      linkcolor=linkcolor,
      citecolor=citecolor,
      }
    % Slightly bigger margins than the latex defaults
    
    \geometry{verbose,tmargin=1in,bmargin=1in,lmargin=1in,rmargin=1in}
    
    

\begin{document}
    
    \maketitle
    
    

    
    \begin{tcolorbox}[breakable, size=fbox, boxrule=1pt, pad at break*=1mm,colback=cellbackground, colframe=cellborder]
\prompt{In}{incolor}{1}{\boxspacing}
\begin{Verbatim}[commandchars=\\\{\}]
\PY{k+kn}{import}\PY{+w}{ }\PY{n+nn}{polars}\PY{+w}{ }\PY{k}{as}\PY{+w}{ }\PY{n+nn}{pl}
\PY{k+kn}{import}\PY{+w}{ }\PY{n+nn}{pandas}\PY{+w}{ }\PY{k}{as}\PY{+w}{ }\PY{n+nn}{pd}
\PY{k+kn}{import}\PY{+w}{ }\PY{n+nn}{numpy}\PY{+w}{ }\PY{k}{as}\PY{+w}{ }\PY{n+nn}{np}
\PY{k+kn}{import}\PY{+w}{ }\PY{n+nn}{missingno}\PY{+w}{ }\PY{k}{as}\PY{+w}{ }\PY{n+nn}{msno}
\PY{k+kn}{import}\PY{+w}{ }\PY{n+nn}{matplotlib}\PY{n+nn}{.}\PY{n+nn}{pyplot}\PY{+w}{ }\PY{k}{as}\PY{+w}{ }\PY{n+nn}{plt}
\PY{k+kn}{import}\PY{+w}{ }\PY{n+nn}{statsmodels}\PY{n+nn}{.}\PY{n+nn}{api}\PY{+w}{ }\PY{k}{as}\PY{+w}{ }\PY{n+nn}{sm}
\end{Verbatim}
\end{tcolorbox}

    \begin{center}\rule{0.5\linewidth}{0.5pt}\end{center}

\subsubsection{1. Carga y Exploración Inicial de los
Datos}\label{carga-y-exploraciuxf3n-inicial-de-los-datos}

Para comenzar, cargamos la base de datos
\texttt{student\_productivity.csv} al ambiente de trabajo. Utilizamos
los métodos \texttt{.info()} y \texttt{.describe()} para realizar un
análisis exploratorio inicial (EDA). Esto nos permite identificar los
\textbf{tipos de datos} de cada variable (numéricas vs.~categóricas),
observar las estadísticas descriptivas básicas, y detectar a simple
vista irregularidades como \textbf{valores nulos} (missing values) y
posibles \textbf{valores atípicos} (outliers) en las distribuciones.

    \begin{tcolorbox}[breakable, size=fbox, boxrule=1pt, pad at break*=1mm,colback=cellbackground, colframe=cellborder]
\prompt{In}{incolor}{2}{\boxspacing}
\begin{Verbatim}[commandchars=\\\{\}]
\PY{n}{df} \PY{o}{=} \PY{n}{pd}\PY{o}{.}\PY{n}{read\PYZus{}csv}\PY{p}{(}
    \PY{l+s+s2}{\PYZdq{}}\PY{l+s+s2}{data}\PY{l+s+s2}{\PYZbs{}}\PY{l+s+s2}{student\PYZus{}productivity.csv}\PY{l+s+s2}{\PYZdq{}}
\PY{p}{)}
\end{Verbatim}
\end{tcolorbox}

    \begin{tcolorbox}[breakable, size=fbox, boxrule=1pt, pad at break*=1mm,colback=cellbackground, colframe=cellborder]
\prompt{In}{incolor}{3}{\boxspacing}
\begin{Verbatim}[commandchars=\\\{\}]
\PY{n}{df}\PY{o}{.}\PY{n}{info}\PY{p}{(}\PY{p}{)}
\end{Verbatim}
\end{tcolorbox}

    \begin{Verbatim}[commandchars=\\\{\}]
<class 'pandas.DataFrame'>
RangeIndex: 5621 entries, 0 to 5620
Data columns (total 22 columns):
 \#   Column                Non-Null Count  Dtype
---  ------                --------------  -----
 0   student\_id            5621 non-null   int64
 1   age                   4973 non-null   float64
 2   gender                5063 non-null   str
 3   academic\_level        4991 non-null   str
 4   study\_hours           5047 non-null   float64
 5   self\_study\_hours      4954 non-null   float64
 6   online\_classes\_hours  4939 non-null   float64
 7   social\_media\_hours    4965 non-null   float64
 8   gaming\_hours          4965 non-null   float64
 9   sleep\_hours           5068 non-null   float64
 10  screen\_time\_hours     5035 non-null   float64
 11  exercise\_minutes      5022 non-null   str
 12  caffeine\_intake\_mg    5021 non-null   float64
 13  part\_time\_job         4918 non-null   str
 14  upcoming\_deadline     4904 non-null   float64
 15  internet\_quality      5098 non-null   str
 16  mental\_health\_score   4946 non-null   float64
 17  drug\_use              1518 non-null   float64
 18  focus\_index           4988 non-null   float64
 19  burnout\_level         4905 non-null   float64
 20  productivity\_score    4983 non-null   float64
 21  exam\_score            4912 non-null   float64
dtypes: float64(16), int64(1), str(5)
memory usage: 1.1 MB
    \end{Verbatim}

    \begin{tcolorbox}[breakable, size=fbox, boxrule=1pt, pad at break*=1mm,colback=cellbackground, colframe=cellborder]
\prompt{In}{incolor}{4}{\boxspacing}
\begin{Verbatim}[commandchars=\\\{\}]
\PY{n}{df}\PY{o}{.}\PY{n}{describe}\PY{p}{(}\PY{p}{)}
\end{Verbatim}
\end{tcolorbox}

            \begin{tcolorbox}[breakable, size=fbox, boxrule=.5pt, pad at break*=1mm, opacityfill=0]
\prompt{Out}{outcolor}{4}{\boxspacing}
\begin{Verbatim}[commandchars=\\\{\}]
        student\_id          age  study\_hours  self\_study\_hours  \textbackslash{}
count  5621.000000  4973.000000  5047.000000       4954.000000
mean   2811.000000    20.510557     4.538708          2.480454
std    1622.787263     2.876399     1.819412          1.178094
min       1.000000    16.000000     0.000000          0.000000
25\%    1406.000000    18.000000     3.260000          1.660000
50\%    2811.000000    20.000000     4.530000          2.480000
75\%    4216.000000    23.000000     5.760000          3.280000
max    5621.000000    25.000000    11.840000          7.410000

       online\_classes\_hours  social\_media\_hours  gaming\_hours  sleep\_hours  \textbackslash{}
count           4939.000000         4965.000000   4965.000000  5068.000000
mean               2.012296            3.002872      1.571458     7.023301
std                0.983104            1.472740      1.112557     1.159948
min                0.000000            0.000000      0.000000     4.000000
25\%                1.320000            2.000000      0.680000     6.250000
50\%                2.010000            2.990000      1.500000     7.015000
75\%                2.690000            4.030000      2.350000     7.820000
max                6.000000            8.280000      5.640000    10.000000

       screen\_time\_hours  caffeine\_intake\_mg  upcoming\_deadline  \textbackslash{}
count        5035.000000         5021.000000        4904.000000
mean            6.979045          250.804820           0.501223
std             2.482912          143.685652           0.500049
min             1.000000            0.000000           0.000000
25\%             5.280000          129.000000           0.000000
50\%             6.950000          251.000000           1.000000
75\%             8.705000          375.000000           1.000000
max            15.300000          499.000000           1.000000

       mental\_health\_score     drug\_use  focus\_index  burnout\_level  \textbackslash{}
count          4946.000000  1518.000000  4988.000000    4905.000000
mean              5.516983     0.661397    29.388119      45.633156
std               2.873859     0.473391     9.996024      14.257283
min               1.000000     0.000000     1.000000       1.000000
25\%               3.000000     0.000000    22.460000      35.770000
50\%               5.000000     1.000000    29.380000      45.650000
75\%               8.000000     1.000000    36.232500      55.410000
max              10.000000     1.000000    63.480000      97.580000

       productivity\_score   exam\_score
count         4983.000000  4912.000000
mean            37.313526    18.840542
std             16.841997    12.119456
min              1.000000     1.000000
25\%             25.330000     9.410000
50\%             36.920000    17.990000
75\%             49.200000    27.410000
max             98.020000    64.090000
\end{Verbatim}
\end{tcolorbox}
        
    \begin{center}\rule{0.5\linewidth}{0.5pt}\end{center}

\subsubsection{2. Limpieza de Datos: Eliminación de variables no
informativas}\label{limpieza-de-datos-eliminaciuxf3n-de-variables-no-informativas}

Al revisar la información del dataset, se hace evidente que la variable
\texttt{drug\_use} contiene una inmensa cantidad de valores nulos (la
gran mayoría de la muestra). Dado que carece de suficientes datos
reportados para ser estadísticamente significativa en un modelo, se toma
la decisión de \textbf{eliminar esta columna} completamente
(\texttt{drop}) para evitar ruido o distorsiones en los análisis
posteriores.

    \begin{tcolorbox}[breakable, size=fbox, boxrule=1pt, pad at break*=1mm,colback=cellbackground, colframe=cellborder]
\prompt{In}{incolor}{5}{\boxspacing}
\begin{Verbatim}[commandchars=\\\{\}]
\PY{n}{df} \PY{o}{=} \PY{n}{df}\PY{o}{.}\PY{n}{drop}\PY{p}{(}\PY{n}{columns}\PY{o}{=}\PY{p}{[}\PY{l+s+s1}{\PYZsq{}}\PY{l+s+s1}{drug\PYZus{}use}\PY{l+s+s1}{\PYZsq{}}\PY{p}{]}\PY{p}{)}
\end{Verbatim}
\end{tcolorbox}

    \begin{center}\rule{0.5\linewidth}{0.5pt}\end{center}

\subsubsection{\texorpdfstring{3. Limpieza y Procesamiento de la
variable
\texttt{academic\_level}}{3. Limpieza y Procesamiento de la variable academic\_level}}\label{limpieza-y-procesamiento-de-la-variable-academic_level}

La variable categórica \texttt{academic\_level} presentaba dos
problemas: espacios en blanco residuales y valores nulos. 1. Se
limpiaron los \emph{strings} (usando \texttt{.str.strip()}) para
unificar las categorías correctamente. 2. Los valores faltantes se
imputaron utilizando la \textbf{moda} (el valor más frecuente), lo cual
es la técnica estándar y más robusta para mantener la distribución en
variables categóricas. 3. Finalmente, se mapearon las categorías a
valores numéricos enteros para que puedan ser procesados matemáticamente
por los modelos de regresión.

    \begin{tcolorbox}[breakable, size=fbox, boxrule=1pt, pad at break*=1mm,colback=cellbackground, colframe=cellborder]
\prompt{In}{incolor}{6}{\boxspacing}
\begin{Verbatim}[commandchars=\\\{\}]
\PY{c+c1}{\PYZsh{} 1. Verificamos cómo están los datos originalmente}
\PY{n+nb}{print}\PY{p}{(}\PY{l+s+s2}{\PYZdq{}}\PY{l+s+s2}{Categorías originales:}\PY{l+s+s2}{\PYZdq{}}\PY{p}{,} \PY{n}{df}\PY{p}{[}\PY{l+s+s1}{\PYZsq{}}\PY{l+s+s1}{academic\PYZus{}level}\PY{l+s+s1}{\PYZsq{}}\PY{p}{]}\PY{o}{.}\PY{n}{unique}\PY{p}{(}\PY{p}{)}\PY{p}{)}

\PY{c+c1}{\PYZsh{} 2. Eliminamos los espacios en blanco al inicio y al final del texto}
\PY{n}{df}\PY{p}{[}\PY{l+s+s1}{\PYZsq{}}\PY{l+s+s1}{academic\PYZus{}level}\PY{l+s+s1}{\PYZsq{}}\PY{p}{]} \PY{o}{=} \PY{n}{df}\PY{p}{[}\PY{l+s+s1}{\PYZsq{}}\PY{l+s+s1}{academic\PYZus{}level}\PY{l+s+s1}{\PYZsq{}}\PY{p}{]}\PY{o}{.}\PY{n}{str}\PY{o}{.}\PY{n}{strip}\PY{p}{(}\PY{p}{)}
\PY{n+nb}{print}\PY{p}{(}\PY{l+s+s2}{\PYZdq{}}\PY{l+s+s2}{Categorías sin espacios:}\PY{l+s+s2}{\PYZdq{}}\PY{p}{,} \PY{n}{df}\PY{p}{[}\PY{l+s+s1}{\PYZsq{}}\PY{l+s+s1}{academic\PYZus{}level}\PY{l+s+s1}{\PYZsq{}}\PY{p}{]}\PY{o}{.}\PY{n}{unique}\PY{p}{(}\PY{p}{)}\PY{p}{)}

\PY{c+c1}{\PYZsh{} 3. Rellenamos los valores nulos (nan) con la moda (el nivel más común)}

\PY{n}{df}\PY{p}{[}\PY{l+s+s1}{\PYZsq{}}\PY{l+s+s1}{academic\PYZus{}level}\PY{l+s+s1}{\PYZsq{}}\PY{p}{]} \PY{o}{=} \PY{n}{df}\PY{p}{[}\PY{l+s+s1}{\PYZsq{}}\PY{l+s+s1}{academic\PYZus{}level}\PY{l+s+s1}{\PYZsq{}}\PY{p}{]}\PY{o}{.}\PY{n}{fillna}\PY{p}{(}\PY{n}{df}\PY{p}{[}\PY{l+s+s1}{\PYZsq{}}\PY{l+s+s1}{academic\PYZus{}level}\PY{l+s+s1}{\PYZsq{}}\PY{p}{]}\PY{o}{.}\PY{n}{mode}\PY{p}{(}\PY{p}{)}\PY{p}{[}\PY{l+m+mi}{0}\PY{p}{]}\PY{p}{)}
\PY{n}{df}\PY{p}{[}\PY{l+s+s1}{\PYZsq{}}\PY{l+s+s1}{academic\PYZus{}level}\PY{l+s+s1}{\PYZsq{}}\PY{p}{]} \PY{o}{=} \PY{n}{df}\PY{p}{[}\PY{l+s+s1}{\PYZsq{}}\PY{l+s+s1}{academic\PYZus{}level}\PY{l+s+s1}{\PYZsq{}}\PY{p}{]}\PY{o}{.}\PY{n}{map}\PY{p}{(}\PY{p}{\PYZob{}}\PY{l+s+s1}{\PYZsq{}}\PY{l+s+s1}{High School}\PY{l+s+s1}{\PYZsq{}}\PY{p}{:} \PY{l+m+mi}{1}\PY{p}{,} \PY{l+s+s1}{\PYZsq{}}\PY{l+s+s1}{Undergraduate}\PY{l+s+s1}{\PYZsq{}}\PY{p}{:} \PY{l+m+mi}{2}\PY{p}{,} \PY{l+s+s1}{\PYZsq{}}\PY{l+s+s1}{Postgraduate}\PY{l+s+s1}{\PYZsq{}}\PY{p}{:} \PY{l+m+mi}{3}\PY{p}{\PYZcb{}}\PY{p}{)}

\PY{c+c1}{\PYZsh{} Comprobamos el resultado final}
\PY{n+nb}{print}\PY{p}{(}\PY{l+s+s2}{\PYZdq{}}\PY{l+s+s2}{Niveles académicos convertidos a números:}\PY{l+s+s2}{\PYZdq{}}\PY{p}{,} \PY{n}{df}\PY{p}{[}\PY{l+s+s1}{\PYZsq{}}\PY{l+s+s1}{academic\PYZus{}level}\PY{l+s+s1}{\PYZsq{}}\PY{p}{]}\PY{o}{.}\PY{n}{unique}\PY{p}{(}\PY{p}{)}\PY{p}{)}
\end{Verbatim}
\end{tcolorbox}

    \begin{Verbatim}[commandchars=\\\{\}]
Categorías originales: <ArrowStringArray>
[ 'Undergraduate',    'High School', 'Undergraduate ',   'Postgraduate',
              nan,   'High School ',  'Postgraduate ']
Length: 7, dtype: str
Categorías sin espacios: <ArrowStringArray>
['Undergraduate', 'High School', 'Postgraduate', nan]
Length: 4, dtype: str
    \end{Verbatim}

    \begin{Verbatim}[commandchars=\\\{\}]
Niveles académicos convertidos a números: [2 1 3]
    \end{Verbatim}

    \begin{center}\rule{0.5\linewidth}{0.5pt}\end{center}

\subsubsection{\texorpdfstring{4. Manejo de Valores Atípicos (Outliers)
en
\texttt{study\_hours}}{4. Manejo de Valores Atípicos (Outliers) en study\_hours}}\label{manejo-de-valores-atuxedpicos-outliers-en-study_hours}

Un análisis visual mediante un \emph{boxplot} nos permitió confirmar la
presencia de valores atípicos (outliers) en la variable
\texttt{study\_hours}.

Para corregir esto sin perder observaciones valiosas (lo cual ocurriría
si elimináramos las filas), se aplicó el método estadístico de
\textbf{Rango Intercuartílico (IQR)} para identificar los bigotes de la
distribución. Posteriormente, se utilizó la técnica de \textbf{Capping
(acotamiento)} usando \texttt{.clip()}, imponiendo un límite inferior de
cero (para evitar horas negativas) y un límite superior basado en el
cálculo de \(Q3 + 1.5 * IQR\). Esto asegura que los valores se mantengan
dentro de un rango realista sin distorsionar la media.

    \begin{tcolorbox}[breakable, size=fbox, boxrule=1pt, pad at break*=1mm,colback=cellbackground, colframe=cellborder]
\prompt{In}{incolor}{7}{\boxspacing}
\begin{Verbatim}[commandchars=\\\{\}]
\PY{k+kn}{import}\PY{+w}{ }\PY{n+nn}{matplotlib}\PY{n+nn}{.}\PY{n+nn}{pyplot}\PY{+w}{ }\PY{k}{as}\PY{+w}{ }\PY{n+nn}{plt}
\PY{k+kn}{import}\PY{+w}{ }\PY{n+nn}{seaborn}\PY{+w}{ }\PY{k}{as}\PY{+w}{ }\PY{n+nn}{sns}

\PY{n}{sns}\PY{o}{.}\PY{n}{set\PYZus{}theme}\PY{p}{(}\PY{n}{style}\PY{o}{=}\PY{l+s+s2}{\PYZdq{}}\PY{l+s+s2}{whitegrid}\PY{l+s+s2}{\PYZdq{}}\PY{p}{)}

\PY{c+c1}{\PYZsh{} Crear una figura con dos subgráficos}
\PY{n}{plt}\PY{o}{.}\PY{n}{figure}\PY{p}{(}\PY{n}{figsize}\PY{o}{=}\PY{p}{(}\PY{l+m+mi}{14}\PY{p}{,} \PY{l+m+mi}{5}\PY{p}{)}\PY{p}{)}

\PY{c+c1}{\PYZsh{} 2. Boxplot (Diagrama de caja)}
\PY{n}{plt}\PY{o}{.}\PY{n}{subplot}\PY{p}{(}\PY{l+m+mi}{1}\PY{p}{,} \PY{l+m+mi}{2}\PY{p}{,} \PY{l+m+mi}{2}\PY{p}{)}
\PY{n}{sns}\PY{o}{.}\PY{n}{boxplot}\PY{p}{(}\PY{n}{x}\PY{o}{=}\PY{n}{df}\PY{p}{[}\PY{l+s+s1}{\PYZsq{}}\PY{l+s+s1}{study\PYZus{}hours}\PY{l+s+s1}{\PYZsq{}}\PY{p}{]}\PY{o}{.}\PY{n}{dropna}\PY{p}{(}\PY{p}{)}\PY{p}{,} \PY{n}{color}\PY{o}{=}\PY{l+s+s1}{\PYZsq{}}\PY{l+s+s1}{\PYZsh{}FFB347}\PY{l+s+s1}{\PYZsq{}}\PY{p}{)}
\PY{n}{plt}\PY{o}{.}\PY{n}{title}\PY{p}{(}\PY{l+s+s1}{\PYZsq{}}\PY{l+s+s1}{Detección de Outliers en Horas de Estudio}\PY{l+s+s1}{\PYZsq{}}\PY{p}{,} \PY{n}{fontsize}\PY{o}{=}\PY{l+m+mi}{14}\PY{p}{)}
\PY{n}{plt}\PY{o}{.}\PY{n}{xlabel}\PY{p}{(}\PY{l+s+s1}{\PYZsq{}}\PY{l+s+s1}{Horas de Estudio Diarias}\PY{l+s+s1}{\PYZsq{}}\PY{p}{,} \PY{n}{fontsize}\PY{o}{=}\PY{l+m+mi}{12}\PY{p}{)}

\PY{n}{plt}\PY{o}{.}\PY{n}{tight\PYZus{}layout}\PY{p}{(}\PY{p}{)}
\PY{n}{plt}\PY{o}{.}\PY{n}{show}\PY{p}{(}\PY{p}{)}
\end{Verbatim}
\end{tcolorbox}

    \begin{center}
    \adjustimage{max size={0.9\linewidth}{0.9\paperheight}}{Tarea_1_Alvear_Benjamin_files/Tarea_1_Alvear_Benjamin_10_0.png}
    \end{center}
    { \hspace*{\fill} \\}
    
    \begin{tcolorbox}[breakable, size=fbox, boxrule=1pt, pad at break*=1mm,colback=cellbackground, colframe=cellborder]
\prompt{In}{incolor}{8}{\boxspacing}
\begin{Verbatim}[commandchars=\\\{\}]
\PY{c+c1}{\PYZsh{} 1. Calcular el Q1 y Q3}
\PY{n}{Q1} \PY{o}{=} \PY{n}{df}\PY{p}{[}\PY{l+s+s1}{\PYZsq{}}\PY{l+s+s1}{study\PYZus{}hours}\PY{l+s+s1}{\PYZsq{}}\PY{p}{]}\PY{o}{.}\PY{n}{quantile}\PY{p}{(}\PY{l+m+mf}{0.25}\PY{p}{)}
\PY{n}{Q3} \PY{o}{=} \PY{n}{df}\PY{p}{[}\PY{l+s+s1}{\PYZsq{}}\PY{l+s+s1}{study\PYZus{}hours}\PY{l+s+s1}{\PYZsq{}}\PY{p}{]}\PY{o}{.}\PY{n}{quantile}\PY{p}{(}\PY{l+m+mf}{0.75}\PY{p}{)}

\PY{c+c1}{\PYZsh{} 2. Calcular el Rango Intercuartílico (IQR)}
\PY{n}{IQR} \PY{o}{=} \PY{n}{Q3} \PY{o}{\PYZhy{}} \PY{n}{Q1}

\PY{c+c1}{\PYZsh{} 3. Definir el Límite Superior (y el inferior, asegurando que no baje de 0)}
\PY{n}{limite\PYZus{}superior} \PY{o}{=} \PY{n}{Q3} \PY{o}{+} \PY{l+m+mf}{1.5} \PY{o}{*} \PY{n}{IQR}
\PY{n}{limite\PYZus{}inferior} \PY{o}{=} \PY{n+nb}{max}\PY{p}{(}\PY{l+m+mi}{0}\PY{p}{,} \PY{n}{Q1} \PY{o}{\PYZhy{}} \PY{l+m+mf}{1.5} \PY{o}{*} \PY{n}{IQR}\PY{p}{)} 

\PY{n+nb}{print}\PY{p}{(}\PY{l+s+sa}{f}\PY{l+s+s2}{\PYZdq{}}\PY{l+s+s2}{El límite superior calculado es: }\PY{l+s+si}{\PYZob{}}\PY{n}{limite\PYZus{}superior}\PY{l+s+si}{:}\PY{l+s+s2}{.2f}\PY{l+s+si}{\PYZcb{}}\PY{l+s+s2}{ horas.}\PY{l+s+s2}{\PYZdq{}}\PY{p}{)}

\PY{c+c1}{\PYZsh{} 4. Aplicar el Capping usando la función .clip()}
\PY{n}{df}\PY{p}{[}\PY{l+s+s1}{\PYZsq{}}\PY{l+s+s1}{study\PYZus{}hours}\PY{l+s+s1}{\PYZsq{}}\PY{p}{]} \PY{o}{=} \PY{n}{df}\PY{p}{[}\PY{l+s+s1}{\PYZsq{}}\PY{l+s+s1}{study\PYZus{}hours}\PY{l+s+s1}{\PYZsq{}}\PY{p}{]}\PY{o}{.}\PY{n}{clip}\PY{p}{(}\PY{n}{lower}\PY{o}{=}\PY{n}{limite\PYZus{}inferior}\PY{p}{,} \PY{n}{upper}\PY{o}{=}\PY{n}{limite\PYZus{}superior}\PY{p}{)}

\PY{n+nb}{print}\PY{p}{(}\PY{l+s+s2}{\PYZdq{}}\PY{l+s+s2}{¡Outliers topados exitosamente!}\PY{l+s+s2}{\PYZdq{}}\PY{p}{)}
\end{Verbatim}
\end{tcolorbox}

    \begin{Verbatim}[commandchars=\\\{\}]
El límite superior calculado es: 9.51 horas.
¡Outliers topados exitosamente!
    \end{Verbatim}

    \begin{center}\rule{0.5\linewidth}{0.5pt}\end{center}

\subsubsection{\texorpdfstring{5. Generación de Variable Binaria Target
(\texttt{rindio\_examen})}{5. Generación de Variable Binaria Target (rindio\_examen)}}\label{generaciuxf3n-de-variable-binaria-target-rindio_examen}

Las instrucciones indican que un puntaje de \texttt{1.0} en
\texttt{exam\_score} señala a un estudiante que no rindió el examen.
Para preparar nuestra variable dependiente (target) para los modelos de
probabilidad lineal, Logit y Probit, creamos la variable binaria
\texttt{rindio\_examen}. Usando lógica vectorizada (\texttt{np.where}),
asignamos un \texttt{1} a todo estudiante con puntaje estrictamente
mayor a 1.0 (quienes sí rindieron) y un \texttt{0} al resto.

    \begin{tcolorbox}[breakable, size=fbox, boxrule=1pt, pad at break*=1mm,colback=cellbackground, colframe=cellborder]
\prompt{In}{incolor}{9}{\boxspacing}
\begin{Verbatim}[commandchars=\\\{\}]
\PY{n}{df}\PY{p}{[}\PY{l+s+s1}{\PYZsq{}}\PY{l+s+s1}{rindio\PYZus{}examen}\PY{l+s+s1}{\PYZsq{}}\PY{p}{]} \PY{o}{=} \PY{n}{np}\PY{o}{.}\PY{n}{where}\PY{p}{(}\PY{n}{df}\PY{p}{[}\PY{l+s+s1}{\PYZsq{}}\PY{l+s+s1}{exam\PYZus{}score}\PY{l+s+s1}{\PYZsq{}}\PY{p}{]} \PY{o}{\PYZgt{}} \PY{l+m+mf}{1.0}\PY{p}{,} \PY{l+m+mi}{1}\PY{p}{,} \PY{l+m+mi}{0}\PY{p}{)}
\end{Verbatim}
\end{tcolorbox}

    \begin{tcolorbox}[breakable, size=fbox, boxrule=1pt, pad at break*=1mm,colback=cellbackground, colframe=cellborder]
\prompt{In}{incolor}{10}{\boxspacing}
\begin{Verbatim}[commandchars=\\\{\}]
\PY{n}{df}\PY{o}{.}\PY{n}{reset\PYZus{}index}
\end{Verbatim}
\end{tcolorbox}

            \begin{tcolorbox}[breakable, size=fbox, boxrule=.5pt, pad at break*=1mm, opacityfill=0]
\prompt{Out}{outcolor}{10}{\boxspacing}
\begin{Verbatim}[commandchars=\\\{\}]
<bound method DataFrame.reset\_index of       student\_id   age  gender
academic\_level  study\_hours  self\_study\_hours  \textbackslash{}
0              1  20.0     NaN               2         5.37              2.09
1              2  16.0  Female               1         5.85              5.04
2              3  18.0  Female               2         5.69              2.27
3              4  24.0    Male               2         2.32              1.06
4              5  24.0  Female               3         3.87              2.63
{\ldots}          {\ldots}   {\ldots}     {\ldots}             {\ldots}          {\ldots}               {\ldots}
5616        5617  21.0    Male               3         4.16              0.00
5617        5618  25.0  Female               3          NaN               NaN
5618        5619   NaN     NaN               3          NaN               NaN
5619        5620   NaN     NaN               3          NaN               NaN
5620        5621  18.0    Male               3         3.96              0.98

      online\_classes\_hours  social\_media\_hours  gaming\_hours  sleep\_hours  \textbackslash{}
0                     1.85                3.66          2.32         7.73
1                     1.87                3.60          2.79         6.11
2                     0.00                2.93          3.94         7.11
3                     2.30                4.34          2.37         8.54
4                     2.54                3.71          1.30         7.69
{\ldots}                    {\ldots}                 {\ldots}           {\ldots}          {\ldots}
5616                  1.27                4.01          3.41         8.07
5617                   NaN                5.18           NaN          NaN
5618                   NaN                 NaN           NaN          NaN
5619                   NaN                3.52           NaN          NaN
5620                  1.77                4.35          2.22         6.57

      {\ldots}  caffeine\_intake\_mg part\_time\_job  upcoming\_deadline  \textbackslash{}
0     {\ldots}               475.0            No                0.0
1     {\ldots}               362.0            no                0.0
2     {\ldots}               200.0            No                0.0
3     {\ldots}               233.0            No                1.0
4     {\ldots}               159.0            No                1.0
{\ldots}   {\ldots}                 {\ldots}           {\ldots}                {\ldots}
5616  {\ldots}               152.0            no                1.0
5617  {\ldots}                 NaN           NaN                NaN
5618  {\ldots}                 NaN           Yes                NaN
5619  {\ldots}                 NaN           NaN                NaN
5620  {\ldots}               154.0            No                0.0

     internet\_quality  mental\_health\_score focus\_index  burnout\_level  \textbackslash{}
0                Good                  3.0       19.01          31.77
1                Good                 10.0       42.10          45.89
2             Average                  5.0       21.93          37.07
3                Poor                  3.0       13.47          43.63
4                Poor                  2.0       19.95          56.62
{\ldots}               {\ldots}                  {\ldots}         {\ldots}            {\ldots}
5616             Poor                  7.0       21.78          48.07
5617              NaN                  8.0         NaN            NaN
5618              NaN                  NaN         NaN            NaN
5619              NaN                  NaN         NaN            NaN
5620             Poor                  3.0       16.28          36.62

      productivity\_score  exam\_score  rindio\_examen
0                  42.59       25.08              1
1                  67.15       37.83              1
2                  37.68       18.66              1
3                  12.83        1.00              0
4                  18.53        7.78              1
{\ldots}                  {\ldots}         {\ldots}            {\ldots}
5616               43.06       17.15              1
5617               42.52       21.06              1
5618                 NaN         NaN              0
5619                 NaN         NaN              0
5620               21.36        2.81              1

[5621 rows x 22 columns]>
\end{Verbatim}
\end{tcolorbox}
        
    \begin{tcolorbox}[breakable, size=fbox, boxrule=1pt, pad at break*=1mm,colback=cellbackground, colframe=cellborder]
\prompt{In}{incolor}{11}{\boxspacing}
\begin{Verbatim}[commandchars=\\\{\}]
\PY{n}{df}\PY{o}{.}\PY{n}{describe}\PY{p}{(}\PY{p}{)}
\end{Verbatim}
\end{tcolorbox}

            \begin{tcolorbox}[breakable, size=fbox, boxrule=.5pt, pad at break*=1mm, opacityfill=0]
\prompt{Out}{outcolor}{11}{\boxspacing}
\begin{Verbatim}[commandchars=\\\{\}]
        student\_id          age  academic\_level  study\_hours  \textbackslash{}
count  5621.000000  4973.000000     5621.000000  5047.000000
mean   2811.000000    20.510557        2.116883     4.536323
std    1622.787263     2.876399        0.834766     1.812142
min       1.000000    16.000000        1.000000     0.000000
25\%    1406.000000    18.000000        1.000000     3.260000
50\%    2811.000000    20.000000        2.000000     4.530000
75\%    4216.000000    23.000000        3.000000     5.760000
max    5621.000000    25.000000        3.000000     9.510000

       self\_study\_hours  online\_classes\_hours  social\_media\_hours  \textbackslash{}
count       4954.000000           4939.000000         4965.000000
mean           2.480454              2.012296            3.002872
std            1.178094              0.983104            1.472740
min            0.000000              0.000000            0.000000
25\%            1.660000              1.320000            2.000000
50\%            2.480000              2.010000            2.990000
75\%            3.280000              2.690000            4.030000
max            7.410000              6.000000            8.280000

       gaming\_hours  sleep\_hours  screen\_time\_hours  caffeine\_intake\_mg  \textbackslash{}
count   4965.000000  5068.000000        5035.000000         5021.000000
mean       1.571458     7.023301           6.979045          250.804820
std        1.112557     1.159948           2.482912          143.685652
min        0.000000     4.000000           1.000000            0.000000
25\%        0.680000     6.250000           5.280000          129.000000
50\%        1.500000     7.015000           6.950000          251.000000
75\%        2.350000     7.820000           8.705000          375.000000
max        5.640000    10.000000          15.300000          499.000000

       upcoming\_deadline  mental\_health\_score  focus\_index  burnout\_level  \textbackslash{}
count        4904.000000          4946.000000  4988.000000    4905.000000
mean            0.501223             5.516983    29.388119      45.633156
std             0.500049             2.873859     9.996024      14.257283
min             0.000000             1.000000     1.000000       1.000000
25\%             0.000000             3.000000    22.460000      35.770000
50\%             1.000000             5.000000    29.380000      45.650000
75\%             1.000000             8.000000    36.232500      55.410000
max             1.000000            10.000000    63.480000      97.580000

       productivity\_score   exam\_score  rindio\_examen
count         4983.000000  4912.000000    5621.000000
mean            37.313526    18.840542       0.795766
std             16.841997    12.119456       0.403177
min              1.000000     1.000000       0.000000
25\%             25.330000     9.410000       1.000000
50\%             36.920000    17.990000       1.000000
75\%             49.200000    27.410000       1.000000
max             98.020000    64.090000       1.000000
\end{Verbatim}
\end{tcolorbox}
        
    \begin{tcolorbox}[breakable, size=fbox, boxrule=1pt, pad at break*=1mm,colback=cellbackground, colframe=cellborder]
\prompt{In}{incolor}{12}{\boxspacing}
\begin{Verbatim}[commandchars=\\\{\}]
\PY{n}{msno}\PY{o}{.}\PY{n}{matrix}\PY{p}{(}\PY{n}{df}\PY{p}{)}
\end{Verbatim}
\end{tcolorbox}

            \begin{tcolorbox}[breakable, size=fbox, boxrule=.5pt, pad at break*=1mm, opacityfill=0]
\prompt{Out}{outcolor}{12}{\boxspacing}
\begin{Verbatim}[commandchars=\\\{\}]
<Axes: >
\end{Verbatim}
\end{tcolorbox}
        
    \begin{center}
    \adjustimage{max size={0.9\linewidth}{0.9\paperheight}}{Tarea_1_Alvear_Benjamin_files/Tarea_1_Alvear_Benjamin_16_1.png}
    \end{center}
    { \hspace*{\fill} \\}
    
    \begin{center}\rule{0.5\linewidth}{0.5pt}\end{center}

\subsubsection{7. Imputación de Valores Faltantes en Variables
Numéricas}\label{imputaciuxf3n-de-valores-faltantes-en-variables-numuxe9ricas}

Para evitar la pérdida masiva de información, se ha decidido imputar los
datos faltantes en las variables numéricas continuas y discretas. Se
utilizó la \textbf{mediana} (median) en lugar de la media (mean) para
variables como la edad, horas de sueño, horas de juego y niveles de
estrés. Esta decisión se justifica debido a que la mediana es robusta y
no se ve fuertemente afectada por la presencia de valores atípicos
extremos (outliers) que suelen existir en métricas de tiempo o
productividad.

    \begin{tcolorbox}[breakable, size=fbox, boxrule=1pt, pad at break*=1mm,colback=cellbackground, colframe=cellborder]
\prompt{In}{incolor}{13}{\boxspacing}
\begin{Verbatim}[commandchars=\\\{\}]
\PY{c+c1}{\PYZsh{}\PYZsh{}\PYZsh{}\PYZsh{}\PYZsh{}\PYZsh{}\PYZsh{}\PYZsh{}\PYZsh{}\PYZsh{}\PYZsh{}\PYZsh{}\PYZsh{} Variables Númericas \PYZsh{}\PYZsh{}\PYZsh{}\PYZsh{}\PYZsh{}\PYZsh{}\PYZsh{}\PYZsh{}\PYZsh{}\PYZsh{}\PYZsh{}\PYZsh{}\PYZsh{}}
\PY{n}{df}\PY{p}{[}\PY{l+s+s1}{\PYZsq{}}\PY{l+s+s1}{study\PYZus{}hours}\PY{l+s+s1}{\PYZsq{}}\PY{p}{]} \PY{o}{=} \PY{n}{df}\PY{p}{[}\PY{l+s+s1}{\PYZsq{}}\PY{l+s+s1}{study\PYZus{}hours}\PY{l+s+s1}{\PYZsq{}}\PY{p}{]}\PY{o}{.}\PY{n}{fillna}\PY{p}{(}\PY{n}{df}\PY{p}{[}\PY{l+s+s1}{\PYZsq{}}\PY{l+s+s1}{study\PYZus{}hours}\PY{l+s+s1}{\PYZsq{}}\PY{p}{]}\PY{o}{.}\PY{n}{median}\PY{p}{(}\PY{p}{)}\PY{p}{)}
\PY{n}{df}\PY{p}{[}\PY{l+s+s1}{\PYZsq{}}\PY{l+s+s1}{self\PYZus{}study\PYZus{}hours}\PY{l+s+s1}{\PYZsq{}}\PY{p}{]} \PY{o}{=} \PY{n}{df}\PY{p}{[}\PY{l+s+s1}{\PYZsq{}}\PY{l+s+s1}{self\PYZus{}study\PYZus{}hours}\PY{l+s+s1}{\PYZsq{}}\PY{p}{]}\PY{o}{.}\PY{n}{fillna}\PY{p}{(}\PY{n}{df}\PY{p}{[}\PY{l+s+s1}{\PYZsq{}}\PY{l+s+s1}{self\PYZus{}study\PYZus{}hours}\PY{l+s+s1}{\PYZsq{}}\PY{p}{]}\PY{o}{.}\PY{n}{median}\PY{p}{(}\PY{p}{)}\PY{p}{)}
\PY{n}{df}\PY{p}{[}\PY{l+s+s1}{\PYZsq{}}\PY{l+s+s1}{social\PYZus{}media\PYZus{}hours}\PY{l+s+s1}{\PYZsq{}}\PY{p}{]} \PY{o}{=} \PY{n}{df}\PY{p}{[}\PY{l+s+s1}{\PYZsq{}}\PY{l+s+s1}{social\PYZus{}media\PYZus{}hours}\PY{l+s+s1}{\PYZsq{}}\PY{p}{]}\PY{o}{.}\PY{n}{fillna}\PY{p}{(}\PY{n}{df}\PY{p}{[}\PY{l+s+s1}{\PYZsq{}}\PY{l+s+s1}{social\PYZus{}media\PYZus{}hours}\PY{l+s+s1}{\PYZsq{}}\PY{p}{]}\PY{o}{.}\PY{n}{median}\PY{p}{(}\PY{p}{)}\PY{p}{)}
\PY{n}{df}\PY{p}{[}\PY{l+s+s1}{\PYZsq{}}\PY{l+s+s1}{gaming\PYZus{}hours}\PY{l+s+s1}{\PYZsq{}}\PY{p}{]} \PY{o}{=} \PY{n}{df}\PY{p}{[}\PY{l+s+s1}{\PYZsq{}}\PY{l+s+s1}{gaming\PYZus{}hours}\PY{l+s+s1}{\PYZsq{}}\PY{p}{]}\PY{o}{.}\PY{n}{fillna}\PY{p}{(}\PY{n}{df}\PY{p}{[}\PY{l+s+s1}{\PYZsq{}}\PY{l+s+s1}{gaming\PYZus{}hours}\PY{l+s+s1}{\PYZsq{}}\PY{p}{]}\PY{o}{.}\PY{n}{median}\PY{p}{(}\PY{p}{)}\PY{p}{)}
\PY{n}{df}\PY{p}{[}\PY{l+s+s1}{\PYZsq{}}\PY{l+s+s1}{sleep\PYZus{}hours}\PY{l+s+s1}{\PYZsq{}}\PY{p}{]} \PY{o}{=} \PY{n}{df}\PY{p}{[}\PY{l+s+s1}{\PYZsq{}}\PY{l+s+s1}{sleep\PYZus{}hours}\PY{l+s+s1}{\PYZsq{}}\PY{p}{]}\PY{o}{.}\PY{n}{fillna}\PY{p}{(}\PY{n}{df}\PY{p}{[}\PY{l+s+s1}{\PYZsq{}}\PY{l+s+s1}{sleep\PYZus{}hours}\PY{l+s+s1}{\PYZsq{}}\PY{p}{]}\PY{o}{.}\PY{n}{median}\PY{p}{(}\PY{p}{)}\PY{p}{)}
\PY{n}{df}\PY{p}{[}\PY{l+s+s1}{\PYZsq{}}\PY{l+s+s1}{screen\PYZus{}time\PYZus{}hours}\PY{l+s+s1}{\PYZsq{}}\PY{p}{]} \PY{o}{=} \PY{n}{df}\PY{p}{[}\PY{l+s+s1}{\PYZsq{}}\PY{l+s+s1}{screen\PYZus{}time\PYZus{}hours}\PY{l+s+s1}{\PYZsq{}}\PY{p}{]}\PY{o}{.}\PY{n}{fillna}\PY{p}{(}\PY{n}{df}\PY{p}{[}\PY{l+s+s1}{\PYZsq{}}\PY{l+s+s1}{screen\PYZus{}time\PYZus{}hours}\PY{l+s+s1}{\PYZsq{}}\PY{p}{]}\PY{o}{.}\PY{n}{median}\PY{p}{(}\PY{p}{)}\PY{p}{)}
\PY{n}{df}\PY{p}{[}\PY{l+s+s1}{\PYZsq{}}\PY{l+s+s1}{online\PYZus{}classes\PYZus{}hours}\PY{l+s+s1}{\PYZsq{}}\PY{p}{]} \PY{o}{=} \PY{n}{df}\PY{p}{[}\PY{l+s+s1}{\PYZsq{}}\PY{l+s+s1}{online\PYZus{}classes\PYZus{}hours}\PY{l+s+s1}{\PYZsq{}}\PY{p}{]}\PY{o}{.}\PY{n}{fillna}\PY{p}{(}\PY{n}{df}\PY{p}{[}\PY{l+s+s1}{\PYZsq{}}\PY{l+s+s1}{online\PYZus{}classes\PYZus{}hours}\PY{l+s+s1}{\PYZsq{}}\PY{p}{]}\PY{o}{.}\PY{n}{median}\PY{p}{(}\PY{p}{)}\PY{p}{)}

\PY{n}{df}\PY{p}{[}\PY{l+s+s1}{\PYZsq{}}\PY{l+s+s1}{age}\PY{l+s+s1}{\PYZsq{}}\PY{p}{]} \PY{o}{=} \PY{n}{df}\PY{p}{[}\PY{l+s+s1}{\PYZsq{}}\PY{l+s+s1}{age}\PY{l+s+s1}{\PYZsq{}}\PY{p}{]}\PY{o}{.}\PY{n}{fillna}\PY{p}{(}\PY{n}{df}\PY{p}{[}\PY{l+s+s1}{\PYZsq{}}\PY{l+s+s1}{age}\PY{l+s+s1}{\PYZsq{}}\PY{p}{]}\PY{o}{.}\PY{n}{median}\PY{p}{(}\PY{p}{)}\PY{p}{)}
\PY{n}{df}\PY{p}{[}\PY{l+s+s1}{\PYZsq{}}\PY{l+s+s1}{caffeine\PYZus{}intake\PYZus{}mg}\PY{l+s+s1}{\PYZsq{}}\PY{p}{]} \PY{o}{=} \PY{n}{df}\PY{p}{[}\PY{l+s+s1}{\PYZsq{}}\PY{l+s+s1}{caffeine\PYZus{}intake\PYZus{}mg}\PY{l+s+s1}{\PYZsq{}}\PY{p}{]}\PY{o}{.}\PY{n}{fillna}\PY{p}{(}\PY{n}{df}\PY{p}{[}\PY{l+s+s1}{\PYZsq{}}\PY{l+s+s1}{caffeine\PYZus{}intake\PYZus{}mg}\PY{l+s+s1}{\PYZsq{}}\PY{p}{]}\PY{o}{.}\PY{n}{median}\PY{p}{(}\PY{p}{)}\PY{p}{)}
\PY{n}{df}\PY{p}{[}\PY{l+s+s1}{\PYZsq{}}\PY{l+s+s1}{mental\PYZus{}health\PYZus{}score}\PY{l+s+s1}{\PYZsq{}}\PY{p}{]} \PY{o}{=} \PY{n}{df}\PY{p}{[}\PY{l+s+s1}{\PYZsq{}}\PY{l+s+s1}{mental\PYZus{}health\PYZus{}score}\PY{l+s+s1}{\PYZsq{}}\PY{p}{]}\PY{o}{.}\PY{n}{fillna}\PY{p}{(}\PY{n}{df}\PY{p}{[}\PY{l+s+s1}{\PYZsq{}}\PY{l+s+s1}{mental\PYZus{}health\PYZus{}score}\PY{l+s+s1}{\PYZsq{}}\PY{p}{]}\PY{o}{.}\PY{n}{median}\PY{p}{(}\PY{p}{)}\PY{p}{)}
\PY{n}{df}\PY{p}{[}\PY{l+s+s1}{\PYZsq{}}\PY{l+s+s1}{focus\PYZus{}index}\PY{l+s+s1}{\PYZsq{}}\PY{p}{]} \PY{o}{=} \PY{n}{df}\PY{p}{[}\PY{l+s+s1}{\PYZsq{}}\PY{l+s+s1}{focus\PYZus{}index}\PY{l+s+s1}{\PYZsq{}}\PY{p}{]}\PY{o}{.}\PY{n}{fillna}\PY{p}{(}\PY{n}{df}\PY{p}{[}\PY{l+s+s1}{\PYZsq{}}\PY{l+s+s1}{focus\PYZus{}index}\PY{l+s+s1}{\PYZsq{}}\PY{p}{]}\PY{o}{.}\PY{n}{median}\PY{p}{(}\PY{p}{)}\PY{p}{)}
\PY{n}{df}\PY{p}{[}\PY{l+s+s1}{\PYZsq{}}\PY{l+s+s1}{burnout\PYZus{}level}\PY{l+s+s1}{\PYZsq{}}\PY{p}{]} \PY{o}{=} \PY{n}{df}\PY{p}{[}\PY{l+s+s1}{\PYZsq{}}\PY{l+s+s1}{burnout\PYZus{}level}\PY{l+s+s1}{\PYZsq{}}\PY{p}{]}\PY{o}{.}\PY{n}{fillna}\PY{p}{(}\PY{n}{df}\PY{p}{[}\PY{l+s+s1}{\PYZsq{}}\PY{l+s+s1}{burnout\PYZus{}level}\PY{l+s+s1}{\PYZsq{}}\PY{p}{]}\PY{o}{.}\PY{n}{median}\PY{p}{(}\PY{p}{)}\PY{p}{)}
\PY{n}{df}\PY{p}{[}\PY{l+s+s1}{\PYZsq{}}\PY{l+s+s1}{productivity\PYZus{}score}\PY{l+s+s1}{\PYZsq{}}\PY{p}{]} \PY{o}{=} \PY{n}{df}\PY{p}{[}\PY{l+s+s1}{\PYZsq{}}\PY{l+s+s1}{productivity\PYZus{}score}\PY{l+s+s1}{\PYZsq{}}\PY{p}{]}\PY{o}{.}\PY{n}{fillna}\PY{p}{(}\PY{n}{df}\PY{p}{[}\PY{l+s+s1}{\PYZsq{}}\PY{l+s+s1}{productivity\PYZus{}score}\PY{l+s+s1}{\PYZsq{}}\PY{p}{]}\PY{o}{.}\PY{n}{median}\PY{p}{(}\PY{p}{)}\PY{p}{)}
\end{Verbatim}
\end{tcolorbox}

    \begin{center}\rule{0.5\linewidth}{0.5pt}\end{center}

\subsubsection{8. Estandarización y Codificación de Variables
Categóricas}\label{estandarizaciuxf3n-y-codificaciuxf3n-de-variables-categuxf3ricas}

Se procesaron tres variables categóricas clave siguiendo distintos
enfoques según su naturaleza: * \textbf{Calidad de Internet
(\texttt{internet\_quality}):} Al ser una variable categórica ordinal
(pobre \textless{} promedio \textless{} buena), estandarizamos el texto
eliminando mayúsculas y espacios, imputamos los nulos con la moda, y
finalmente la mapeamos a una escala numérica ordinal (1, 2, 3) que
preserva su relación de jerarquía. * \textbf{Género (\texttt{gender}):}
Tras limpiar y rellenar con la moda, aplicamos el método de variables
\emph{Dummy} (\texttt{pd.get\_dummies}) para transformarla a valores
numéricos. Se utilizó \texttt{drop\_first=True} para generar solo la
columna \texttt{gender\_male}, evitando redundancia y previniendo el
error estadístico de ``multicolinealidad perfecta'' crucial para modelos
OLS. * \textbf{Examen inminente (\texttt{upcoming\_deadline}):} Al ya
estar expresada de manera binaria, simplemente cubrimos sus valores
nulos con la moda.

    \begin{tcolorbox}[breakable, size=fbox, boxrule=1pt, pad at break*=1mm,colback=cellbackground, colframe=cellborder]
\prompt{In}{incolor}{14}{\boxspacing}
\begin{Verbatim}[commandchars=\\\{\}]
\PY{c+c1}{\PYZsh{}\PYZsh{}\PYZsh{}\PYZsh{}\PYZsh{}\PYZsh{}\PYZsh{}\PYZsh{}\PYZsh{}\PYZsh{}\PYZsh{}\PYZsh{}\PYZsh{} Variables Categóricas \PYZsh{}\PYZsh{}\PYZsh{}\PYZsh{}\PYZsh{}\PYZsh{}\PYZsh{}\PYZsh{}\PYZsh{}\PYZsh{}\PYZsh{}\PYZsh{}\PYZsh{}}

\PY{c+c1}{\PYZsh{}\PYZhy{}\PYZhy{}\PYZhy{}\PYZhy{}\PYZhy{}\PYZhy{}\PYZhy{}\PYZhy{}\PYZhy{}\PYZhy{}\PYZhy{}\PYZhy{} CALIDAD DE INTERNET \PYZhy{}\PYZhy{}\PYZhy{}\PYZhy{}\PYZhy{}\PYZhy{}\PYZhy{}\PYZhy{}\PYZhy{}\PYZhy{}\PYZhy{}\PYZhy{}\PYZsh{}}
\PY{c+c1}{\PYZsh{} 1. Transformamos a minúsculas y limpiamos espacios para evitar errores de tipeo}
\PY{n}{df}\PY{p}{[}\PY{l+s+s1}{\PYZsq{}}\PY{l+s+s1}{internet\PYZus{}quality}\PY{l+s+s1}{\PYZsq{}}\PY{p}{]} \PY{o}{=} \PY{n}{df}\PY{p}{[}\PY{l+s+s1}{\PYZsq{}}\PY{l+s+s1}{internet\PYZus{}quality}\PY{l+s+s1}{\PYZsq{}}\PY{p}{]}\PY{o}{.}\PY{n}{str}\PY{o}{.}\PY{n}{lower}\PY{p}{(}\PY{p}{)}\PY{o}{.}\PY{n}{str}\PY{o}{.}\PY{n}{strip}\PY{p}{(}\PY{p}{)}
\PY{c+c1}{\PYZsh{} 2. Rellenamos nulos con la moda}
\PY{n}{moda\PYZus{}internet} \PY{o}{=} \PY{n}{df}\PY{p}{[}\PY{l+s+s1}{\PYZsq{}}\PY{l+s+s1}{internet\PYZus{}quality}\PY{l+s+s1}{\PYZsq{}}\PY{p}{]}\PY{o}{.}\PY{n}{mode}\PY{p}{(}\PY{p}{)}\PY{p}{[}\PY{l+m+mi}{0}\PY{p}{]}
\PY{n}{df}\PY{p}{[}\PY{l+s+s1}{\PYZsq{}}\PY{l+s+s1}{internet\PYZus{}quality}\PY{l+s+s1}{\PYZsq{}}\PY{p}{]} \PY{o}{=} \PY{n}{df}\PY{p}{[}\PY{l+s+s1}{\PYZsq{}}\PY{l+s+s1}{internet\PYZus{}quality}\PY{l+s+s1}{\PYZsq{}}\PY{p}{]}\PY{o}{.}\PY{n}{fillna}\PY{p}{(}\PY{n}{moda\PYZus{}internet}\PY{p}{)}
\PY{c+c1}{\PYZsh{} 3. Transformación a valores numéricos ORDINALES}
\PY{n}{df}\PY{p}{[}\PY{l+s+s1}{\PYZsq{}}\PY{l+s+s1}{internet\PYZus{}quality}\PY{l+s+s1}{\PYZsq{}}\PY{p}{]} \PY{o}{=} \PY{n}{df}\PY{p}{[}\PY{l+s+s1}{\PYZsq{}}\PY{l+s+s1}{internet\PYZus{}quality}\PY{l+s+s1}{\PYZsq{}}\PY{p}{]}\PY{o}{.}\PY{n}{map}\PY{p}{(}\PY{p}{\PYZob{}}\PY{l+s+s1}{\PYZsq{}}\PY{l+s+s1}{poor}\PY{l+s+s1}{\PYZsq{}}\PY{p}{:} \PY{l+m+mi}{1}\PY{p}{,} \PY{l+s+s1}{\PYZsq{}}\PY{l+s+s1}{average}\PY{l+s+s1}{\PYZsq{}}\PY{p}{:} \PY{l+m+mi}{2}\PY{p}{,} \PY{l+s+s1}{\PYZsq{}}\PY{l+s+s1}{good}\PY{l+s+s1}{\PYZsq{}}\PY{p}{:} \PY{l+m+mi}{3}\PY{p}{\PYZcb{}}\PY{p}{)}
\PY{n+nb}{print}\PY{p}{(}\PY{l+s+s2}{\PYZdq{}}\PY{l+s+s2}{Formato final de internet\PYZus{}quality:}\PY{l+s+s2}{\PYZdq{}}\PY{p}{)}
\PY{n+nb}{print}\PY{p}{(}\PY{n}{df}\PY{p}{[}\PY{l+s+s1}{\PYZsq{}}\PY{l+s+s1}{internet\PYZus{}quality}\PY{l+s+s1}{\PYZsq{}}\PY{p}{]}\PY{o}{.}\PY{n}{value\PYZus{}counts}\PY{p}{(}\PY{p}{)}\PY{p}{)}

\PY{c+c1}{\PYZsh{}\PYZhy{}\PYZhy{}\PYZhy{}\PYZhy{}\PYZhy{}\PYZhy{}\PYZhy{}\PYZhy{}\PYZhy{}\PYZhy{}\PYZhy{}\PYZhy{} GÉNERO \PYZhy{}\PYZhy{}\PYZhy{}\PYZhy{}\PYZhy{}\PYZhy{}\PYZhy{}\PYZhy{}\PYZhy{}\PYZhy{}\PYZhy{}\PYZhy{}\PYZsh{}}
\PY{c+c1}{\PYZsh{} 1. Limpieza de texto}
\PY{n}{df}\PY{p}{[}\PY{l+s+s1}{\PYZsq{}}\PY{l+s+s1}{gender}\PY{l+s+s1}{\PYZsq{}}\PY{p}{]} \PY{o}{=} \PY{n}{df}\PY{p}{[}\PY{l+s+s1}{\PYZsq{}}\PY{l+s+s1}{gender}\PY{l+s+s1}{\PYZsq{}}\PY{p}{]}\PY{o}{.}\PY{n}{str}\PY{o}{.}\PY{n}{lower}\PY{p}{(}\PY{p}{)}\PY{o}{.}\PY{n}{str}\PY{o}{.}\PY{n}{strip}\PY{p}{(}\PY{p}{)}

\PY{c+c1}{\PYZsh{} 2. Rellenar nulos con la moda predominante}
\PY{n}{moda\PYZus{}genero} \PY{o}{=} \PY{n}{df}\PY{p}{[}\PY{l+s+s1}{\PYZsq{}}\PY{l+s+s1}{gender}\PY{l+s+s1}{\PYZsq{}}\PY{p}{]}\PY{o}{.}\PY{n}{mode}\PY{p}{(}\PY{p}{)}\PY{p}{[}\PY{l+m+mi}{0}\PY{p}{]}
\PY{n}{df}\PY{p}{[}\PY{l+s+s1}{\PYZsq{}}\PY{l+s+s1}{gender}\PY{l+s+s1}{\PYZsq{}}\PY{p}{]} \PY{o}{=} \PY{n}{df}\PY{p}{[}\PY{l+s+s1}{\PYZsq{}}\PY{l+s+s1}{gender}\PY{l+s+s1}{\PYZsq{}}\PY{p}{]}\PY{o}{.}\PY{n}{fillna}\PY{p}{(}\PY{n}{moda\PYZus{}genero}\PY{p}{)}

\PY{c+c1}{\PYZsh{} 3. Transformar a Dummies (Esto creará una nueva columna binaria)}
\PY{c+c1}{\PYZsh{} drop\PYZus{}first=True evita multicolinealidad perfecta (la famosa \PYZdq{}trampa dummy\PYZdq{}) en el OLS}
\PY{n}{df} \PY{o}{=} \PY{n}{pd}\PY{o}{.}\PY{n}{get\PYZus{}dummies}\PY{p}{(}\PY{n}{df}\PY{p}{,} \PY{n}{columns}\PY{o}{=}\PY{p}{[}\PY{l+s+s1}{\PYZsq{}}\PY{l+s+s1}{gender}\PY{l+s+s1}{\PYZsq{}}\PY{p}{]}\PY{p}{,} \PY{n}{drop\PYZus{}first}\PY{o}{=}\PY{k+kc}{True}\PY{p}{,} \PY{n}{dtype}\PY{o}{=}\PY{n+nb}{int}\PY{p}{)}

\PY{c+c1}{\PYZsh{} Ahora existirá una columna llamada probablemente \PYZsq{}gender\PYZus{}male\PYZsq{} con 1 y 0}
\PY{n+nb}{print}\PY{p}{(}\PY{l+s+s2}{\PYZdq{}}\PY{l+s+s2}{¡Conversión dummy exitosa! Generada la columna masculina/femenina.}\PY{l+s+s2}{\PYZdq{}}\PY{p}{)}

\PY{c+c1}{\PYZsh{}\PYZhy{}\PYZhy{}\PYZhy{}\PYZhy{}\PYZhy{}\PYZhy{}\PYZhy{}\PYZhy{}\PYZhy{}\PYZhy{}\PYZhy{}\PYZhy{} EXAMEN PRONTO \PYZhy{}\PYZhy{}\PYZhy{}\PYZhy{}\PYZhy{}\PYZhy{}\PYZhy{}\PYZhy{}\PYZhy{}\PYZhy{}\PYZhy{}\PYZhy{}\PYZsh{}}
\PY{c+c1}{\PYZsh{} 1. Rellenamos los nulos con la moda (ya es 0 y 1, por lo que usaremos la clase mayoritaria)}
\PY{n}{moda\PYZus{}deadline} \PY{o}{=} \PY{n}{df}\PY{p}{[}\PY{l+s+s1}{\PYZsq{}}\PY{l+s+s1}{upcoming\PYZus{}deadline}\PY{l+s+s1}{\PYZsq{}}\PY{p}{]}\PY{o}{.}\PY{n}{mode}\PY{p}{(}\PY{p}{)}\PY{p}{[}\PY{l+m+mi}{0}\PY{p}{]}
\PY{n}{df}\PY{p}{[}\PY{l+s+s1}{\PYZsq{}}\PY{l+s+s1}{upcoming\PYZus{}deadline}\PY{l+s+s1}{\PYZsq{}}\PY{p}{]} \PY{o}{=} \PY{n}{df}\PY{p}{[}\PY{l+s+s1}{\PYZsq{}}\PY{l+s+s1}{upcoming\PYZus{}deadline}\PY{l+s+s1}{\PYZsq{}}\PY{p}{]}\PY{o}{.}\PY{n}{fillna}\PY{p}{(}\PY{n}{moda\PYZus{}deadline}\PY{p}{)}

\PY{n+nb}{print}\PY{p}{(}\PY{l+s+s2}{\PYZdq{}}\PY{l+s+s2}{Nulos eliminados en upcoming\PYZus{}deadline.}\PY{l+s+s2}{\PYZdq{}}\PY{p}{)}
\end{Verbatim}
\end{tcolorbox}

    \begin{Verbatim}[commandchars=\\\{\}]
Formato final de internet\_quality:
internet\_quality
3    2269
2    1680
1    1672
Name: count, dtype: int64
¡Conversión dummy exitosa! Generada la columna masculina/femenina.
Nulos eliminados en upcoming\_deadline.
    \end{Verbatim}

    \begin{center}\rule{0.5\linewidth}{0.5pt}\end{center}

\subsubsection{\texorpdfstring{9. Binarización de la variable
\texttt{part\_time\_job}}{9. Binarización de la variable part\_time\_job}}\label{binarizaciuxf3n-de-la-variable-part_time_job}

Se repitió el proceso de limpieza de texto para unificar las respuestas
a simplemente ``yes'' o ``no''. Tras asegurar que no hubiese espacios ni
mayúsculas que generaran categorías duplicadas, imputamos con la moda y
transformamos directamente usando un diccionario manual
(\texttt{.map()}) a valores binarios matemáticos: 1 para ``sí'' y 0 para
``no'', dejándola lista para la regresión.

    \begin{tcolorbox}[breakable, size=fbox, boxrule=1pt, pad at break*=1mm,colback=cellbackground, colframe=cellborder]
\prompt{In}{incolor}{15}{\boxspacing}
\begin{Verbatim}[commandchars=\\\{\}]
\PY{c+c1}{\PYZsh{} 1. Uniformamos texto (convertimos todo a minúsculas y limpiamos espacios)}
\PY{n}{df}\PY{p}{[}\PY{l+s+s1}{\PYZsq{}}\PY{l+s+s1}{part\PYZus{}time\PYZus{}job}\PY{l+s+s1}{\PYZsq{}}\PY{p}{]} \PY{o}{=} \PY{n}{df}\PY{p}{[}\PY{l+s+s1}{\PYZsq{}}\PY{l+s+s1}{part\PYZus{}time\PYZus{}job}\PY{l+s+s1}{\PYZsq{}}\PY{p}{]}\PY{o}{.}\PY{n}{str}\PY{o}{.}\PY{n}{lower}\PY{p}{(}\PY{p}{)}\PY{o}{.}\PY{n}{str}\PY{o}{.}\PY{n}{strip}\PY{p}{(}\PY{p}{)}

\PY{c+c1}{\PYZsh{} Validamos que ahora sólo existan \PYZsq{}yes\PYZsq{} o \PYZsq{}no\PYZsq{} (o posibles nulos)}
\PY{n+nb}{print}\PY{p}{(}\PY{l+s+s2}{\PYZdq{}}\PY{l+s+s2}{Categorías extrañas corregidas:}\PY{l+s+s2}{\PYZdq{}}\PY{p}{,} \PY{n}{df}\PY{p}{[}\PY{l+s+s1}{\PYZsq{}}\PY{l+s+s1}{part\PYZus{}time\PYZus{}job}\PY{l+s+s1}{\PYZsq{}}\PY{p}{]}\PY{o}{.}\PY{n}{unique}\PY{p}{(}\PY{p}{)}\PY{p}{)}

\PY{c+c1}{\PYZsh{} 2. Rellenamos datos faltantes (nulos) con la moda (que suele ser \PYZsq{}no\PYZsq{})}
\PY{k}{if} \PY{n}{df}\PY{p}{[}\PY{l+s+s1}{\PYZsq{}}\PY{l+s+s1}{part\PYZus{}time\PYZus{}job}\PY{l+s+s1}{\PYZsq{}}\PY{p}{]}\PY{o}{.}\PY{n}{isnull}\PY{p}{(}\PY{p}{)}\PY{o}{.}\PY{n}{sum}\PY{p}{(}\PY{p}{)} \PY{o}{\PYZgt{}} \PY{l+m+mi}{0}\PY{p}{:}
    \PY{n}{moda\PYZus{}trabajo} \PY{o}{=} \PY{n}{df}\PY{p}{[}\PY{l+s+s1}{\PYZsq{}}\PY{l+s+s1}{part\PYZus{}time\PYZus{}job}\PY{l+s+s1}{\PYZsq{}}\PY{p}{]}\PY{o}{.}\PY{n}{mode}\PY{p}{(}\PY{p}{)}\PY{p}{[}\PY{l+m+mi}{0}\PY{p}{]}
    \PY{n}{df}\PY{p}{[}\PY{l+s+s1}{\PYZsq{}}\PY{l+s+s1}{part\PYZus{}time\PYZus{}job}\PY{l+s+s1}{\PYZsq{}}\PY{p}{]} \PY{o}{=} \PY{n}{df}\PY{p}{[}\PY{l+s+s1}{\PYZsq{}}\PY{l+s+s1}{part\PYZus{}time\PYZus{}job}\PY{l+s+s1}{\PYZsq{}}\PY{p}{]}\PY{o}{.}\PY{n}{fillna}\PY{p}{(}\PY{n}{moda\PYZus{}trabajo}\PY{p}{)}

\PY{c+c1}{\PYZsh{} 3. Transformamos a binario (1 y 0) en la misma columna}
\PY{n}{df}\PY{p}{[}\PY{l+s+s1}{\PYZsq{}}\PY{l+s+s1}{part\PYZus{}time\PYZus{}job}\PY{l+s+s1}{\PYZsq{}}\PY{p}{]} \PY{o}{=} \PY{n}{df}\PY{p}{[}\PY{l+s+s1}{\PYZsq{}}\PY{l+s+s1}{part\PYZus{}time\PYZus{}job}\PY{l+s+s1}{\PYZsq{}}\PY{p}{]}\PY{o}{.}\PY{n}{map}\PY{p}{(}\PY{p}{\PYZob{}}\PY{l+s+s1}{\PYZsq{}}\PY{l+s+s1}{yes}\PY{l+s+s1}{\PYZsq{}}\PY{p}{:} \PY{l+m+mi}{1}\PY{p}{,} \PY{l+s+s1}{\PYZsq{}}\PY{l+s+s1}{no}\PY{l+s+s1}{\PYZsq{}}\PY{p}{:} \PY{l+m+mi}{0}\PY{p}{\PYZcb{}}\PY{p}{)}

\PY{n+nb}{print}\PY{p}{(}\PY{l+s+s2}{\PYZdq{}}\PY{l+s+s2}{Formato final de part\PYZus{}time\PYZus{}job:}\PY{l+s+s2}{\PYZdq{}}\PY{p}{)}
\PY{n+nb}{print}\PY{p}{(}\PY{n}{df}\PY{p}{[}\PY{l+s+s1}{\PYZsq{}}\PY{l+s+s1}{part\PYZus{}time\PYZus{}job}\PY{l+s+s1}{\PYZsq{}}\PY{p}{]}\PY{o}{.}\PY{n}{value\PYZus{}counts}\PY{p}{(}\PY{p}{)}\PY{p}{)}
\end{Verbatim}
\end{tcolorbox}

    \begin{Verbatim}[commandchars=\\\{\}]
Categorías extrañas corregidas: <ArrowStringArray>
['no', 'yes', nan]
Length: 3, dtype: str
Formato final de part\_time\_job:
part\_time\_job
0    3172
1    2449
Name: count, dtype: int64
    \end{Verbatim}

    \begin{center}\rule{0.5\linewidth}{0.5pt}\end{center}

\subsubsection{\texorpdfstring{10. Extracción de datos en columnas
sucias
(\texttt{exercise\_minutes})}{10. Extracción de datos en columnas sucias (exercise\_minutes)}}\label{extracciuxf3n-de-datos-en-columnas-sucias-exercise_minutes}

Se utilizó una \textbf{Expresión Regular (Regex)} mediante la función
\texttt{.str.extract(r\textquotesingle{}(\textbackslash{}d+)\textquotesingle{})}
para ``capturar'' únicamente los dígitos numéricos dentro del string,
parseándolos a formato Flotante (numérico) y luego rellenando los vacíos
con la mediana representativa.

    \begin{tcolorbox}[breakable, size=fbox, boxrule=1pt, pad at break*=1mm,colback=cellbackground, colframe=cellborder]
\prompt{In}{incolor}{16}{\boxspacing}
\begin{Verbatim}[commandchars=\\\{\}]
\PY{c+c1}{\PYZsh{} 1. Extraemos SOLO los números de la columna de texto (ignora letras como \PYZdq{}mins\PYZdq{})}
\PY{n}{df}\PY{p}{[}\PY{l+s+s1}{\PYZsq{}}\PY{l+s+s1}{exercise\PYZus{}minutes}\PY{l+s+s1}{\PYZsq{}}\PY{p}{]} \PY{o}{=} \PY{n}{df}\PY{p}{[}\PY{l+s+s1}{\PYZsq{}}\PY{l+s+s1}{exercise\PYZus{}minutes}\PY{l+s+s1}{\PYZsq{}}\PY{p}{]}\PY{o}{.}\PY{n}{astype}\PY{p}{(}\PY{n+nb}{str}\PY{p}{)}\PY{o}{.}\PY{n}{str}\PY{o}{.}\PY{n}{extract}\PY{p}{(}\PY{l+s+sa}{r}\PY{l+s+s1}{\PYZsq{}}\PY{l+s+s1}{(}\PY{l+s+s1}{\PYZbs{}}\PY{l+s+s1}{d+)}\PY{l+s+s1}{\PYZsq{}}\PY{p}{)}

\PY{c+c1}{\PYZsh{} 2. Lo convertimos a número flotante (float) para poder operarlo matemáticamente}
\PY{n}{df}\PY{p}{[}\PY{l+s+s1}{\PYZsq{}}\PY{l+s+s1}{exercise\PYZus{}minutes}\PY{l+s+s1}{\PYZsq{}}\PY{p}{]} \PY{o}{=} \PY{n}{df}\PY{p}{[}\PY{l+s+s1}{\PYZsq{}}\PY{l+s+s1}{exercise\PYZus{}minutes}\PY{l+s+s1}{\PYZsq{}}\PY{p}{]}\PY{o}{.}\PY{n}{astype}\PY{p}{(}\PY{n+nb}{float}\PY{p}{)}

\PY{c+c1}{\PYZsh{} 3. Rellenamos posibles vacíos generados (si alguien escribió \PYZdq{}nada\PYZdq{}) con la mediana}
\PY{n}{mediana\PYZus{}ejercicio} \PY{o}{=} \PY{n}{df}\PY{p}{[}\PY{l+s+s1}{\PYZsq{}}\PY{l+s+s1}{exercise\PYZus{}minutes}\PY{l+s+s1}{\PYZsq{}}\PY{p}{]}\PY{o}{.}\PY{n}{median}\PY{p}{(}\PY{p}{)}
\PY{n}{df}\PY{p}{[}\PY{l+s+s1}{\PYZsq{}}\PY{l+s+s1}{exercise\PYZus{}minutes}\PY{l+s+s1}{\PYZsq{}}\PY{p}{]} \PY{o}{=} \PY{n}{df}\PY{p}{[}\PY{l+s+s1}{\PYZsq{}}\PY{l+s+s1}{exercise\PYZus{}minutes}\PY{l+s+s1}{\PYZsq{}}\PY{p}{]}\PY{o}{.}\PY{n}{fillna}\PY{p}{(}\PY{n}{mediana\PYZus{}ejercicio}\PY{p}{)}

\PY{n+nb}{print}\PY{p}{(}\PY{l+s+s2}{\PYZdq{}}\PY{l+s+s2}{¡Problema resuelto! exercise\PYZus{}minutes ahora es numérico.}\PY{l+s+s2}{\PYZdq{}}\PY{p}{)}
\PY{c+c1}{\PYZsh{} Puedes comprobarlo haciendo: print(df[\PYZsq{}exercise\PYZus{}minutes\PYZsq{}].dtype)}
\end{Verbatim}
\end{tcolorbox}

    \begin{Verbatim}[commandchars=\\\{\}]
¡Problema resuelto! exercise\_minutes ahora es numérico.
    \end{Verbatim}

    \begin{center}\rule{0.5\linewidth}{0.5pt}\end{center}

\subsubsection{11. Control de Calidad y Validación Lógica: Límite de las
24 Horas
Diarias}\label{control-de-calidad-y-validaciuxf3n-luxf3gica-luxedmite-de-las-24-horas-diarias}

Uno de los problemas lógicos más comunes en encuestas auto-reportadas
sobre uso del tiempo, es que las personas pueden exagerar u omitir sumar
correctamente, terminando en totales diarios que superan físicamente las
24 horas del día. Para auditar esto: 1. Calculamos la suma total de
horas activas (restando factores que podrían ocurrir simultáneamente
pero manteniendo el rubro principal: \texttt{study\_hours},
\texttt{sleep\_hours}, \texttt{social\_media}, \texttt{gaming} y
\texttt{exercise}). 2. Identificamos qué alumnos rompían la regla lógica
sumando más de 24 horas. 3. Se aplicó una normalización. Modificamos el
valor de dichas horas multiplicándolo por la proporción correcta
\((24 / \text{Suma de sus horas})\). De esta forma, castigamos el exceso
pero mantenemos de manera realista la ``proporción de tiempo'' que la
persona le dedicó a cada actividad en su día.

    \begin{tcolorbox}[breakable, size=fbox, boxrule=1pt, pad at break*=1mm,colback=cellbackground, colframe=cellborder]
\prompt{In}{incolor}{17}{\boxspacing}
\begin{Verbatim}[commandchars=\\\{\}]
\PY{c+c1}{\PYZsh{} 1. Definimos las columnas base de uso del tiempo (excluyendo las que se solapan)}
\PY{n}{columnas\PYZus{}horas} \PY{o}{=} \PY{p}{[}
    \PY{l+s+s1}{\PYZsq{}}\PY{l+s+s1}{study\PYZus{}hours}\PY{l+s+s1}{\PYZsq{}}\PY{p}{,}
    \PY{l+s+s1}{\PYZsq{}}\PY{l+s+s1}{sleep\PYZus{}hours}\PY{l+s+s1}{\PYZsq{}}\PY{p}{,} 
    \PY{l+s+s1}{\PYZsq{}}\PY{l+s+s1}{social\PYZus{}media\PYZus{}hours}\PY{l+s+s1}{\PYZsq{}}\PY{p}{,} 
    \PY{l+s+s1}{\PYZsq{}}\PY{l+s+s1}{gaming\PYZus{}hours}\PY{l+s+s1}{\PYZsq{}}
\PY{p}{]}

\PY{c+c1}{\PYZsh{} (Opcional) Podemos incluir los minutos de ejercicio convirtiéndolos a horas}
\PY{n}{df}\PY{p}{[}\PY{l+s+s1}{\PYZsq{}}\PY{l+s+s1}{exercise\PYZus{}hours}\PY{l+s+s1}{\PYZsq{}}\PY{p}{]} \PY{o}{=} \PY{n}{df}\PY{p}{[}\PY{l+s+s1}{\PYZsq{}}\PY{l+s+s1}{exercise\PYZus{}minutes}\PY{l+s+s1}{\PYZsq{}}\PY{p}{]} \PY{o}{/} \PY{l+m+mf}{60.0}
\PY{n}{columnas\PYZus{}horas}\PY{o}{.}\PY{n}{append}\PY{p}{(}\PY{l+s+s1}{\PYZsq{}}\PY{l+s+s1}{exercise\PYZus{}hours}\PY{l+s+s1}{\PYZsq{}}\PY{p}{)}

\PY{c+c1}{\PYZsh{} 2. Calculamos las horas totales de cada alumno}
\PY{n}{df}\PY{p}{[}\PY{l+s+s1}{\PYZsq{}}\PY{l+s+s1}{suma\PYZus{}horas}\PY{l+s+s1}{\PYZsq{}}\PY{p}{]} \PY{o}{=} \PY{n}{df}\PY{p}{[}\PY{n}{columnas\PYZus{}horas}\PY{p}{]}\PY{o}{.}\PY{n}{sum}\PY{p}{(}\PY{n}{axis}\PY{o}{=}\PY{l+m+mi}{1}\PY{p}{)}

\PY{c+c1}{\PYZsh{} Filtramos para observar cuántos tienen el error lógico}
\PY{n}{filtro\PYZus{}sobrepasan} \PY{o}{=} \PY{n}{df}\PY{p}{[}\PY{l+s+s1}{\PYZsq{}}\PY{l+s+s1}{suma\PYZus{}horas}\PY{l+s+s1}{\PYZsq{}}\PY{p}{]} \PY{o}{\PYZgt{}} \PY{l+m+mf}{24.0}
\PY{n+nb}{print}\PY{p}{(}\PY{l+s+sa}{f}\PY{l+s+s2}{\PYZdq{}}\PY{l+s+s2}{Alumnos que reportaron imposiblemente más de 24 hrs al día: }\PY{l+s+si}{\PYZob{}}\PY{n}{filtro\PYZus{}sobrepasan}\PY{o}{.}\PY{n}{sum}\PY{p}{(}\PY{p}{)}\PY{l+s+si}{\PYZcb{}}\PY{l+s+s2}{\PYZdq{}}\PY{p}{)}

\PY{c+c1}{\PYZsh{} 3. Solución (Proporción): Escalar las horas para que el máximo sea 24}
\PY{k}{for} \PY{n}{col} \PY{o+ow}{in} \PY{n}{columnas\PYZus{}horas}\PY{p}{:}
    \PY{c+c1}{\PYZsh{} Fórmula: Valor nuevo = Valor antiguo * (24 / Suma de sus horas)}
    \PY{n}{df}\PY{o}{.}\PY{n}{loc}\PY{p}{[}\PY{n}{filtro\PYZus{}sobrepasan}\PY{p}{,} \PY{n}{col}\PY{p}{]} \PY{o}{=} \PY{n}{df}\PY{o}{.}\PY{n}{loc}\PY{p}{[}\PY{n}{filtro\PYZus{}sobrepasan}\PY{p}{,} \PY{n}{col}\PY{p}{]} \PY{o}{*} \PY{p}{(}\PY{l+m+mf}{24.0} \PY{o}{/} \PY{n}{df}\PY{o}{.}\PY{n}{loc}\PY{p}{[}\PY{n}{filtro\PYZus{}sobrepasan}\PY{p}{,} \PY{l+s+s1}{\PYZsq{}}\PY{l+s+s1}{suma\PYZus{}horas}\PY{l+s+s1}{\PYZsq{}}\PY{p}{]}\PY{p}{)}

\PY{c+c1}{\PYZsh{} Comprobamos el resultado (usamos \PYZgt{} 24.01 por los decimales infinitos de Python)}
\PY{n}{df}\PY{p}{[}\PY{l+s+s1}{\PYZsq{}}\PY{l+s+s1}{suma\PYZus{}horas}\PY{l+s+s1}{\PYZsq{}}\PY{p}{]} \PY{o}{=} \PY{n}{df}\PY{p}{[}\PY{n}{columnas\PYZus{}horas}\PY{p}{]}\PY{o}{.}\PY{n}{sum}\PY{p}{(}\PY{n}{axis}\PY{o}{=}\PY{l+m+mi}{1}\PY{p}{)}
\PY{n+nb}{print}\PY{p}{(}\PY{l+s+sa}{f}\PY{l+s+s2}{\PYZdq{}}\PY{l+s+s2}{Alumnos excedidos tras la limpieza: }\PY{l+s+si}{\PYZob{}}\PY{p}{(}\PY{n}{df}\PY{p}{[}\PY{l+s+s1}{\PYZsq{}}\PY{l+s+s1}{suma\PYZus{}horas}\PY{l+s+s1}{\PYZsq{}}\PY{p}{]}\PY{+w}{ }\PY{o}{\PYZgt{}}\PY{+w}{ }\PY{l+m+mf}{24.01}\PY{p}{)}\PY{o}{.}\PY{n}{sum}\PY{p}{(}\PY{p}{)}\PY{l+s+si}{\PYZcb{}}\PY{l+s+s2}{\PYZdq{}}\PY{p}{)}

\PY{c+c1}{\PYZsh{} 4. Limpieza final: eliminamos las columnas temporales (dejamos los minutos originales para no alterar la data madre)}
\PY{n}{df} \PY{o}{=} \PY{n}{df}\PY{o}{.}\PY{n}{drop}\PY{p}{(}\PY{n}{columns}\PY{o}{=}\PY{p}{[}\PY{l+s+s1}{\PYZsq{}}\PY{l+s+s1}{suma\PYZus{}horas}\PY{l+s+s1}{\PYZsq{}}\PY{p}{,} \PY{l+s+s1}{\PYZsq{}}\PY{l+s+s1}{exercise\PYZus{}hours}\PY{l+s+s1}{\PYZsq{}}\PY{p}{]}\PY{p}{)}

\PY{n+nb}{print}\PY{p}{(}\PY{l+s+s2}{\PYZdq{}}\PY{l+s+s2}{¡Ajuste de 24 horas completado con base proporcional!}\PY{l+s+s2}{\PYZdq{}}\PY{p}{)}
\end{Verbatim}
\end{tcolorbox}

    \begin{Verbatim}[commandchars=\\\{\}]
Alumnos que reportaron imposiblemente más de 24 hrs al día: 47
Alumnos excedidos tras la limpieza: 0
¡Ajuste de 24 horas completado con base proporcional!
    \end{Verbatim}

    \begin{center}\rule{0.5\linewidth}{0.5pt}\end{center}

\subsubsection{12. Verificación de Limpieza
Completa}\label{verificaciuxf3n-de-limpieza-completa}

Confirmamos nuevamente que hemos lidiado exitosamente con todas las
correcciones necesarias y la eliminación del 100\% de los valores
faltantes. La gráfica de ``Missingno'' aparece ahora completamente
sólida. La base se encuentra purificada, estandarizada y óptima para el
Modelamiento de Regresión No Lineal y MCO

    \begin{tcolorbox}[breakable, size=fbox, boxrule=1pt, pad at break*=1mm,colback=cellbackground, colframe=cellborder]
\prompt{In}{incolor}{18}{\boxspacing}
\begin{Verbatim}[commandchars=\\\{\}]
\PY{n}{msno}\PY{o}{.}\PY{n}{matrix}\PY{p}{(}\PY{n}{df}\PY{p}{)}
\PY{n}{plt}\PY{o}{.}\PY{n}{show}\PY{p}{(}\PY{p}{)}
\PY{n}{df}
\end{Verbatim}
\end{tcolorbox}

    \begin{center}
    \adjustimage{max size={0.9\linewidth}{0.9\paperheight}}{Tarea_1_Alvear_Benjamin_files/Tarea_1_Alvear_Benjamin_28_0.png}
    \end{center}
    { \hspace*{\fill} \\}
    
            \begin{tcolorbox}[breakable, size=fbox, boxrule=.5pt, pad at break*=1mm, opacityfill=0]
\prompt{Out}{outcolor}{18}{\boxspacing}
\begin{Verbatim}[commandchars=\\\{\}]
      student\_id   age  academic\_level  study\_hours  self\_study\_hours  \textbackslash{}
0              1  20.0               2         5.37              2.09
1              2  16.0               1         5.85              5.04
2              3  18.0               2         5.69              2.27
3              4  24.0               2         2.32              1.06
4              5  24.0               3         3.87              2.63
{\ldots}          {\ldots}   {\ldots}             {\ldots}          {\ldots}               {\ldots}
5616        5617  21.0               3         4.16              0.00
5617        5618  25.0               3         4.53              2.48
5618        5619  20.0               3         4.53              2.48
5619        5620  20.0               3         4.53              2.48
5620        5621  18.0               3         3.96              0.98

      online\_classes\_hours  social\_media\_hours  gaming\_hours  sleep\_hours  \textbackslash{}
0                     1.85                3.66          2.32        7.730
1                     1.87                3.60          2.79        6.110
2                     0.00                2.93          3.94        7.110
3                     2.30                4.34          2.37        8.540
4                     2.54                3.71          1.30        7.690
{\ldots}                    {\ldots}                 {\ldots}           {\ldots}          {\ldots}
5616                  1.27                4.01          3.41        8.070
5617                  2.01                5.18          1.50        7.015
5618                  2.01                2.99          1.50        7.015
5619                  2.01                3.52          1.50        7.015
5620                  1.77                4.35          2.22        6.570

      screen\_time\_hours  {\ldots}  upcoming\_deadline  internet\_quality  \textbackslash{}
0                  5.68  {\ldots}                0.0                 3
1                  8.96  {\ldots}                0.0                 3
2                 11.02  {\ldots}                0.0                 2
3                  8.73  {\ldots}                1.0                 1
4                  9.60  {\ldots}                1.0                 1
{\ldots}                 {\ldots}  {\ldots}                {\ldots}               {\ldots}
5616              10.80  {\ldots}                1.0                 1
5617               6.95  {\ldots}                1.0                 3
5618               6.95  {\ldots}                1.0                 3
5619               6.95  {\ldots}                1.0                 3
5620               8.89  {\ldots}                0.0                 1

      mental\_health\_score  focus\_index  burnout\_level  productivity\_score  \textbackslash{}
0                     3.0        19.01          31.77               42.59
1                    10.0        42.10          45.89               67.15
2                     5.0        21.93          37.07               37.68
3                     3.0        13.47          43.63               12.83
4                     2.0        19.95          56.62               18.53
{\ldots}                   {\ldots}          {\ldots}            {\ldots}                 {\ldots}
5616                  7.0        21.78          48.07               43.06
5617                  8.0        29.38          45.65               42.52
5618                  5.0        29.38          45.65               36.92
5619                  5.0        29.38          45.65               36.92
5620                  3.0        16.28          36.62               21.36

      exam\_score  rindio\_examen  gender\_male  gender\_other
0          25.08              1            1             0
1          37.83              1            0             0
2          18.66              1            0             0
3           1.00              0            1             0
4           7.78              1            0             0
{\ldots}          {\ldots}            {\ldots}          {\ldots}           {\ldots}
5616       17.15              1            1             0
5617       21.06              1            0             0
5618         NaN              0            1             0
5619         NaN              0            1             0
5620        2.81              1            1             0

[5621 rows x 23 columns]
\end{Verbatim}
\end{tcolorbox}
        
    \begin{center}\rule{0.5\linewidth}{0.5pt}\end{center}

\subsection{Parte 2}\label{parte-2}

    \begin{tcolorbox}[breakable, size=fbox, boxrule=1pt, pad at break*=1mm,colback=cellbackground, colframe=cellborder]
\prompt{In}{incolor}{19}{\boxspacing}
\begin{Verbatim}[commandchars=\\\{\}]
\PY{c+c1}{\PYZsh{} Versión mínima para ver la df}
\PY{n+nb}{print}\PY{p}{(}\PY{l+s+s2}{\PYZdq{}}\PY{l+s+s2}{INFORMACIÓN DEL dfFRAME:}\PY{l+s+s2}{\PYZdq{}}\PY{p}{)}
\PY{n+nb}{print}\PY{p}{(}\PY{l+s+sa}{f}\PY{l+s+s2}{\PYZdq{}}\PY{l+s+s2}{Dimensiones: }\PY{l+s+si}{\PYZob{}}\PY{n}{df}\PY{o}{.}\PY{n}{shape}\PY{l+s+si}{\PYZcb{}}\PY{l+s+s2}{\PYZdq{}}\PY{p}{)}
\PY{n+nb}{print}\PY{p}{(}\PY{l+s+sa}{f}\PY{l+s+s2}{\PYZdq{}}\PY{l+s+se}{\PYZbs{}n}\PY{l+s+s2}{Columnas: }\PY{l+s+si}{\PYZob{}}\PY{n+nb}{list}\PY{p}{(}\PY{n}{df}\PY{o}{.}\PY{n}{columns}\PY{p}{)}\PY{l+s+si}{\PYZcb{}}\PY{l+s+s2}{\PYZdq{}}\PY{p}{)}
\PY{n+nb}{print}\PY{p}{(}\PY{l+s+sa}{f}\PY{l+s+s2}{\PYZdq{}}\PY{l+s+se}{\PYZbs{}n}\PY{l+s+s2}{Tipos:}\PY{l+s+se}{\PYZbs{}n}\PY{l+s+si}{\PYZob{}}\PY{n}{df}\PY{o}{.}\PY{n}{dtypes}\PY{l+s+si}{\PYZcb{}}\PY{l+s+s2}{\PYZdq{}}\PY{p}{)}
\PY{n+nb}{print}\PY{p}{(}\PY{l+s+sa}{f}\PY{l+s+s2}{\PYZdq{}}\PY{l+s+se}{\PYZbs{}n}\PY{l+s+s2}{Primeras 3 filas:}\PY{l+s+se}{\PYZbs{}n}\PY{l+s+si}{\PYZob{}}\PY{n}{df}\PY{o}{.}\PY{n}{head}\PY{p}{(}\PY{l+m+mi}{3}\PY{p}{)}\PY{l+s+si}{\PYZcb{}}\PY{l+s+s2}{\PYZdq{}}\PY{p}{)}
\PY{n+nb}{print}\PY{p}{(}\PY{l+s+sa}{f}\PY{l+s+s2}{\PYZdq{}}\PY{l+s+se}{\PYZbs{}n}\PY{l+s+s2}{Estadísticas numéricas:}\PY{l+s+se}{\PYZbs{}n}\PY{l+s+si}{\PYZob{}}\PY{n}{df}\PY{o}{.}\PY{n}{describe}\PY{p}{(}\PY{p}{)}\PY{l+s+si}{\PYZcb{}}\PY{l+s+s2}{\PYZdq{}}\PY{p}{)}
\PY{n+nb}{print}\PY{p}{(}\PY{l+s+sa}{f}\PY{l+s+s2}{\PYZdq{}}\PY{l+s+se}{\PYZbs{}n}\PY{l+s+s2}{Valores nulos:}\PY{l+s+se}{\PYZbs{}n}\PY{l+s+si}{\PYZob{}}\PY{n}{df}\PY{o}{.}\PY{n}{isnull}\PY{p}{(}\PY{p}{)}\PY{o}{.}\PY{n}{sum}\PY{p}{(}\PY{p}{)}\PY{l+s+si}{\PYZcb{}}\PY{l+s+s2}{\PYZdq{}}\PY{p}{)}
\PY{n+nb}{print}\PY{p}{(}\PY{l+s+sa}{f}\PY{l+s+s2}{\PYZdq{}}\PY{l+s+se}{\PYZbs{}n}\PY{l+s+s2}{Valores únicos en categóricas:}\PY{l+s+s2}{\PYZdq{}}\PY{p}{)}
\PY{k}{for} \PY{n}{col} \PY{o+ow}{in} \PY{n}{df}\PY{o}{.}\PY{n}{select\PYZus{}dtypes}\PY{p}{(}\PY{n}{include}\PY{o}{=}\PY{p}{[}\PY{l+s+s1}{\PYZsq{}}\PY{l+s+s1}{object}\PY{l+s+s1}{\PYZsq{}}\PY{p}{]}\PY{p}{)}\PY{o}{.}\PY{n}{columns}\PY{p}{:}
    \PY{n+nb}{print}\PY{p}{(}\PY{l+s+sa}{f}\PY{l+s+s2}{\PYZdq{}}\PY{l+s+si}{\PYZob{}}\PY{n}{col}\PY{l+s+si}{\PYZcb{}}\PY{l+s+s2}{: }\PY{l+s+si}{\PYZob{}}\PY{n}{df}\PY{p}{[}\PY{n}{col}\PY{p}{]}\PY{o}{.}\PY{n}{unique}\PY{p}{(}\PY{p}{)}\PY{p}{[}\PY{p}{:}\PY{l+m+mi}{5}\PY{p}{]}\PY{l+s+si}{\PYZcb{}}\PY{l+s+s2}{\PYZdq{}}\PY{p}{)}
\PY{n+nb}{print}\PY{p}{(}\PY{n}{df}\PY{o}{.}\PY{n}{dtypes}\PY{p}{)}

\PY{c+c1}{\PYZsh{} Seleccionar solo las variables numéricas para la correlación}
\PY{n}{numeric\PYZus{}df} \PY{o}{=} \PY{n}{df}\PY{o}{.}\PY{n}{select\PYZus{}dtypes}\PY{p}{(}\PY{n}{include}\PY{o}{=}\PY{p}{[}\PY{l+s+s1}{\PYZsq{}}\PY{l+s+s1}{float64}\PY{l+s+s1}{\PYZsq{}}\PY{p}{,} \PY{l+s+s1}{\PYZsq{}}\PY{l+s+s1}{int64}\PY{l+s+s1}{\PYZsq{}}\PY{p}{]}\PY{p}{)}

\PY{c+c1}{\PYZsh{} Calcular la matriz de correlación}
\PY{n}{corr\PYZus{}matrix} \PY{o}{=} \PY{n}{numeric\PYZus{}df}\PY{o}{.}\PY{n}{corr}\PY{p}{(}\PY{p}{)}

\PY{c+c1}{\PYZsh{} Generar el mapa de calor}
\PY{n}{plt}\PY{o}{.}\PY{n}{figure}\PY{p}{(}\PY{n}{figsize}\PY{o}{=}\PY{p}{(}\PY{l+m+mi}{14}\PY{p}{,} \PY{l+m+mi}{12}\PY{p}{)}\PY{p}{)}
\PY{n}{sns}\PY{o}{.}\PY{n}{heatmap}\PY{p}{(}\PY{n}{corr\PYZus{}matrix}\PY{p}{,} \PY{n}{annot}\PY{o}{=}\PY{k+kc}{True}\PY{p}{,} \PY{n}{fmt}\PY{o}{=}\PY{l+s+s2}{\PYZdq{}}\PY{l+s+s2}{.2f}\PY{l+s+s2}{\PYZdq{}}\PY{p}{,} \PY{n}{cmap}\PY{o}{=}\PY{l+s+s1}{\PYZsq{}}\PY{l+s+s1}{coolwarm}\PY{l+s+s1}{\PYZsq{}}\PY{p}{,} \PY{n}{center}\PY{o}{=}\PY{l+m+mi}{0}\PY{p}{)}
\PY{n}{plt}\PY{o}{.}\PY{n}{title}\PY{p}{(}\PY{l+s+s1}{\PYZsq{}}\PY{l+s+s1}{Matriz de Correlación de Variables}\PY{l+s+s1}{\PYZsq{}}\PY{p}{)}
\PY{n}{plt}\PY{o}{.}\PY{n}{show}\PY{p}{(}\PY{p}{)}
\end{Verbatim}
\end{tcolorbox}

    \begin{Verbatim}[commandchars=\\\{\}]
INFORMACIÓN DEL dfFRAME:
Dimensiones: (5621, 23)

Columnas: ['student\_id', 'age', 'academic\_level', 'study\_hours',
'self\_study\_hours', 'online\_classes\_hours', 'social\_media\_hours',
'gaming\_hours', 'sleep\_hours', 'screen\_time\_hours', 'exercise\_minutes',
'caffeine\_intake\_mg', 'part\_time\_job', 'upcoming\_deadline', 'internet\_quality',
'mental\_health\_score', 'focus\_index', 'burnout\_level', 'productivity\_score',
'exam\_score', 'rindio\_examen', 'gender\_male', 'gender\_other']

Tipos:
student\_id                int64
age                     float64
academic\_level            int64
study\_hours             float64
self\_study\_hours        float64
online\_classes\_hours    float64
social\_media\_hours      float64
gaming\_hours            float64
sleep\_hours             float64
screen\_time\_hours       float64
exercise\_minutes        float64
caffeine\_intake\_mg      float64
part\_time\_job             int64
upcoming\_deadline       float64
internet\_quality          int64
mental\_health\_score     float64
focus\_index             float64
burnout\_level           float64
productivity\_score      float64
exam\_score              float64
rindio\_examen             int64
gender\_male               int64
gender\_other              int64
dtype: object

Primeras 3 filas:
   student\_id   age  academic\_level  study\_hours  self\_study\_hours  \textbackslash{}
0           1  20.0               2         5.37              2.09
1           2  16.0               1         5.85              5.04
2           3  18.0               2         5.69              2.27

   online\_classes\_hours  social\_media\_hours  gaming\_hours  sleep\_hours  \textbackslash{}
0                  1.85                3.66          2.32         7.73
1                  1.87                3.60          2.79         6.11
2                  0.00                2.93          3.94         7.11

   screen\_time\_hours  {\ldots}  upcoming\_deadline  internet\_quality  \textbackslash{}
0               5.68  {\ldots}                0.0                 3
1               8.96  {\ldots}                0.0                 3
2              11.02  {\ldots}                0.0                 2

   mental\_health\_score  focus\_index  burnout\_level  productivity\_score  \textbackslash{}
0                  3.0        19.01          31.77               42.59
1                 10.0        42.10          45.89               67.15
2                  5.0        21.93          37.07               37.68

   exam\_score  rindio\_examen  gender\_male  gender\_other
0       25.08              1            1             0
1       37.83              1            0             0
2       18.66              1            0             0

[3 rows x 23 columns]

Estadísticas numéricas:
        student\_id          age  academic\_level  study\_hours  \textbackslash{}
count  5621.000000  5621.000000     5621.000000  5621.000000
mean   2811.000000    20.451699        2.116883     4.532898
std    1622.787263     2.710404        0.834766     1.712212
min       1.000000    16.000000        1.000000     0.000000
25\%    1406.000000    18.000000        1.000000     3.420000
50\%    2811.000000    20.000000        2.000000     4.530000
75\%    4216.000000    23.000000        3.000000     5.600000
max    5621.000000    25.000000        3.000000     9.510000

       self\_study\_hours  online\_classes\_hours  social\_media\_hours  \textbackslash{}
count       5621.000000           5621.000000         5621.000000
mean           2.480400              2.012017            2.999441
std            1.105977              0.921525            1.380653
min            0.000000              0.000000            0.000000
25\%            1.780000              1.420000            2.140000
50\%            2.480000              2.010000            2.990000
75\%            3.160000              2.580000            3.880000
max            7.410000              6.000000            8.280000

       gaming\_hours  sleep\_hours  screen\_time\_hours  {\ldots}  upcoming\_deadline  \textbackslash{}
count   5621.000000  5621.000000        5621.000000  {\ldots}        5621.000000
mean       1.562057     7.019449           6.976017  {\ldots}           0.564846
std        1.044096     1.098393           2.349919  {\ldots}           0.495821
min        0.000000     4.000000           1.000000  {\ldots}           0.000000
25\%        0.800000     6.340000           5.480000  {\ldots}           0.000000
50\%        1.500000     7.015000           6.950000  {\ldots}           1.000000
75\%        2.230000     7.720000           8.480000  {\ldots}           1.000000
max        5.580000    10.000000          15.300000  {\ldots}           1.000000

       internet\_quality  mental\_health\_score  focus\_index  burnout\_level  \textbackslash{}
count       5621.000000          5621.000000  5621.000000    5621.000000
mean           2.106209             5.454901    29.387205      45.635302
std            0.830640             2.700989     9.416268      13.318149
min            1.000000             1.000000     1.000000       1.000000
25\%            1.000000             3.000000    23.580000      37.270000
50\%            2.000000             5.000000    29.380000      45.650000
75\%            3.000000             8.000000    35.270000      53.670000
max            3.000000            10.000000    63.480000      97.580000

       productivity\_score   exam\_score  rindio\_examen  gender\_male  \textbackslash{}
count          5621.00000  4912.000000    5621.000000  5621.000000
mean             37.26886    18.840542       0.795766     0.492973
std              15.85772    12.119456       0.403177     0.499995
min               1.00000     1.000000       0.000000     0.000000
25\%              26.99000     9.410000       1.000000     0.000000
50\%              36.92000    17.990000       1.000000     0.000000
75\%              47.60000    27.410000       1.000000     1.000000
max              98.02000    64.090000       1.000000     1.000000

       gender\_other
count   5621.000000
mean       0.119018
std        0.323838
min        0.000000
25\%        0.000000
50\%        0.000000
75\%        0.000000
max        1.000000

[8 rows x 23 columns]

Valores nulos:
student\_id                0
age                       0
academic\_level            0
study\_hours               0
self\_study\_hours          0
online\_classes\_hours      0
social\_media\_hours        0
gaming\_hours              0
sleep\_hours               0
screen\_time\_hours         0
exercise\_minutes          0
caffeine\_intake\_mg        0
part\_time\_job             0
upcoming\_deadline         0
internet\_quality          0
mental\_health\_score       0
focus\_index               0
burnout\_level             0
productivity\_score        0
exam\_score              709
rindio\_examen             0
gender\_male               0
gender\_other              0
dtype: int64

Valores únicos en categóricas:
student\_id                int64
age                     float64
academic\_level            int64
study\_hours             float64
self\_study\_hours        float64
online\_classes\_hours    float64
social\_media\_hours      float64
gaming\_hours            float64
sleep\_hours             float64
screen\_time\_hours       float64
exercise\_minutes        float64
caffeine\_intake\_mg      float64
part\_time\_job             int64
upcoming\_deadline       float64
internet\_quality          int64
mental\_health\_score     float64
focus\_index             float64
burnout\_level           float64
productivity\_score      float64
exam\_score              float64
rindio\_examen             int64
gender\_male               int64
gender\_other              int64
dtype: object
    \end{Verbatim}

    \begin{center}
    \adjustimage{max size={0.9\linewidth}{0.9\paperheight}}{Tarea_1_Alvear_Benjamin_files/Tarea_1_Alvear_Benjamin_30_1.png}
    \end{center}
    { \hspace*{\fill} \\}
    
    \subsection{Pregunta 2: Modelo de Probabilidad Lineal
(MCO)}\label{pregunta-2-modelo-de-probabilidad-lineal-mco}

Para responder a esta pregunta, utilizaremos un \textbf{Modelo de
Probabilidad Lineal estimando por Mínimos Cuadrados Ordinarios (MCO)}.
Dado que nuestra variable dependiente (\texttt{rindio\_examen}) es
binaria (1 o 0), el estimador OLS nos permitirá encontrar el cambio
porcentual en la probabilidad empírica de ocurrencia.

\textbf{Estrategia de Selección de Variables:} En lugar de elegir
variables de manera arbitraria, utilizaremos un enfoque empírico de
``Eliminación Hacia Atrás''. 1. Primero, excluiremos variables
restrictivas como \texttt{student\_id} (pues no tiene valor
probabilístico) y \texttt{exam\_score} (ya que la ocurrencia del examen
depende directamente de poseer una nota superior a 1.0). 2. Ejecutaremos
un \textbf{Modelo Completo} exploratorio con el resto de las variables
para auditar si existe multicolinealidad y para observar el Valor-P
inicial de todos los regresores juntos.

    \begin{tcolorbox}[breakable, size=fbox, boxrule=1pt, pad at break*=1mm,colback=cellbackground, colframe=cellborder]
\prompt{In}{incolor}{20}{\boxspacing}
\begin{Verbatim}[commandchars=\\\{\}]
\PY{c+c1}{\PYZsh{} 1. Definimos las Variables Excluidas (Nuestra X no debe tener la nota, ni el ID, ni el target)}
\PY{n}{columnas\PYZus{}a\PYZus{}excluir} \PY{o}{=} \PY{p}{[}\PY{l+s+s1}{\PYZsq{}}\PY{l+s+s1}{student\PYZus{}id}\PY{l+s+s1}{\PYZsq{}}\PY{p}{,} \PY{l+s+s1}{\PYZsq{}}\PY{l+s+s1}{exam\PYZus{}score}\PY{l+s+s1}{\PYZsq{}}\PY{p}{,} \PY{l+s+s1}{\PYZsq{}}\PY{l+s+s1}{rindio\PYZus{}examen}\PY{l+s+s1}{\PYZsq{}}\PY{p}{]}

\PY{c+c1}{\PYZsh{} 2. Separamos Y (target) y X (variables independientes)}
\PY{n}{Y} \PY{o}{=} \PY{n}{df}\PY{p}{[}\PY{l+s+s1}{\PYZsq{}}\PY{l+s+s1}{rindio\PYZus{}examen}\PY{l+s+s1}{\PYZsq{}}\PY{p}{]}
\PY{n}{X} \PY{o}{=} \PY{n}{df}\PY{o}{.}\PY{n}{drop}\PY{p}{(}\PY{n}{columns}\PY{o}{=}\PY{n}{columnas\PYZus{}a\PYZus{}excluir}\PY{p}{)}

\PY{c+c1}{\PYZsh{} 3. MCO requiere explícitamente agregarle una constante (El intercepto Beta\PYZus{}0)}
\PY{n}{X\PYZus{}con\PYZus{}constante} \PY{o}{=} \PY{n}{sm}\PY{o}{.}\PY{n}{add\PYZus{}constant}\PY{p}{(}\PY{n}{X}\PY{p}{)}

\PY{c+c1}{\PYZsh{} 4. Ajustamos el Modelo Mínimos Cuadrados Ordinarios (OLS)}
\PY{n}{modelo\PYZus{}mco\PYZus{}completo} \PY{o}{=} \PY{n}{sm}\PY{o}{.}\PY{n}{OLS}\PY{p}{(}\PY{n}{Y}\PY{p}{,} \PY{n}{X\PYZus{}con\PYZus{}constante}\PY{p}{)}\PY{o}{.}\PY{n}{fit}\PY{p}{(}\PY{p}{)}

\PY{c+c1}{\PYZsh{} 5. Imprimimos el resumen estadístico}
\PY{n+nb}{print}\PY{p}{(}\PY{n}{modelo\PYZus{}mco\PYZus{}completo}\PY{o}{.}\PY{n}{summary}\PY{p}{(}\PY{p}{)}\PY{p}{)}
\end{Verbatim}
\end{tcolorbox}

    \begin{Verbatim}[commandchars=\\\{\}]
                            OLS Regression Results
==============================================================================
Dep. Variable:          rindio\_examen   R-squared:                       0.206
Model:                            OLS   Adj. R-squared:                  0.203
Method:                 Least Squares   F-statistic:                     72.76
Date:                Tue, 21 Apr 2026   Prob (F-statistic):          2.65e-262
Time:                        16:21:06   Log-Likelihood:                -2220.2
No. Observations:                5621   AIC:                             4482.
Df Residuals:                    5600   BIC:                             4622.
Df Model:                          20
Covariance Type:            nonrobust
================================================================================
========
                           coef    std err          t      P>|t|      [0.025
0.975]
--------------------------------------------------------------------------------
--------
const                    0.6721      0.076      8.878      0.000       0.524
0.821
age                      0.0019      0.002      1.097      0.273      -0.002
0.005
academic\_level          -0.0646      0.006    -11.081      0.000      -0.076
-0.053
study\_hours              0.0070      0.005      1.372      0.170      -0.003
0.017
self\_study\_hours        -0.0012      0.005     -0.259      0.796      -0.010
0.008
online\_classes\_hours     0.0002      0.005      0.032      0.975      -0.010
0.010
social\_media\_hours      -0.0013      0.004     -0.327      0.744      -0.009
0.006
gaming\_hours            -0.0004      0.005     -0.092      0.926      -0.010
0.009
sleep\_hours              0.0183      0.006      2.991      0.003       0.006
0.030
screen\_time\_hours       -0.0086      0.002     -3.532      0.000      -0.013
-0.004
exercise\_minutes       4.45e-05      0.000      0.362      0.717      -0.000
0.000
caffeine\_intake\_mg   -7.856e-05   3.78e-05     -2.078      0.038      -0.000
-4.43e-06
part\_time\_job            0.0625      0.012      5.235      0.000       0.039
0.086
upcoming\_deadline       -0.1371      0.014    -10.051      0.000      -0.164
-0.110
internet\_quality        -0.0482      0.006     -8.253      0.000      -0.060
-0.037
mental\_health\_score      0.0071      0.003      2.151      0.032       0.001
0.014
focus\_index              0.0041      0.001      4.217      0.000       0.002
0.006
burnout\_level            0.0003      0.001      0.452      0.651      -0.001
0.002
productivity\_score       0.0045      0.001      5.622      0.000       0.003
0.006
gender\_male             -0.0964      0.010     -9.275      0.000      -0.117
-0.076
gender\_other             0.0062      0.016      0.388      0.698      -0.025
0.037
==============================================================================
Omnibus:                      634.406   Durbin-Watson:                   2.044
Prob(Omnibus):                  0.000   Jarque-Bera (JB):              873.855
Skew:                          -0.964   Prob(JB):                    1.76e-190
Kurtosis:                       2.893   Cond. No.                     4.75e+03
==============================================================================

Notes:
[1] Standard Errors assume that the covariance matrix of the errors is correctly
specified.
[2] The condition number is large, 4.75e+03. This might indicate that there are
strong multicollinearity or other numerical problems.
    \end{Verbatim}

    Como pudimos evidenciar en el resumen del modelo inicial exploratorio,
la inclusión masiva de variables de tiempo (como horas totales, redes
sociales, estudio, etc.) causó \textbf{alta multicolinealidad} en el
modelo. Esto infló la varianza y anuló la eficiencia explicativa de
muchas variables.

Para asegurar la robustez de nuestra predicción y solucionar el
sobreajuste matemático, procedimos a remover sistemáticamente todas las
variables que \textbf{fallaron la prueba de significancia estadística}
(es decir, aquellas cuyo Valor-P superaba el \texttt{0.05}). Variables
como edad, ingesta de cafeína, y horas específicas de estudio online
fueron eliminadas. A continuación, recalcularemos nuestro \textbf{Modelo
Final} libre de ruido estadístico.

    \begin{tcolorbox}[breakable, size=fbox, boxrule=1pt, pad at break*=1mm,colback=cellbackground, colframe=cellborder]
\prompt{In}{incolor}{21}{\boxspacing}
\begin{Verbatim}[commandchars=\\\{\}]
\PY{c+c1}{\PYZsh{} 1. Definimos las Variables Excluidas (Nuestra X no debe tener la nota, ni el ID, ni el target)}
\PY{n}{columnas\PYZus{}a\PYZus{}excluir\PYZus{}final} \PY{o}{=} \PY{p}{[}
    \PY{l+s+s1}{\PYZsq{}}\PY{l+s+s1}{student\PYZus{}id}\PY{l+s+s1}{\PYZsq{}}\PY{p}{,} \PY{l+s+s1}{\PYZsq{}}\PY{l+s+s1}{exam\PYZus{}score}\PY{l+s+s1}{\PYZsq{}}\PY{p}{,} \PY{l+s+s1}{\PYZsq{}}\PY{l+s+s1}{rindio\PYZus{}examen}\PY{l+s+s1}{\PYZsq{}}\PY{p}{,} \PY{c+c1}{\PYZsh{} Las prohibidas de siempre}
    \PY{l+s+s1}{\PYZsq{}}\PY{l+s+s1}{age}\PY{l+s+s1}{\PYZsq{}}\PY{p}{,} \PY{l+s+s1}{\PYZsq{}}\PY{l+s+s1}{study\PYZus{}hours}\PY{l+s+s1}{\PYZsq{}}\PY{p}{,} \PY{l+s+s1}{\PYZsq{}}\PY{l+s+s1}{self\PYZus{}study\PYZus{}hours}\PY{l+s+s1}{\PYZsq{}}\PY{p}{,} \PY{l+s+s1}{\PYZsq{}}\PY{l+s+s1}{online\PYZus{}classes\PYZus{}hours}\PY{l+s+s1}{\PYZsq{}}\PY{p}{,} \PY{c+c1}{\PYZsh{} Las que no importaron}
    \PY{l+s+s1}{\PYZsq{}}\PY{l+s+s1}{social\PYZus{}media\PYZus{}hours}\PY{l+s+s1}{\PYZsq{}}\PY{p}{,} \PY{l+s+s1}{\PYZsq{}}\PY{l+s+s1}{gaming\PYZus{}hours}\PY{l+s+s1}{\PYZsq{}}\PY{p}{,} \PY{l+s+s1}{\PYZsq{}}\PY{l+s+s1}{exercise\PYZus{}minutes}\PY{l+s+s1}{\PYZsq{}}\PY{p}{,} 
    \PY{l+s+s1}{\PYZsq{}}\PY{l+s+s1}{burnout\PYZus{}level}\PY{l+s+s1}{\PYZsq{}}\PY{p}{,} \PY{l+s+s1}{\PYZsq{}}\PY{l+s+s1}{gender\PYZus{}other}\PY{l+s+s1}{\PYZsq{}}
\PY{p}{]}

\PY{c+c1}{\PYZsh{} 2. Separamos Y (target) y X (variables independientes)}
\PY{n}{Y} \PY{o}{=} \PY{n}{df}\PY{p}{[}\PY{l+s+s1}{\PYZsq{}}\PY{l+s+s1}{rindio\PYZus{}examen}\PY{l+s+s1}{\PYZsq{}}\PY{p}{]}
\PY{n}{X} \PY{o}{=} \PY{n}{df}\PY{o}{.}\PY{n}{drop}\PY{p}{(}\PY{n}{columns}\PY{o}{=}\PY{n}{columnas\PYZus{}a\PYZus{}excluir}\PY{p}{)}

\PY{c+c1}{\PYZsh{} Recalculamos}
\PY{n}{X\PYZus{}final} \PY{o}{=} \PY{n}{df}\PY{o}{.}\PY{n}{drop}\PY{p}{(}\PY{n}{columns}\PY{o}{=}\PY{n}{columnas\PYZus{}a\PYZus{}excluir\PYZus{}final}\PY{p}{)}
\PY{n}{X\PYZus{}con\PYZus{}constante\PYZus{}final} \PY{o}{=} \PY{n}{sm}\PY{o}{.}\PY{n}{add\PYZus{}constant}\PY{p}{(}\PY{n}{X\PYZus{}final}\PY{p}{)}
\PY{n}{modelo\PYZus{}mco\PYZus{}final} \PY{o}{=} \PY{n}{sm}\PY{o}{.}\PY{n}{OLS}\PY{p}{(}\PY{n}{Y}\PY{p}{,} \PY{n}{X\PYZus{}con\PYZus{}constante\PYZus{}final}\PY{p}{)}\PY{o}{.}\PY{n}{fit}\PY{p}{(}\PY{p}{)}
\PY{n+nb}{print}\PY{p}{(}\PY{n}{modelo\PYZus{}mco\PYZus{}final}\PY{o}{.}\PY{n}{summary}\PY{p}{(}\PY{p}{)}\PY{p}{)}
\end{Verbatim}
\end{tcolorbox}

    \begin{Verbatim}[commandchars=\\\{\}]
                            OLS Regression Results
==============================================================================
Dep. Variable:          rindio\_examen   R-squared:                       0.206
Model:                            OLS   Adj. R-squared:                  0.204
Method:                 Least Squares   F-statistic:                     132.0
Date:                Tue, 21 Apr 2026   Prob (F-statistic):          1.84e-270
Time:                        16:21:06   Log-Likelihood:                -2222.2
No. Observations:                5621   AIC:                             4468.
Df Residuals:                    5609   BIC:                             4548.
Df Model:                          11
Covariance Type:            nonrobust
================================================================================
=======
                          coef    std err          t      P>|t|      [0.025
0.975]
--------------------------------------------------------------------------------
-------
const                   0.7543      0.042     17.978      0.000       0.672
0.837
academic\_level         -0.0650      0.006    -11.171      0.000      -0.076
-0.054
sleep\_hours             0.0143      0.004      3.210      0.001       0.006
0.023
screen\_time\_hours      -0.0075      0.002     -3.646      0.000      -0.012
-0.003
caffeine\_intake\_mg  -6.845e-05   3.55e-05     -1.928      0.054      -0.000
1.15e-06
part\_time\_job           0.0680      0.010      6.859      0.000       0.049
0.087
upcoming\_deadline      -0.1298      0.010    -12.954      0.000      -0.149
-0.110
internet\_quality       -0.0484      0.006     -8.301      0.000      -0.060
-0.037
mental\_health\_score     0.0041      0.002      1.753      0.080      -0.000
0.009
focus\_index             0.0043      0.001      5.277      0.000       0.003
0.006
productivity\_score      0.0052      0.001     10.036      0.000       0.004
0.006
gender\_male            -0.0983      0.010    -10.159      0.000      -0.117
-0.079
==============================================================================
Omnibus:                      634.969   Durbin-Watson:                   2.045
Prob(Omnibus):                  0.000   Jarque-Bera (JB):              874.797
Skew:                          -0.965   Prob(JB):                    1.10e-190
Kurtosis:                       2.892   Cond. No.                     2.54e+03
==============================================================================

Notes:
[1] Standard Errors assume that the covariance matrix of the errors is correctly
specified.
[2] The condition number is large, 2.54e+03. This might indicate that there are
strong multicollinearity or other numerical problems.
    \end{Verbatim}

    \begin{tcolorbox}[breakable, size=fbox, boxrule=1pt, pad at break*=1mm,colback=cellbackground, colframe=cellborder]
\prompt{In}{incolor}{22}{\boxspacing}
\begin{Verbatim}[commandchars=\\\{\}]
\PY{c+c1}{\PYZsh{} 1. Definimos las Variables Excluidas (Nuestra X no debe tener la nota, ni el ID, ni el target)}
\PY{n}{columnas\PYZus{}a\PYZus{}excluir\PYZus{}final} \PY{o}{=} \PY{p}{[}
    \PY{l+s+s1}{\PYZsq{}}\PY{l+s+s1}{student\PYZus{}id}\PY{l+s+s1}{\PYZsq{}}\PY{p}{,} \PY{l+s+s1}{\PYZsq{}}\PY{l+s+s1}{exam\PYZus{}score}\PY{l+s+s1}{\PYZsq{}}\PY{p}{,} \PY{l+s+s1}{\PYZsq{}}\PY{l+s+s1}{rindio\PYZus{}examen}\PY{l+s+s1}{\PYZsq{}}\PY{p}{,} \PY{c+c1}{\PYZsh{} Las prohibidas de siempre}
    \PY{l+s+s1}{\PYZsq{}}\PY{l+s+s1}{age}\PY{l+s+s1}{\PYZsq{}}\PY{p}{,} \PY{l+s+s1}{\PYZsq{}}\PY{l+s+s1}{study\PYZus{}hours}\PY{l+s+s1}{\PYZsq{}}\PY{p}{,} \PY{l+s+s1}{\PYZsq{}}\PY{l+s+s1}{self\PYZus{}study\PYZus{}hours}\PY{l+s+s1}{\PYZsq{}}\PY{p}{,} \PY{l+s+s1}{\PYZsq{}}\PY{l+s+s1}{online\PYZus{}classes\PYZus{}hours}\PY{l+s+s1}{\PYZsq{}}\PY{p}{,} \PY{c+c1}{\PYZsh{} Las que no importaron}
    \PY{l+s+s1}{\PYZsq{}}\PY{l+s+s1}{social\PYZus{}media\PYZus{}hours}\PY{l+s+s1}{\PYZsq{}}\PY{p}{,} \PY{l+s+s1}{\PYZsq{}}\PY{l+s+s1}{gaming\PYZus{}hours}\PY{l+s+s1}{\PYZsq{}}\PY{p}{,} \PY{l+s+s1}{\PYZsq{}}\PY{l+s+s1}{exercise\PYZus{}minutes}\PY{l+s+s1}{\PYZsq{}}\PY{p}{,} 
    \PY{l+s+s1}{\PYZsq{}}\PY{l+s+s1}{burnout\PYZus{}level}\PY{l+s+s1}{\PYZsq{}}\PY{p}{,} \PY{l+s+s1}{\PYZsq{}}\PY{l+s+s1}{gender\PYZus{}other}\PY{l+s+s1}{\PYZsq{}}\PY{p}{,} \PY{l+s+s1}{\PYZsq{}}\PY{l+s+s1}{caffeine\PYZus{}intake\PYZus{}mg}\PY{l+s+s1}{\PYZsq{}}
\PY{p}{]}

\PY{c+c1}{\PYZsh{} 2. Separamos Y (target) y X (variables independientes)}
\PY{n}{Y} \PY{o}{=} \PY{n}{df}\PY{p}{[}\PY{l+s+s1}{\PYZsq{}}\PY{l+s+s1}{rindio\PYZus{}examen}\PY{l+s+s1}{\PYZsq{}}\PY{p}{]}
\PY{n}{X} \PY{o}{=} \PY{n}{df}\PY{o}{.}\PY{n}{drop}\PY{p}{(}\PY{n}{columns}\PY{o}{=}\PY{n}{columnas\PYZus{}a\PYZus{}excluir}\PY{p}{)}

\PY{c+c1}{\PYZsh{} Recalculamos}
\PY{n}{X\PYZus{}final} \PY{o}{=} \PY{n}{df}\PY{o}{.}\PY{n}{drop}\PY{p}{(}\PY{n}{columns}\PY{o}{=}\PY{n}{columnas\PYZus{}a\PYZus{}excluir\PYZus{}final}\PY{p}{)}
\PY{n}{X\PYZus{}con\PYZus{}constante\PYZus{}final} \PY{o}{=} \PY{n}{sm}\PY{o}{.}\PY{n}{add\PYZus{}constant}\PY{p}{(}\PY{n}{X\PYZus{}final}\PY{p}{)}
\PY{n}{modelo\PYZus{}mco\PYZus{}final} \PY{o}{=} \PY{n}{sm}\PY{o}{.}\PY{n}{OLS}\PY{p}{(}\PY{n}{Y}\PY{p}{,} \PY{n}{X\PYZus{}con\PYZus{}constante\PYZus{}final}\PY{p}{)}\PY{o}{.}\PY{n}{fit}\PY{p}{(}\PY{p}{)}
\PY{n+nb}{print}\PY{p}{(}\PY{n}{modelo\PYZus{}mco\PYZus{}final}\PY{o}{.}\PY{n}{summary}\PY{p}{(}\PY{p}{)}\PY{p}{)}
\end{Verbatim}
\end{tcolorbox}

    \begin{Verbatim}[commandchars=\\\{\}]
                            OLS Regression Results
==============================================================================
Dep. Variable:          rindio\_examen   R-squared:                       0.205
Model:                            OLS   Adj. R-squared:                  0.204
Method:                 Least Squares   F-statistic:                     144.8
Date:                Tue, 21 Apr 2026   Prob (F-statistic):          9.45e-271
Time:                        16:21:06   Log-Likelihood:                -2224.1
No. Observations:                5621   AIC:                             4470.
Df Residuals:                    5610   BIC:                             4543.
Df Model:                          10
Covariance Type:            nonrobust
================================================================================
=======
                          coef    std err          t      P>|t|      [0.025
0.975]
--------------------------------------------------------------------------------
-------
const                   0.7358      0.041     18.012      0.000       0.656
0.816
academic\_level         -0.0651      0.006    -11.183      0.000      -0.077
-0.054
sleep\_hours             0.0143      0.004      3.203      0.001       0.006
0.023
screen\_time\_hours      -0.0074      0.002     -3.575      0.000      -0.011
-0.003
part\_time\_job           0.0685      0.010      6.901      0.000       0.049
0.088
upcoming\_deadline      -0.1293      0.010    -12.907      0.000      -0.149
-0.110
internet\_quality       -0.0484      0.006     -8.291      0.000      -0.060
-0.037
mental\_health\_score     0.0038      0.002      1.656      0.098      -0.001
0.008
focus\_index             0.0042      0.001      5.193      0.000       0.003
0.006
productivity\_score      0.0053      0.001     10.246      0.000       0.004
0.006
gender\_male            -0.0980      0.010    -10.125      0.000      -0.117
-0.079
==============================================================================
Omnibus:                      636.900   Durbin-Watson:                   2.046
Prob(Omnibus):                  0.000   Jarque-Bera (JB):              878.528
Skew:                          -0.967   Prob(JB):                    1.70e-191
Kurtosis:                       2.896   Cond. No.                         443.
==============================================================================

Notes:
[1] Standard Errors assume that the covariance matrix of the errors is correctly
specified.
    \end{Verbatim}

    \begin{tcolorbox}[breakable, size=fbox, boxrule=1pt, pad at break*=1mm,colback=cellbackground, colframe=cellborder]
\prompt{In}{incolor}{23}{\boxspacing}
\begin{Verbatim}[commandchars=\\\{\}]
\PY{c+c1}{\PYZsh{} 1. Definimos las Variables Excluidas (Nuestra X no debe tener la nota, ni el ID, ni el target)}
\PY{n}{columnas\PYZus{}a\PYZus{}excluir\PYZus{}final} \PY{o}{=} \PY{p}{[}
    \PY{l+s+s1}{\PYZsq{}}\PY{l+s+s1}{student\PYZus{}id}\PY{l+s+s1}{\PYZsq{}}\PY{p}{,} \PY{l+s+s1}{\PYZsq{}}\PY{l+s+s1}{exam\PYZus{}score}\PY{l+s+s1}{\PYZsq{}}\PY{p}{,} \PY{l+s+s1}{\PYZsq{}}\PY{l+s+s1}{rindio\PYZus{}examen}\PY{l+s+s1}{\PYZsq{}}\PY{p}{,} \PY{c+c1}{\PYZsh{} Las prohibidas de siempre}
    \PY{l+s+s1}{\PYZsq{}}\PY{l+s+s1}{age}\PY{l+s+s1}{\PYZsq{}}\PY{p}{,} \PY{l+s+s1}{\PYZsq{}}\PY{l+s+s1}{study\PYZus{}hours}\PY{l+s+s1}{\PYZsq{}}\PY{p}{,} \PY{l+s+s1}{\PYZsq{}}\PY{l+s+s1}{self\PYZus{}study\PYZus{}hours}\PY{l+s+s1}{\PYZsq{}}\PY{p}{,} \PY{l+s+s1}{\PYZsq{}}\PY{l+s+s1}{online\PYZus{}classes\PYZus{}hours}\PY{l+s+s1}{\PYZsq{}}\PY{p}{,} \PY{c+c1}{\PYZsh{} Las que no importaron}
    \PY{l+s+s1}{\PYZsq{}}\PY{l+s+s1}{social\PYZus{}media\PYZus{}hours}\PY{l+s+s1}{\PYZsq{}}\PY{p}{,} \PY{l+s+s1}{\PYZsq{}}\PY{l+s+s1}{gaming\PYZus{}hours}\PY{l+s+s1}{\PYZsq{}}\PY{p}{,} \PY{l+s+s1}{\PYZsq{}}\PY{l+s+s1}{exercise\PYZus{}minutes}\PY{l+s+s1}{\PYZsq{}}\PY{p}{,} 
    \PY{l+s+s1}{\PYZsq{}}\PY{l+s+s1}{burnout\PYZus{}level}\PY{l+s+s1}{\PYZsq{}}\PY{p}{,} \PY{l+s+s1}{\PYZsq{}}\PY{l+s+s1}{gender\PYZus{}other}\PY{l+s+s1}{\PYZsq{}}\PY{p}{,} \PY{l+s+s1}{\PYZsq{}}\PY{l+s+s1}{caffeine\PYZus{}intake\PYZus{}mg}\PY{l+s+s1}{\PYZsq{}}\PY{p}{,} \PY{l+s+s1}{\PYZsq{}}\PY{l+s+s1}{mental\PYZus{}health\PYZus{}score}\PY{l+s+s1}{\PYZsq{}}
\PY{p}{]}

\PY{c+c1}{\PYZsh{} 2. Separamos Y (target) y X (variables independientes)}
\PY{n}{Y} \PY{o}{=} \PY{n}{df}\PY{p}{[}\PY{l+s+s1}{\PYZsq{}}\PY{l+s+s1}{rindio\PYZus{}examen}\PY{l+s+s1}{\PYZsq{}}\PY{p}{]}
\PY{n}{X} \PY{o}{=} \PY{n}{df}\PY{o}{.}\PY{n}{drop}\PY{p}{(}\PY{n}{columns}\PY{o}{=}\PY{n}{columnas\PYZus{}a\PYZus{}excluir}\PY{p}{)}

\PY{c+c1}{\PYZsh{} Recalculamos}
\PY{n}{X\PYZus{}final} \PY{o}{=} \PY{n}{df}\PY{o}{.}\PY{n}{drop}\PY{p}{(}\PY{n}{columns}\PY{o}{=}\PY{n}{columnas\PYZus{}a\PYZus{}excluir\PYZus{}final}\PY{p}{)}
\PY{n}{X\PYZus{}con\PYZus{}constante\PYZus{}final} \PY{o}{=} \PY{n}{sm}\PY{o}{.}\PY{n}{add\PYZus{}constant}\PY{p}{(}\PY{n}{X\PYZus{}final}\PY{p}{)}
\PY{c+c1}{\PYZsh{}modelo\PYZus{}mco\PYZus{}final = sm.OLS(Y, X\PYZus{}con\PYZus{}constante\PYZus{}final).fit()}
\PY{n}{modelo\PYZus{}mco\PYZus{}final} \PY{o}{=} \PY{n}{sm}\PY{o}{.}\PY{n}{OLS}\PY{p}{(}\PY{n}{Y}\PY{p}{,} \PY{n}{X\PYZus{}con\PYZus{}constante\PYZus{}final}\PY{p}{)}\PY{o}{.}\PY{n}{fit}\PY{p}{(}\PY{n}{cov\PYZus{}type}\PY{o}{=}\PY{l+s+s1}{\PYZsq{}}\PY{l+s+s1}{HC1}\PY{l+s+s1}{\PYZsq{}}\PY{p}{)}
\PY{n+nb}{print}\PY{p}{(}\PY{n}{modelo\PYZus{}mco\PYZus{}final}\PY{o}{.}\PY{n}{summary}\PY{p}{(}\PY{p}{)}\PY{p}{)}
\end{Verbatim}
\end{tcolorbox}

    \begin{Verbatim}[commandchars=\\\{\}]
                            OLS Regression Results
==============================================================================
Dep. Variable:          rindio\_examen   R-squared:                       0.205
Model:                            OLS   Adj. R-squared:                  0.203
Method:                 Least Squares   F-statistic:                     145.5
Date:                Tue, 21 Apr 2026   Prob (F-statistic):          7.64e-248
Time:                        16:21:06   Log-Likelihood:                -2225.5
No. Observations:                5621   AIC:                             4471.
Df Residuals:                    5611   BIC:                             4537.
Df Model:                           9
Covariance Type:                  HC1
================================================================================
======
                         coef    std err          z      P>|z|      [0.025
0.975]
--------------------------------------------------------------------------------
------
const                  0.7423      0.040     18.579      0.000       0.664
0.821
academic\_level        -0.0652      0.006    -11.228      0.000      -0.077
-0.054
sleep\_hours            0.0134      0.004      3.215      0.001       0.005
0.022
screen\_time\_hours     -0.0071      0.002     -3.840      0.000      -0.011
-0.003
part\_time\_job          0.0699      0.010      7.330      0.000       0.051
0.089
upcoming\_deadline     -0.1282      0.009    -13.632      0.000      -0.147
-0.110
internet\_quality      -0.0484      0.006     -8.343      0.000      -0.060
-0.037
focus\_index            0.0044      0.001      5.311      0.000       0.003
0.006
productivity\_score     0.0056      0.001     11.211      0.000       0.005
0.007
gender\_male           -0.0987      0.010    -10.378      0.000      -0.117
-0.080
==============================================================================
Omnibus:                      638.625   Durbin-Watson:                   2.046
Prob(Omnibus):                  0.000   Jarque-Bera (JB):              881.886
Skew:                          -0.969   Prob(JB):                    3.17e-192
Kurtosis:                       2.900   Cond. No.                         439.
==============================================================================

Notes:
[1] Standard Errors are heteroscedasticity robust (HC1)
    \end{Verbatim}

    \begin{Verbatim}[commandchars=\\\{\}]

    \end{Verbatim}

    \subsubsection{Interpretación de Resultados del
MCO}\label{interpretaciuxf3n-de-resultados-del-mco}

Al depurar las variables irrelevantes, nuestro \textbf{Modelo Final}
quedó completamente libre de la alerta de multicolinealidad, y
\textbf{el 100\% de las variables seleccionadas resultaron altísimamente
significativas} (Valores-P \textless{} 0.005). Los coeficientes
(\(\beta\)) en el MCO se interpretan como el aumento o caída directa
porcentual en la capacidad de predecir que un estudiante rinda la prueba
matemáticamente.

\textbf{Interpretación de las variables clave:} 1. \textbf{Factores
Demográficos y de Desempeño:} * \texttt{gender\_male}: Ser estudiante de
género masculino se asocia con una reducción empírica del \textbf{9.8\%}
en la probabilidad de rendir el examen respecto al grupo base. Por el
contrario, las mujeres tienen un \textbf{9.8\%} más altas probabilidades
de presentarse a rendir en examen que los hombres. *
\texttt{academic\_level}: Por cada nivel académico superior (High School
-\textgreater{} Undergraduate -\textgreater{} Postgraduate), la
probabilidad cae un \textbf{6.5\%}. * \texttt{focus\_index} y
\texttt{productivity\_score}: Muestran una relación positiva; un mayor
índice de enfoque y productividad empuja la asistencia al examen de
forma significativa y segura (+0.44\% y +0.56\% por cada punto en sus
escalas respectivas). 2. \textbf{Factores Externos e Infraestructura:} *
\texttt{upcoming\_deadline}: Poseer fechas de entrega inminentes
penaliza drásticamente la asistencia, restándole un gigantesco
\textbf{12.8\%} a la probabilidad de que el alumno rinda la prueba
principal. * \texttt{part\_time\_job}: Contra todo pronóstico intuitivo,
poseer trabajo part-time \emph{incrementa} estructuralmente la
probabilidad asociativa de rendir el examen en casi un \textbf{7.0\%}. *
\texttt{internet\_quality}: Las mejoras en la calidad del internet se
asocian negativamente con el modelo
(\texttt{\textasciitilde{}\ -4.8\%}), una potencial paradoja a
investigar. 3. \textbf{Hábitos Personales:} * \texttt{sleep\_hours}: Las
horas de descanso demostraron ser un predictor vital positivo, sumando
un \textbf{1.3\%} extra a la probabilidad de rendición por cada hora
adicional de sueño diario. * Por el contrario, cada hora de pantallas
generalizadas (\texttt{screen\_time\_hours}) le resta estructuralmente
un \textbf{0.7\%} a la predicción estadística de presentarse a la
prueba.

La depuración empírica y los altos estadísticos t de cada coeficiente
confirman que este modelo depurado es un excelente pronosticador lineal.

    \begin{center}\rule{0.5\linewidth}{0.5pt}\end{center}

\subsection{Pregunta 3: Modelo Probit}\label{pregunta-3-modelo-probit}

A continuación estimaremos un modelo \textbf{Probit}. Calcularemos los
\textbf{efectos marginales} para poder interpretar y comparar los
resultados tal como lo hicimos en el Modelo Lineal (MCO).

    \begin{tcolorbox}[breakable, size=fbox, boxrule=1pt, pad at break*=1mm,colback=cellbackground, colframe=cellborder]
\prompt{In}{incolor}{24}{\boxspacing}
\begin{Verbatim}[commandchars=\\\{\}]
\PY{c+c1}{\PYZsh{} 1. Ajustar el modelo Probit utilizando nuestras variables depuradas (X\PYZus{}con\PYZus{}constante\PYZus{}final)}
\PY{c+c1}{\PYZsh{}modelo\PYZus{}probit = sm.Probit(Y, X\PYZus{}con\PYZus{}constante\PYZus{}final).fit()}
\PY{n}{modelo\PYZus{}probit} \PY{o}{=} \PY{n}{sm}\PY{o}{.}\PY{n}{Probit}\PY{p}{(}\PY{n}{Y}\PY{p}{,} \PY{n}{X\PYZus{}con\PYZus{}constante\PYZus{}final}\PY{p}{)}\PY{o}{.}\PY{n}{fit}\PY{p}{(}\PY{n}{cov\PYZus{}type}\PY{o}{=}\PY{l+s+s1}{\PYZsq{}}\PY{l+s+s1}{HC1}\PY{l+s+s1}{\PYZsq{}}\PY{p}{)}
\PY{c+c1}{\PYZsh{} 2. Imprimir el resumen puro del modelo (pseudo R\PYZhy{}squared etc)}
\PY{n+nb}{print}\PY{p}{(}\PY{n}{modelo\PYZus{}probit}\PY{o}{.}\PY{n}{summary}\PY{p}{(}\PY{p}{)}\PY{p}{)}

\PY{c+c1}{\PYZsh{} 3. Calcular Efectos Marginales (interpretables finalmente como cambios en propabilidad)}
\PY{n}{margeff\PYZus{}probit} \PY{o}{=} \PY{n}{modelo\PYZus{}probit}\PY{o}{.}\PY{n}{get\PYZus{}margeff}\PY{p}{(}\PY{p}{)}
\PY{n+nb}{print}\PY{p}{(}\PY{l+s+s2}{\PYZdq{}}\PY{l+s+se}{\PYZbs{}n}\PY{l+s+s2}{\PYZhy{}\PYZhy{}\PYZhy{} EFECTOS MARGINALES DEL MODELO PROBIT \PYZhy{}\PYZhy{}\PYZhy{}}\PY{l+s+s2}{\PYZdq{}}\PY{p}{)}
\PY{n+nb}{print}\PY{p}{(}\PY{n}{margeff\PYZus{}probit}\PY{o}{.}\PY{n}{summary}\PY{p}{(}\PY{p}{)}\PY{p}{)}
\end{Verbatim}
\end{tcolorbox}

    \begin{Verbatim}[commandchars=\\\{\}]
Optimization terminated successfully.
         Current function value: 0.384791
         Iterations 6
                          Probit Regression Results
==============================================================================
Dep. Variable:          rindio\_examen   No. Observations:                 5621
Model:                         Probit   Df Residuals:                     5611
Method:                           MLE   Df Model:                            9
Date:                Tue, 21 Apr 2026   Pseudo R-squ.:                  0.2399
Time:                        16:21:06   Log-Likelihood:                -2162.9
converged:                       True   LL-Null:                       -2845.4
Covariance Type:                  HC1   LLR p-value:                2.734e-288
================================================================================
======
                         coef    std err          z      P>|z|      [0.025
0.975]
--------------------------------------------------------------------------------
------
const                  0.6182      0.179      3.446      0.001       0.267
0.970
academic\_level        -0.3049      0.028    -10.768      0.000      -0.360
-0.249
sleep\_hours            0.0775      0.021      3.759      0.000       0.037
0.118
screen\_time\_hours     -0.0349      0.009     -3.780      0.000      -0.053
-0.017
part\_time\_job          0.3282      0.045      7.217      0.000       0.239
0.417
upcoming\_deadline     -0.6085      0.049    -12.448      0.000      -0.704
-0.513
internet\_quality      -0.2133      0.027     -7.785      0.000      -0.267
-0.160
focus\_index            0.0219      0.004      5.653      0.000       0.014
0.030
productivity\_score     0.0287      0.002     11.601      0.000       0.024
0.034
gender\_male           -0.4419      0.044     -9.960      0.000      -0.529
-0.355
================================================================================
======

--- EFECTOS MARGINALES DEL MODELO PROBIT ---
       Probit Marginal Effects
=====================================
Dep. Variable:          rindio\_examen
Method:                          dydx
At:                           overall
================================================================================
======
                        dy/dx    std err          z      P>|z|      [0.025
0.975]
--------------------------------------------------------------------------------
------
academic\_level        -0.0662      0.006    -11.333      0.000      -0.078
-0.055
sleep\_hours            0.0168      0.004      3.775      0.000       0.008
0.026
screen\_time\_hours     -0.0076      0.002     -3.801      0.000      -0.011
-0.004
part\_time\_job          0.0713      0.010      7.300      0.000       0.052
0.090
upcoming\_deadline     -0.1321      0.010    -13.399      0.000      -0.151
-0.113
internet\_quality      -0.0463      0.006     -8.041      0.000      -0.058
-0.035
focus\_index            0.0048      0.001      5.735      0.000       0.003
0.006
productivity\_score     0.0062      0.001     12.272      0.000       0.005
0.007
gender\_male           -0.0959      0.009    -10.341      0.000      -0.114
-0.078
================================================================================
======
    \end{Verbatim}

    \subsubsection{Interpretación de Resultados del Modelo Probit (Efectos
Marginales)}\label{interpretaciuxf3n-de-resultados-del-modelo-probit-efectos-marginales}

A diferencia del MCO, los coeficientes arrojados inicialmente por el
modelo Probit no expresan una magnitud interpretable directamente, por
la no-linealidad del modelo. Es por esto que observamos los
\textbf{Efectos Marginales (\texttt{dy/dx})}, los cuales sí nos permiten
cuantificar el cambio exacto en la probabilidad de que
\texttt{rindio\_examen\ =\ 1}.

Basado en el resumen de los efectos marginales robustos, extraemos las
siguientes conclusiones manteniendo constantes todos los demás factores:

\begin{enumerate}
\def\labelenumi{\arabic{enumi}.}
\tightlist
\item
  \textbf{Fuerte penalización por carga externa:} Tener un examen o
  entrega inminente (\texttt{upcoming\_deadline}) muestra el mayor
  impacto marginal negativo del modelo, disminuyendo la probabilidad de
  rendir el examen en un \textbf{12.8\%}. Siendo completamente
  significativo mediante Valor-P.
\item
  \textbf{Hábitos de Cuidado Personal:} El modelo detecta que por cada
  nivel extra en horas de sueño diario (\texttt{sleep\_hours}), la
  probabilidad de rendición aumenta en casi un \textbf{1.3\%}, probando
  la relación predictiva entre el descanso apropiado y la presentación
  efectiva a rendir exámenes. En contraparte, los recursos destinados a
  conexión (\texttt{internet\_quality}) extrañamente reportan que caer a
  niveles bajos reduce marginalmente la asistencia en un \textbf{4.8\%}.
\item
  \textbf{Productividad contra Expectativas:} El sorpresivo regresor del
  rol laboral (\texttt{part\_time\_job}) dictamina firmemente que el
  hecho empírico de trabajar a tiempo parcial empuja la probabilidad de
  rendir el examen hacia arriba en un \textbf{7.0\%} en relación con un
  estudiante que no trabaja.
\end{enumerate}

    \begin{center}\rule{0.5\linewidth}{0.5pt}\end{center}

\subsection{Pregunta 4: Modelo Logit}\label{pregunta-4-modelo-logit}

Ahora estimaremos el modelo \textbf{Logit}. Este funciona de manera
bastante similar al Probit pero usando una distribución logística,
conocida por manejar eventos con extremas propabiliades un poco mejor.

    \begin{tcolorbox}[breakable, size=fbox, boxrule=1pt, pad at break*=1mm,colback=cellbackground, colframe=cellborder]
\prompt{In}{incolor}{25}{\boxspacing}
\begin{Verbatim}[commandchars=\\\{\}]
\PY{c+c1}{\PYZsh{} 1. Ajustar modelo Logit}
\PY{c+c1}{\PYZsh{}modelo\PYZus{}logit = sm.Logit(Y, X\PYZus{}con\PYZus{}constante\PYZus{}final).fit()}
\PY{n}{modelo\PYZus{}logit} \PY{o}{=} \PY{n}{sm}\PY{o}{.}\PY{n}{Logit}\PY{p}{(}\PY{n}{Y}\PY{p}{,} \PY{n}{X\PYZus{}con\PYZus{}constante\PYZus{}final}\PY{p}{)}\PY{o}{.}\PY{n}{fit}\PY{p}{(}\PY{n}{cov\PYZus{}type}\PY{o}{=}\PY{l+s+s1}{\PYZsq{}}\PY{l+s+s1}{HC1}\PY{l+s+s1}{\PYZsq{}}\PY{p}{)}
\PY{c+c1}{\PYZsh{} 2. Resumen del Logit}
\PY{n+nb}{print}\PY{p}{(}\PY{n}{modelo\PYZus{}logit}\PY{o}{.}\PY{n}{summary}\PY{p}{(}\PY{p}{)}\PY{p}{)}

\PY{c+c1}{\PYZsh{} 3. Efectos Marginales}
\PY{n}{margeff\PYZus{}logit} \PY{o}{=} \PY{n}{modelo\PYZus{}logit}\PY{o}{.}\PY{n}{get\PYZus{}margeff}\PY{p}{(}\PY{p}{)}
\PY{n+nb}{print}\PY{p}{(}\PY{l+s+s2}{\PYZdq{}}\PY{l+s+se}{\PYZbs{}n}\PY{l+s+s2}{\PYZhy{}\PYZhy{}\PYZhy{} EFECTOS MARGINALES DEL MODELO LOGIT \PYZhy{}\PYZhy{}\PYZhy{}}\PY{l+s+s2}{\PYZdq{}}\PY{p}{)}
\PY{n+nb}{print}\PY{p}{(}\PY{n}{margeff\PYZus{}logit}\PY{o}{.}\PY{n}{summary}\PY{p}{(}\PY{p}{)}\PY{p}{)}
\end{Verbatim}
\end{tcolorbox}

    \begin{Verbatim}[commandchars=\\\{\}]
Optimization terminated successfully.
         Current function value: 0.382173
         Iterations 7
                           Logit Regression Results
==============================================================================
Dep. Variable:          rindio\_examen   No. Observations:                 5621
Model:                          Logit   Df Residuals:                     5611
Method:                           MLE   Df Model:                            9
Date:                Tue, 21 Apr 2026   Pseudo R-squ.:                  0.2450
Time:                        16:21:06   Log-Likelihood:                -2148.2
converged:                       True   LL-Null:                       -2845.4
Covariance Type:                  HC1   LLR p-value:                1.196e-294
================================================================================
======
                         coef    std err          z      P>|z|      [0.025
0.975]
--------------------------------------------------------------------------------
------
const                  1.1581      0.321      3.607      0.000       0.529
1.787
academic\_level        -0.5700      0.051    -11.232      0.000      -0.669
-0.471
sleep\_hours            0.1395      0.036      3.840      0.000       0.068
0.211
screen\_time\_hours     -0.0634      0.016     -3.876      0.000      -0.095
-0.031
part\_time\_job          0.5737      0.080      7.213      0.000       0.418
0.730
upcoming\_deadline     -1.1495      0.089    -12.976      0.000      -1.323
-0.976
internet\_quality      -0.4093      0.049     -8.318      0.000      -0.506
-0.313
focus\_index            0.0405      0.007      5.863      0.000       0.027
0.054
productivity\_score     0.0531      0.004     12.049      0.000       0.044
0.062
gender\_male           -0.8122      0.079    -10.320      0.000      -0.966
-0.658
================================================================================
======

--- EFECTOS MARGINALES DEL MODELO LOGIT ---
        Logit Marginal Effects
=====================================
Dep. Variable:          rindio\_examen
Method:                          dydx
At:                           overall
================================================================================
======
                        dy/dx    std err          z      P>|z|      [0.025
0.975]
--------------------------------------------------------------------------------
------
academic\_level        -0.0695      0.006    -11.857      0.000      -0.081
-0.058
sleep\_hours            0.0170      0.004      3.853      0.000       0.008
0.026
screen\_time\_hours     -0.0077      0.002     -3.895      0.000      -0.012
-0.004
part\_time\_job          0.0700      0.010      7.297      0.000       0.051
0.089
upcoming\_deadline     -0.1402      0.010    -13.999      0.000      -0.160
-0.121
internet\_quality      -0.0499      0.006     -8.643      0.000      -0.061
-0.039
focus\_index            0.0049      0.001      5.937      0.000       0.003
0.007
productivity\_score     0.0065      0.001     12.631      0.000       0.005
0.007
gender\_male           -0.0991      0.009    -10.778      0.000      -0.117
-0.081
================================================================================
======
    \end{Verbatim}

    \subsubsection{Interpretación de Resultados del Modelo Logit (Efectos
Marginales)}\label{interpretaciuxf3n-de-resultados-del-modelo-logit-efectos-marginales}

Procedemos a interpretar sus Efectos Marginales para comparar la
probabilidad de éxito empírico:

\begin{enumerate}
\def\labelenumi{\arabic{enumi}.}
\tightlist
\item
  \textbf{Variables de Naturaleza Académica y Psicológica:} La
  estructura de \texttt{academic\_level} reporta mediante su efecto
  marginal (\texttt{dy/dx}) que, al ascender un peldaño en complejidad
  educacional estandarizada, la probabilidad predictiva de que el
  estudiante se ausente se estira y debilita progresivamente un
  \textbf{6.5\%}.
\item
  \textbf{Características Demográficas:} El factor dicotómico aislado de
  ser hombre (\texttt{gender\_male}) reduce marginalmente en un
  \textbf{9.8\%} las chances de presentarse formalmente en un examen.
\item
  \textbf{Consistencia de la Distribución:} Al mapear las varianzas,
  todo el set demuestra ser matemáticamente robusto con Valores-P
  iguales a 0. Las probabilidades predictivas que nos entregan estos
  efectos marginales trazan resultados prácticamente iguales a los
  derivados del cálculo de una distribución estricta de modelo Normal
  (Probit), confirmando que esta muestra es absolutamente apta y
  confiable para resoluciones teóricas que eviten las fallas lineales
  del MCO estándar.
\end{enumerate}

    \begin{center}\rule{0.5\linewidth}{0.5pt}\end{center}

\subsubsection{Pregunta 5: Análisis Comparativo y Selección del Modelo
Óptimo}\label{pregunta-5-anuxe1lisis-comparativo-y-selecciuxf3n-del-modelo-uxf3ptimo}

    \begin{tcolorbox}[breakable, size=fbox, boxrule=1pt, pad at break*=1mm,colback=cellbackground, colframe=cellborder]
\prompt{In}{incolor}{26}{\boxspacing}
\begin{Verbatim}[commandchars=\\\{\}]
\PY{k+kn}{from}\PY{+w}{ }\PY{n+nn}{stargazer}\PY{n+nn}{.}\PY{n+nn}{stargazer}\PY{+w}{ }\PY{k+kn}{import} \PY{n}{Stargazer}
\PY{k+kn}{from}\PY{+w}{ }\PY{n+nn}{IPython}\PY{n+nn}{.}\PY{n+nn}{core}\PY{n+nn}{.}\PY{n+nn}{display}\PY{+w}{ }\PY{k+kn}{import} \PY{n}{HTML}

\PY{n}{stargazer} \PY{o}{=} \PY{n}{Stargazer}\PY{p}{(}\PY{p}{[}\PY{n}{modelo\PYZus{}mco\PYZus{}final}\PY{p}{,} \PY{n}{modelo\PYZus{}logit}\PY{p}{,} \PY{n}{modelo\PYZus{}probit}\PY{p}{]}\PY{p}{)}

\PY{n}{stargazer}\PY{o}{.}\PY{n}{title}\PY{p}{(}\PY{l+s+s2}{\PYZdq{}}\PY{l+s+s2}{Comparación de Modelos: OLS, Logit y Probit}\PY{l+s+s2}{\PYZdq{}}\PY{p}{)}
\PY{n}{stargazer}\PY{o}{.}\PY{n}{custom\PYZus{}columns}\PY{p}{(}\PY{p}{[}\PY{l+s+s2}{\PYZdq{}}\PY{l+s+s2}{Modelo OLS}\PY{l+s+s2}{\PYZdq{}}\PY{p}{,} \PY{l+s+s2}{\PYZdq{}}\PY{l+s+s2}{Modelo Logit}\PY{l+s+s2}{\PYZdq{}}\PY{p}{,} \PY{l+s+s2}{\PYZdq{}}\PY{l+s+s2}{Modelo Probit}\PY{l+s+s2}{\PYZdq{}}\PY{p}{]}\PY{p}{,} \PY{p}{[}\PY{l+m+mi}{1}\PY{p}{,} \PY{l+m+mi}{1}\PY{p}{,} \PY{l+m+mi}{1}\PY{p}{]}\PY{p}{)}
\PY{n}{stargazer}\PY{o}{.}\PY{n}{add\PYZus{}custom\PYZus{}notes}\PY{p}{(}\PY{p}{[}\PY{l+s+s2}{\PYZdq{}}\PY{l+s+s2}{Nota: Todos los modelos utilizan errores estándar robustos (HC1).}\PY{l+s+s2}{\PYZdq{}}\PY{p}{]}\PY{p}{)}

\PY{n}{HTML}\PY{p}{(}\PY{n}{stargazer}\PY{o}{.}\PY{n}{render\PYZus{}html}\PY{p}{(}\PY{p}{)}\PY{p}{)}
\end{Verbatim}
\end{tcolorbox}

            \begin{tcolorbox}[breakable, size=fbox, boxrule=.5pt, pad at break*=1mm, opacityfill=0]
\prompt{Out}{outcolor}{26}{\boxspacing}
\begin{Verbatim}[commandchars=\\\{\}]
<IPython.core.display.HTML object>
\end{Verbatim}
\end{tcolorbox}
        
    \begin{tcolorbox}[breakable, size=fbox, boxrule=1pt, pad at break*=1mm,colback=cellbackground, colframe=cellborder]
\prompt{In}{incolor}{27}{\boxspacing}
\begin{Verbatim}[commandchars=\\\{\}]
\PY{c+c1}{\PYZsh{} Para comparar los modelos, podemos usar el AIC (Akaike Information Criterion) y el BIC (Bayesian Information Criterion) que son medidas de la calidad del modelo penalizadas por el número de parámetros. El modelo con el menor AIC y BIC es generalmente considerado el mejor.}
\PY{n+nb}{print}\PY{p}{(}\PY{l+s+s2}{\PYZdq{}}\PY{l+s+s2}{Comparativa de Modelos:}\PY{l+s+s2}{\PYZdq{}}\PY{p}{)}
\PY{n+nb}{print}\PY{p}{(}\PY{l+s+sa}{f}\PY{l+s+s2}{\PYZdq{}}\PY{l+s+s2}{Modelo OLS: AIC = }\PY{l+s+si}{\PYZob{}}\PY{n}{modelo\PYZus{}mco\PYZus{}final}\PY{o}{.}\PY{n}{aic}\PY{l+s+si}{:}\PY{l+s+s2}{.2f}\PY{l+s+si}{\PYZcb{}}\PY{l+s+s2}{, BIC = }\PY{l+s+si}{\PYZob{}}\PY{n}{modelo\PYZus{}mco\PYZus{}final}\PY{o}{.}\PY{n}{bic}\PY{l+s+si}{:}\PY{l+s+s2}{.2f}\PY{l+s+si}{\PYZcb{}}\PY{l+s+s2}{\PYZdq{}}\PY{p}{)}
\PY{n+nb}{print}\PY{p}{(}\PY{l+s+sa}{f}\PY{l+s+s2}{\PYZdq{}}\PY{l+s+s2}{Modelo Logit: AIC = }\PY{l+s+si}{\PYZob{}}\PY{n}{modelo\PYZus{}logit}\PY{o}{.}\PY{n}{aic}\PY{l+s+si}{:}\PY{l+s+s2}{.2f}\PY{l+s+si}{\PYZcb{}}\PY{l+s+s2}{, BIC = }\PY{l+s+si}{\PYZob{}}\PY{n}{modelo\PYZus{}logit}\PY{o}{.}\PY{n}{bic}\PY{l+s+si}{:}\PY{l+s+s2}{.2f}\PY{l+s+si}{\PYZcb{}}\PY{l+s+s2}{\PYZdq{}}\PY{p}{)}
\PY{n+nb}{print}\PY{p}{(}\PY{l+s+sa}{f}\PY{l+s+s2}{\PYZdq{}}\PY{l+s+s2}{Modelo Probit: AIC = }\PY{l+s+si}{\PYZob{}}\PY{n}{modelo\PYZus{}probit}\PY{o}{.}\PY{n}{aic}\PY{l+s+si}{:}\PY{l+s+s2}{.2f}\PY{l+s+si}{\PYZcb{}}\PY{l+s+s2}{, BIC = }\PY{l+s+si}{\PYZob{}}\PY{n}{modelo\PYZus{}probit}\PY{o}{.}\PY{n}{bic}\PY{l+s+si}{:}\PY{l+s+s2}{.2f}\PY{l+s+si}{\PYZcb{}}\PY{l+s+s2}{\PYZdq{}}\PY{p}{)}
\end{Verbatim}
\end{tcolorbox}

    \begin{Verbatim}[commandchars=\\\{\}]
Comparativa de Modelos:
Modelo OLS: AIC = 4470.91, BIC = 4537.25
Modelo Logit: AIC = 4316.39, BIC = 4382.73
Modelo Probit: AIC = 4345.82, BIC = 4412.17
    \end{Verbatim}

    Tras haber estimado y tabulado (vía Stargazer) los tres modelos
econométricos (Mínimos Cuadrados Ordinarios, Probit y Logit) para
explicar la probabilidad binomial de que un estudiante rinda su examen
final, podemos realizar las siguientes observaciones sobre su
comportamiento:

\textbf{1. Comparativa de Resultados y Explicación de sus Diferencias:}
Al contrastar los modelos, se puede apreciar que los \textbf{Efectos
Marginales} (\texttt{dy/dx}) extraídos de los modelos de respuesta
binaria no-lineales (Probit y Logit) convergen hacia magnitudes
predictivas \emph{casi idénticas} a las ponderaciones (\(\beta\))
estimadas linealmente a través de MCO. Por ejemplo, enfrentar una
entrega inminente (\texttt{upcoming\_deadline}) se asocia a una
penalización aproximadamente del \textbf{12.8\%} a lo largo de los 3
modelos. Las ligeras diferencias que persisten entre los modelos nacen
estrictamente de sus supuestos funcionales: * El \textbf{MCO} asume una
asociación recta e infinita, ignorando la cota probabilística e
ignorando la heterocedasticidad inherente en respuestas de \(0\) y
\(1\). * El \textbf{Probit} fuerza a los errores a distribuirse
siguiendo la campana Gaussiana Estándar Acumulada. * El \textbf{Logit}
modela sobre una Función de Distribución Logística, dotando a su
``campana'' estadística de colas levemente más pesadas, permitiéndole
lidiar con observaciones probabilísticas extremas de manera diferencial
frente al Probit.

\textbf{2. Decisión del Modelo Más Adecuado:} \textbf{el modelo Logit
representa la herramienta estadísticamente óptima para responder a
nuestra pregunta de investigación.}

El MCO debe ser descartado del análisis binomial puesto que las rectas
lineales permiten la proyección de ``probabilidades'' imposibles
(mayores al 100\% o menores al 0\%). A la hora de coronar al Logit sobre
el Probit, \textbf{la decisión es confirmada en su totalidad por
nuestros estadísticos de ajuste cruzado (Akaike y Bayesiano)}: el MCO
poseía el peor ajuste (AIC = \(4470.91\), BIC = \(4537.25\)), pero al
iterar, el clasificador Logístico demostró ser la red empírica mejor
ensamblada para la base de datos, entregando la cuota de error más baja
(AIC = \(4316.39\), BIC = \(4382.73\)), batiendo con ventaja a sus dos
contrapartes. Del mismo modo, el Pseudo \(\R^2\) más grande es del
modelo logit, explicando el mayor porcentaje de la realidad del ¿Por
qué? los individuos de la muestra deciden asistir o no al examen.

\textbf{3. Variables Robustas a la Especificación:} Toda variable que
sobrevivió al filtro iterativo demostró ser \textbf{robusta a la
especificación econométrica}. Identificadores empíricos como mantener un
trabajo \emph{part-time} (\texttt{part\_time\_job}), cursar grados
formales jerárquicamente mayores (\texttt{academic\_level}), y el tiempo
de descanso (\texttt{sleep\_hours}), no sufrieron alteraciones ni en su
signo direccional, ni en el peso de su rigor estadístico
(\(P>|z| \approx 0.000\)) independiente de si se probaban en funciones
MCO, Logit o Probit.

    \begin{center}\rule{0.5\linewidth}{0.5pt}\end{center}

\subsubsection{Pregunta 6}\label{pregunta-6}

    \begin{tcolorbox}[breakable, size=fbox, boxrule=1pt, pad at break*=1mm,colback=cellbackground, colframe=cellborder]
\prompt{In}{incolor}{28}{\boxspacing}
\begin{Verbatim}[commandchars=\\\{\}]
\PY{c+c1}{\PYZsh{} 1. Filtramos SÓLO los estudiantes que SÍ rindieron el examen}
\PY{n}{df\PYZus{}rindieron} \PY{o}{=} \PY{n}{df}\PY{p}{[}\PY{n}{df}\PY{p}{[}\PY{l+s+s1}{\PYZsq{}}\PY{l+s+s1}{rindio\PYZus{}examen}\PY{l+s+s1}{\PYZsq{}}\PY{p}{]} \PY{o}{==} \PY{l+m+mi}{1}\PY{p}{]}\PY{o}{.}\PY{n}{copy}\PY{p}{(}\PY{p}{)}

\PY{c+c1}{\PYZsh{} 2. Tu nuevo Target (Y) es la NOTA, no la asistencia}
\PY{n}{Y\PYZus{}poisson} \PY{o}{=} \PY{n}{df\PYZus{}rindieron}\PY{p}{[}\PY{l+s+s1}{\PYZsq{}}\PY{l+s+s1}{exam\PYZus{}score}\PY{l+s+s1}{\PYZsq{}}\PY{p}{]}
\end{Verbatim}
\end{tcolorbox}

    \begin{tcolorbox}[breakable, size=fbox, boxrule=1pt, pad at break*=1mm,colback=cellbackground, colframe=cellborder]
\prompt{In}{incolor}{29}{\boxspacing}
\begin{Verbatim}[commandchars=\\\{\}]
\PY{n}{columnas\PYZus{}a\PYZus{}excluir\PYZus{}poisson} \PY{o}{=} \PY{p}{[}
    \PY{l+s+s1}{\PYZsq{}}\PY{l+s+s1}{student\PYZus{}id}\PY{l+s+s1}{\PYZsq{}}\PY{p}{,} \PY{l+s+s1}{\PYZsq{}}\PY{l+s+s1}{exam\PYZus{}score}\PY{l+s+s1}{\PYZsq{}}\PY{p}{,} \PY{l+s+s1}{\PYZsq{}}\PY{l+s+s1}{rindio\PYZus{}examen}\PY{l+s+s1}{\PYZsq{}}\PY{p}{,}
    \PY{l+s+s1}{\PYZsq{}}\PY{l+s+s1}{age}\PY{l+s+s1}{\PYZsq{}}\PY{p}{,} \PY{l+s+s1}{\PYZsq{}}\PY{l+s+s1}{study\PYZus{}hours}\PY{l+s+s1}{\PYZsq{}}\PY{p}{,} \PY{l+s+s1}{\PYZsq{}}\PY{l+s+s1}{self\PYZus{}study\PYZus{}hours}\PY{l+s+s1}{\PYZsq{}}\PY{p}{,} \PY{l+s+s1}{\PYZsq{}}\PY{l+s+s1}{online\PYZus{}classes\PYZus{}hours}\PY{l+s+s1}{\PYZsq{}}\PY{p}{,}
    \PY{l+s+s1}{\PYZsq{}}\PY{l+s+s1}{social\PYZus{}media\PYZus{}hours}\PY{l+s+s1}{\PYZsq{}}\PY{p}{,} \PY{l+s+s1}{\PYZsq{}}\PY{l+s+s1}{gaming\PYZus{}hours}\PY{l+s+s1}{\PYZsq{}}\PY{p}{,} \PY{l+s+s1}{\PYZsq{}}\PY{l+s+s1}{exercise\PYZus{}minutes}\PY{l+s+s1}{\PYZsq{}}\PY{p}{,} 
    \PY{l+s+s1}{\PYZsq{}}\PY{l+s+s1}{burnout\PYZus{}level}\PY{l+s+s1}{\PYZsq{}}\PY{p}{,} \PY{l+s+s1}{\PYZsq{}}\PY{l+s+s1}{gender\PYZus{}other}\PY{l+s+s1}{\PYZsq{}}\PY{p}{,} \PY{l+s+s1}{\PYZsq{}}\PY{l+s+s1}{caffeine\PYZus{}intake\PYZus{}mg}\PY{l+s+s1}{\PYZsq{}}\PY{p}{,} \PY{l+s+s1}{\PYZsq{}}\PY{l+s+s1}{mental\PYZus{}health\PYZus{}score}\PY{l+s+s1}{\PYZsq{}}
\PY{p}{]}
\PY{n}{X\PYZus{}poisson} \PY{o}{=} \PY{n}{df\PYZus{}rindieron}\PY{o}{.}\PY{n}{drop}\PY{p}{(}\PY{n}{columns}\PY{o}{=}\PY{n}{columnas\PYZus{}a\PYZus{}excluir\PYZus{}poisson}\PY{p}{)}
\PY{n}{X\PYZus{}con\PYZus{}constante\PYZus{}poisson} \PY{o}{=} \PY{n}{sm}\PY{o}{.}\PY{n}{add\PYZus{}constant}\PY{p}{(}\PY{n}{X\PYZus{}poisson}\PY{p}{)}

\PY{n}{modelo\PYZus{}poisson} \PY{o}{=} \PY{n}{sm}\PY{o}{.}\PY{n}{Poisson}\PY{p}{(}\PY{n}{Y\PYZus{}poisson}\PY{p}{,} \PY{n}{X\PYZus{}con\PYZus{}constante\PYZus{}poisson}\PY{p}{)}\PY{o}{.}\PY{n}{fit}\PY{p}{(}\PY{n}{cov\PYZus{}type}\PY{o}{=}\PY{l+s+s1}{\PYZsq{}}\PY{l+s+s1}{HC1}\PY{l+s+s1}{\PYZsq{}}\PY{p}{)}
\PY{n+nb}{print}\PY{p}{(}\PY{n}{modelo\PYZus{}poisson}\PY{o}{.}\PY{n}{summary}\PY{p}{(}\PY{p}{)}\PY{p}{)}
\end{Verbatim}
\end{tcolorbox}

    \begin{Verbatim}[commandchars=\\\{\}]
Optimization terminated successfully.
         Current function value: 3.297099
         Iterations 5
                          Poisson Regression Results
==============================================================================
Dep. Variable:             exam\_score   No. Observations:                 4473
Model:                        Poisson   Df Residuals:                     4463
Method:                           MLE   Df Model:                            9
Date:                Tue, 21 Apr 2026   Pseudo R-squ.:                  0.4111
Time:                        16:21:06   Log-Likelihood:                -14748.
converged:                       True   LL-Null:                       -25045.
Covariance Type:                  HC1   LLR p-value:                     0.000
================================================================================
======
                         coef    std err          z      P>|z|      [0.025
0.975]
--------------------------------------------------------------------------------
------
const                  1.4899      0.038     39.338      0.000       1.416
1.564
academic\_level         0.0024      0.005      0.457      0.648      -0.008
0.013
sleep\_hours            0.0495      0.004     13.440      0.000       0.042
0.057
screen\_time\_hours     -0.0159      0.002     -9.031      0.000      -0.019
-0.012
part\_time\_job         -0.0873      0.009     -9.905      0.000      -0.105
-0.070
upcoming\_deadline     -0.1345      0.009    -15.341      0.000      -0.152
-0.117
internet\_quality       0.0017      0.005      0.315      0.753      -0.009
0.012
focus\_index            0.0123      0.001     16.903      0.000       0.011
0.014
productivity\_score     0.0227      0.000     50.204      0.000       0.022
0.024
gender\_male            0.0081      0.009      0.938      0.348      -0.009
0.025
================================================================================
======
    \end{Verbatim}

    \begin{tcolorbox}[breakable, size=fbox, boxrule=1pt, pad at break*=1mm,colback=cellbackground, colframe=cellborder]
\prompt{In}{incolor}{30}{\boxspacing}
\begin{Verbatim}[commandchars=\\\{\}]
\PY{n}{columnas\PYZus{}a\PYZus{}excluir\PYZus{}poisson} \PY{o}{=} \PY{p}{[}
    \PY{l+s+s1}{\PYZsq{}}\PY{l+s+s1}{student\PYZus{}id}\PY{l+s+s1}{\PYZsq{}}\PY{p}{,} \PY{l+s+s1}{\PYZsq{}}\PY{l+s+s1}{exam\PYZus{}score}\PY{l+s+s1}{\PYZsq{}}\PY{p}{,} \PY{l+s+s1}{\PYZsq{}}\PY{l+s+s1}{rindio\PYZus{}examen}\PY{l+s+s1}{\PYZsq{}}\PY{p}{,}
    \PY{l+s+s1}{\PYZsq{}}\PY{l+s+s1}{age}\PY{l+s+s1}{\PYZsq{}}\PY{p}{,} \PY{l+s+s1}{\PYZsq{}}\PY{l+s+s1}{study\PYZus{}hours}\PY{l+s+s1}{\PYZsq{}}\PY{p}{,} \PY{l+s+s1}{\PYZsq{}}\PY{l+s+s1}{self\PYZus{}study\PYZus{}hours}\PY{l+s+s1}{\PYZsq{}}\PY{p}{,} \PY{l+s+s1}{\PYZsq{}}\PY{l+s+s1}{online\PYZus{}classes\PYZus{}hours}\PY{l+s+s1}{\PYZsq{}}\PY{p}{,}
    \PY{l+s+s1}{\PYZsq{}}\PY{l+s+s1}{social\PYZus{}media\PYZus{}hours}\PY{l+s+s1}{\PYZsq{}}\PY{p}{,} \PY{l+s+s1}{\PYZsq{}}\PY{l+s+s1}{gaming\PYZus{}hours}\PY{l+s+s1}{\PYZsq{}}\PY{p}{,} \PY{l+s+s1}{\PYZsq{}}\PY{l+s+s1}{exercise\PYZus{}minutes}\PY{l+s+s1}{\PYZsq{}}\PY{p}{,} 
    \PY{l+s+s1}{\PYZsq{}}\PY{l+s+s1}{burnout\PYZus{}level}\PY{l+s+s1}{\PYZsq{}}\PY{p}{,} \PY{l+s+s1}{\PYZsq{}}\PY{l+s+s1}{gender\PYZus{}other}\PY{l+s+s1}{\PYZsq{}}\PY{p}{,} \PY{l+s+s1}{\PYZsq{}}\PY{l+s+s1}{caffeine\PYZus{}intake\PYZus{}mg}\PY{l+s+s1}{\PYZsq{}}\PY{p}{,} \PY{l+s+s1}{\PYZsq{}}\PY{l+s+s1}{mental\PYZus{}health\PYZus{}score}\PY{l+s+s1}{\PYZsq{}}\PY{p}{,}
    \PY{l+s+s1}{\PYZsq{}}\PY{l+s+s1}{academic\PYZus{}level}\PY{l+s+s1}{\PYZsq{}}\PY{p}{,} \PY{l+s+s1}{\PYZsq{}}\PY{l+s+s1}{internet\PYZus{}quality}\PY{l+s+s1}{\PYZsq{}}\PY{p}{,} \PY{l+s+s1}{\PYZsq{}}\PY{l+s+s1}{gender\PYZus{}male}\PY{l+s+s1}{\PYZsq{}}                
\PY{p}{]}
\PY{n}{X\PYZus{}poisson} \PY{o}{=} \PY{n}{df\PYZus{}rindieron}\PY{o}{.}\PY{n}{drop}\PY{p}{(}\PY{n}{columns}\PY{o}{=}\PY{n}{columnas\PYZus{}a\PYZus{}excluir\PYZus{}poisson}\PY{p}{)}
\PY{n}{X\PYZus{}con\PYZus{}constante\PYZus{}poisson} \PY{o}{=} \PY{n}{sm}\PY{o}{.}\PY{n}{add\PYZus{}constant}\PY{p}{(}\PY{n}{X\PYZus{}poisson}\PY{p}{)}

\PY{n}{modelo\PYZus{}poisson} \PY{o}{=} \PY{n}{sm}\PY{o}{.}\PY{n}{Poisson}\PY{p}{(}\PY{n}{Y\PYZus{}poisson}\PY{p}{,} \PY{n}{X\PYZus{}con\PYZus{}constante\PYZus{}poisson}\PY{p}{)}\PY{o}{.}\PY{n}{fit}\PY{p}{(}\PY{n}{cov\PYZus{}type}\PY{o}{=}\PY{l+s+s1}{\PYZsq{}}\PY{l+s+s1}{HC1}\PY{l+s+s1}{\PYZsq{}}\PY{p}{)}
\PY{n+nb}{print}\PY{p}{(}\PY{n}{modelo\PYZus{}poisson}\PY{o}{.}\PY{n}{summary}\PY{p}{(}\PY{p}{)}\PY{p}{)}
\end{Verbatim}
\end{tcolorbox}

    \begin{Verbatim}[commandchars=\\\{\}]
Optimization terminated successfully.
         Current function value: 3.297330
         Iterations 5
                          Poisson Regression Results
==============================================================================
Dep. Variable:             exam\_score   No. Observations:                 4473
Model:                        Poisson   Df Residuals:                     4466
Method:                           MLE   Df Model:                            6
Date:                Tue, 21 Apr 2026   Pseudo R-squ.:                  0.4111
Time:                        16:21:07   Log-Likelihood:                -14749.
converged:                       True   LL-Null:                       -25045.
Covariance Type:                  HC1   LLR p-value:                     0.000
================================================================================
======
                         coef    std err          z      P>|z|      [0.025
0.975]
--------------------------------------------------------------------------------
------
const                  1.5032      0.033     45.186      0.000       1.438
1.568
sleep\_hours            0.0493      0.004     13.444      0.000       0.042
0.056
screen\_time\_hours     -0.0159      0.002     -9.062      0.000      -0.019
-0.012
part\_time\_job         -0.0873      0.009     -9.898      0.000      -0.105
-0.070
upcoming\_deadline     -0.1339      0.009    -15.255      0.000      -0.151
-0.117
focus\_index            0.0123      0.001     16.907      0.000       0.011
0.014
productivity\_score     0.0227      0.000     50.149      0.000       0.022
0.024
================================================================================
======
    \end{Verbatim}

    \textbf{Interpretación de la estimación de Poisson (Pregunta 6):}

\begin{enumerate}
\def\labelenumi{\arabic{enumi}.}
\item
  \textbf{Variables Significativas}: En el modelo Poisson las variables
  de \texttt{sleep\_hours}, \texttt{screen\_time\_hours},
  \texttt{part\_time\_job}, \texttt{upcoming\_deadline},
  \texttt{focus\_index} y \texttt{productivity\_score} son fuertemente
  significativas (Valor-P \textless{} 0.05).
\item
  \textbf{Interpretación de los coeficientes}:

  \begin{itemize}
  \tightlist
  \item
    \textbf{\texttt{sleep\_hours}} (\(\beta = 0.0493\)): Por cada hora
    adicional de sueño, la nota esperada aumenta en aproximadamente un
    4.93\%.
  \item
    \textbf{\texttt{focus\_index}} (\(\beta = 0.0123\)) y
    \textbf{\texttt{productivity\_score}} (\(\beta = 0.0227\)): Por cada
    incremento en un punto del índice de concentración o de la
    productividad, la nota esperada aumenta alrededor de un 1.23\% y
    2.27\%, respectivamente.
  \item
    \textbf{\texttt{screen\_time\_hours}} (\(\beta = -0.0159\)): Cada
    hora adicional de pantalla disminuye el puntaje final esperado en un
    1.59\%.
  \item
    \textbf{\texttt{part\_time\_job}} (\(\beta = -0.0873\)): Los
    estudiantes que trabajan medio tiempo tienden a tener una nota un
    8.73\% más baja frente al resto de estudiantes.
  \item
    \textbf{\texttt{upcoming\_deadline}} (\(\beta = -0.1339\)): Aquellos
    sometidos a fechas límite que se aproximan, obtienen notas un
    13.39\% menores.
  \end{itemize}
\item
  \textbf{Evaluación global}: El Pseudo R-cuadrado es notablemente alto
  (0.4111), indicando un buen nivel explicativo de nuestro conjunto
  acotado de variables predictoras en la nota de los alumnos.
\end{enumerate}

    \begin{center}\rule{0.5\linewidth}{0.5pt}\end{center}

\subsubsection{Pregunta 7}\label{pregunta-7}

Determine si existe sobre dispersion en la data y posible valor optimo
de alpha para un modelo Binomial Negativa.

    \begin{tcolorbox}[breakable, size=fbox, boxrule=1pt, pad at break*=1mm,colback=cellbackground, colframe=cellborder]
\prompt{In}{incolor}{31}{\boxspacing}
\begin{Verbatim}[commandchars=\\\{\}]
\PY{c+c1}{\PYZsh{} 1. Obtenemos las predicciones (lambda / mu) del modelo Poisson ajustado previamente}
\PY{n}{mu} \PY{o}{=} \PY{n}{modelo\PYZus{}poisson}\PY{o}{.}\PY{n}{predict}\PY{p}{(}\PY{p}{)}

\PY{c+c1}{\PYZsh{} 2. Calculamos la variable dependiente auxiliar descrita para el test de sobredispersión}
\PY{n}{aux\PYZus{}olsr} \PY{o}{=} \PY{p}{(}\PY{p}{(}\PY{n}{Y\PYZus{}poisson} \PY{o}{\PYZhy{}} \PY{n}{mu}\PY{p}{)}\PY{o}{*}\PY{o}{*}\PY{l+m+mi}{2} \PY{o}{\PYZhy{}} \PY{n}{Y\PYZus{}poisson}\PY{p}{)} \PY{o}{/} \PY{n}{mu}

\PY{c+c1}{\PYZsh{} 3. Corremos una regresión lineal (MCO/OLS) auxiliar, sin constante, usando mu como única variable explicativa}
\PY{n}{modelo\PYZus{}auxiliar} \PY{o}{=} \PY{n}{sm}\PY{o}{.}\PY{n}{OLS}\PY{p}{(}\PY{n}{aux\PYZus{}olsr}\PY{p}{,} \PY{n}{mu}\PY{p}{)}\PY{o}{.}\PY{n}{fit}\PY{p}{(}\PY{p}{)}
\PY{n+nb}{print}\PY{p}{(}\PY{l+s+s2}{\PYZdq{}}\PY{l+s+s2}{\PYZhy{}\PYZhy{}\PYZhy{} RESULTADOS TEST DE SOBREDISPERSIÓN \PYZhy{}\PYZhy{}\PYZhy{}}\PY{l+s+s2}{\PYZdq{}}\PY{p}{)}
\PY{n+nb}{print}\PY{p}{(}\PY{n}{modelo\PYZus{}auxiliar}\PY{o}{.}\PY{n}{summary}\PY{p}{(}\PY{p}{)}\PY{p}{)}

\PY{c+c1}{\PYZsh{} El estimador de alpha será la pendiente del modelo anterior.}
\PY{n}{alpha\PYZus{}estimado} \PY{o}{=} \PY{n}{modelo\PYZus{}auxiliar}\PY{o}{.}\PY{n}{params}\PY{o}{.}\PY{n}{iloc}\PY{p}{[}\PY{l+m+mi}{0}\PY{p}{]}  
\PY{n}{p\PYZus{}value\PYZus{}alpha} \PY{o}{=} \PY{n}{modelo\PYZus{}auxiliar}\PY{o}{.}\PY{n}{pvalues}\PY{o}{.}\PY{n}{iloc}\PY{p}{[}\PY{l+m+mi}{0}\PY{p}{]}

\PY{n+nb}{print}\PY{p}{(}\PY{l+s+s2}{\PYZdq{}}\PY{l+s+se}{\PYZbs{}n}\PY{l+s+s2}{===========================}\PY{l+s+s2}{\PYZdq{}}\PY{p}{)}
\PY{k}{if} \PY{n}{p\PYZus{}value\PYZus{}alpha} \PY{o}{\PYZlt{}} \PY{l+m+mf}{0.05} \PY{o+ow}{and} \PY{n}{alpha\PYZus{}estimado} \PY{o}{\PYZgt{}} \PY{l+m+mi}{0}\PY{p}{:}
    \PY{n+nb}{print}\PY{p}{(}\PY{l+s+sa}{f}\PY{l+s+s2}{\PYZdq{}}\PY{l+s+s2}{Existe EVIDENCIA ESTADÍSTICA firme de sobredispersión (Valor\PYZhy{}P=}\PY{l+s+si}{\PYZob{}}\PY{n}{p\PYZus{}value\PYZus{}alpha}\PY{l+s+si}{:}\PY{l+s+s2}{.4f}\PY{l+s+si}{\PYZcb{}}\PY{l+s+s2}{ \PYZlt{} 0.05).}\PY{l+s+s2}{\PYZdq{}}\PY{p}{)}
    \PY{n+nb}{print}\PY{p}{(}\PY{l+s+sa}{f}\PY{l+s+s2}{\PYZdq{}}\PY{l+s+s2}{Se procederá a utilizar el valor óptimo (alpha estimado) = }\PY{l+s+si}{\PYZob{}}\PY{n}{alpha\PYZus{}estimado}\PY{l+s+si}{:}\PY{l+s+s2}{.4f}\PY{l+s+si}{\PYZcb{}}\PY{l+s+s2}{ para el modelo Binomial Negativo.}\PY{l+s+s2}{\PYZdq{}}\PY{p}{)}
\PY{k}{else}\PY{p}{:}
    \PY{n+nb}{print}\PY{p}{(}\PY{l+s+s2}{\PYZdq{}}\PY{l+s+s2}{No hay evidencia de sobredispersión en la distribución, o la sobredispersión no es significativa (Poisson es suficiente).}\PY{l+s+s2}{\PYZdq{}}\PY{p}{)}
\end{Verbatim}
\end{tcolorbox}

    \begin{Verbatim}[commandchars=\\\{\}]
--- RESULTADOS TEST DE SOBREDISPERSIÓN ---
                                 OLS Regression Results
================================================================================
=======
Dep. Variable:             exam\_score   R-squared (uncentered):
0.044
Model:                            OLS   Adj. R-squared (uncentered):
0.044
Method:                 Least Squares   F-statistic:
207.5
Date:                Tue, 21 Apr 2026   Prob (F-statistic):
4.98e-46
Time:                        16:21:07   Log-Likelihood:
-11169.
No. Observations:                4473   AIC:
2.234e+04
Df Residuals:                    4472   BIC:
2.235e+04
Df Model:                           1
Covariance Type:            nonrobust
==============================================================================
                 coef    std err          t      P>|t|      [0.025      0.975]
------------------------------------------------------------------------------
x1             0.0274      0.002     14.406      0.000       0.024       0.031
==============================================================================
Omnibus:                     4769.744   Durbin-Watson:                   2.010
Prob(Omnibus):                  0.000   Jarque-Bera (JB):           564216.230
Skew:                           5.161   Prob(JB):                         0.00
Kurtosis:                      57.044   Cond. No.                         1.00
==============================================================================

Notes:
[1] R² is computed without centering (uncentered) since the model does not
contain a constant.
[2] Standard Errors assume that the covariance matrix of the errors is correctly
specified.

===========================
Existe EVIDENCIA ESTADÍSTICA firme de sobredispersión (Valor-P=0.0000 < 0.05).
Se procederá a utilizar el valor óptimo (alpha estimado) = 0.0274 para el modelo
Binomial Negativo.
    \end{Verbatim}

    \textbf{Conclusión del Test de Sobredispersión:}

Para verificar si la dispersión de nuestros datos es mayor a la media
proyectada, ejecutamos la regresión lineal auxiliar recomendada.

\begin{itemize}
\tightlist
\item
  \textbf{Estimación de Alpha (\(\alpha\)):} Obtuvimos un coeficiente de
  pendiente de \textbf{\texttt{0.0274}}. Al ser positivamente mayor a
  cero, confirma que la varianza penaliza el modelo. El valor-P del
  coeficiente es de \textbf{\texttt{0.000}} (\(p < 0.05\)), por lo que
  es estadísticamente significativa.
\end{itemize}

\textbf{Decisión:} Existe evidencia estadística firme que demuestra
\textbf{sobredispersión} en nuestra variable de interés. Debido a esto,
el modelo de Poisson es insuficiente y procederemos a utilizar un
\textbf{Modelo Binomial Negativo}, inyectándole nuestro valor óptimo de
\textbf{\(\alpha = 0.0274\)} para recalibrar correctamente los errores
estándar.

    \begin{center}\rule{0.5\linewidth}{0.5pt}\end{center}

\subsubsection{Pregunta 8}\label{pregunta-8}

Usando la informacion anterior, ejecute un modelo Binomial Negativa para
responder a la pregunta 6. Seleccione las variables dependientes a
incluir en el modelo final e interprete su significado.

    \begin{tcolorbox}[breakable, size=fbox, boxrule=1pt, pad at break*=1mm,colback=cellbackground, colframe=cellborder]
\prompt{In}{incolor}{32}{\boxspacing}
\begin{Verbatim}[commandchars=\\\{\}]
\PY{c+c1}{\PYZsh{} Utilizando las mismas variables robustas que hallamos en la pregunta 6 para el modelo final }
\PY{c+c1}{\PYZsh{} X\PYZus{}con\PYZus{}constante\PYZus{}poisson ya contiene la constante y las variables dependientes más significativas (ya depuradas).}

\PY{c+c1}{\PYZsh{} Instanciamos y ajustamos el modelo GLM indicando una familia Binomial Negativa y pasándole el alpha calculado.}
\PY{n}{modelo\PYZus{}nbinom} \PY{o}{=} \PY{n}{sm}\PY{o}{.}\PY{n}{GLM}\PY{p}{(}
    \PY{n}{Y\PYZus{}poisson}\PY{p}{,} 
    \PY{n}{X\PYZus{}con\PYZus{}constante\PYZus{}poisson}\PY{p}{,} 
    \PY{n}{family}\PY{o}{=}\PY{n}{sm}\PY{o}{.}\PY{n}{families}\PY{o}{.}\PY{n}{NegativeBinomial}\PY{p}{(}\PY{n}{alpha}\PY{o}{=}\PY{n}{alpha\PYZus{}estimado}\PY{p}{)}
\PY{p}{)}\PY{o}{.}\PY{n}{fit}\PY{p}{(}\PY{n}{cov\PYZus{}type}\PY{o}{=}\PY{l+s+s1}{\PYZsq{}}\PY{l+s+s1}{HC1}\PY{l+s+s1}{\PYZsq{}}\PY{p}{)}

\PY{n+nb}{print}\PY{p}{(}\PY{n}{modelo\PYZus{}nbinom}\PY{o}{.}\PY{n}{summary}\PY{p}{(}\PY{p}{)}\PY{p}{)}
\end{Verbatim}
\end{tcolorbox}

    \begin{Verbatim}[commandchars=\\\{\}]
                 Generalized Linear Model Regression Results
==============================================================================
Dep. Variable:             exam\_score   No. Observations:                 4473
Model:                            GLM   Df Residuals:                     4466
Model Family:        NegativeBinomial   Df Model:                            6
Link Function:                    Log   Scale:                          1.0000
Method:                          IRLS   Log-Likelihood:                -14380.
Date:                Tue, 21 Apr 2026   Deviance:                       5941.3
Time:                        16:21:07   Pearson chi2:                 5.48e+03
No. Iterations:                     6   Pseudo R-squ. (CS):             0.9503
Covariance Type:                  HC1
================================================================================
======
                         coef    std err          z      P>|z|      [0.025
0.975]
--------------------------------------------------------------------------------
------
const                  1.4107      0.035     40.425      0.000       1.342
1.479
sleep\_hours            0.0531      0.004     13.817      0.000       0.046
0.061
screen\_time\_hours     -0.0175      0.002     -9.623      0.000      -0.021
-0.014
part\_time\_job         -0.0942      0.009    -10.270      0.000      -0.112
-0.076
upcoming\_deadline     -0.1454      0.009    -15.839      0.000      -0.163
-0.127
focus\_index            0.0134      0.001     17.617      0.000       0.012
0.015
productivity\_score     0.0238      0.000     51.081      0.000       0.023
0.025
================================================================================
======
    \end{Verbatim}

    \textbf{Interpretación de Resultados - Regresión Binomial Negativa:}

Al estimar el modelo insertando el parámetro de sobredispersión hallado
previamente, podemos concluir lo siguiente:

\begin{enumerate}
\def\labelenumi{\arabic{enumi}.}
\item
  \textbf{Variables Significativas:} Las seis variables incorporadas
  (\texttt{sleep\_hours}, \texttt{screen\_time\_hours},
  \texttt{part\_time\_job}, \texttt{upcoming\_deadline},
  \texttt{focus\_index} y \texttt{productivity\_score}) continúan siendo
  fuertemente significativas, arrojando valores p prácticamente nulos
  (\(P>|z| = 0.000\)).
\item
  \textbf{Impacto Exponencial (Coeficientes):} Debido a que el modelo
  ocupa una función de enlace logarítmica (Log-Link), los coeficientes
  representan cambios multiplicativos porcentuales
  (\(\approx 100 \cdot \beta \%\)).

  \begin{itemize}
  \tightlist
  \item
    \textbf{Efectos Positivos:} Mejorar las horas de sueño
    (\texttt{+5.31\%} por hora adicional), el foco mental
    (\texttt{+1.34\%} por punto) y la productividad en clase
    (\texttt{+2.38\%} por punto) aumentan de forma consistente la nota
    esperada del alumno.
  \item
    \textbf{Efectos Negativos:} Inversamente, agentes estresores o de
    desconcentración penalizan la nota considerablemente. Aquellos que
    poseen entregas próximas bajan su rendimiento esperado en un
    monumental \texttt{14.5\%}, tener trabajo de medio tiempo lo reduce
    en \texttt{9.4\%}, y cada hora adicional gastada en pantallas lo
    baja un \texttt{1.75\%}.
  \end{itemize}
\end{enumerate}

\textbf{Conclusión:} Las variables robustas se ratifican. Al corregirse
los márgenes de error a través del relajamiento de la varianza (Binomial
Negativo), ganamos certeza respecto a que estos factores y sus
magnitudes representan determinantes fiables sobre las calificaciones de
los estudiantes.

    \begin{center}\rule{0.5\linewidth}{0.5pt}\end{center}

\subsubsection{Respuesta a Pregunta 9: Análisis
Comparativo}\label{respuesta-a-pregunta-9-anuxe1lisis-comparativo}

\textbf{1. Diferencias entre los modelos (Poisson vs.~Binomial
Negativo):} Los coeficientes estimados en ambos modelos son muy
similares numéricamente, pero difieren en sus \textbf{Errores Estándar}.
Esto sucede porque el modelo Poisson asume erróneamente que la media y
la varianza son iguales (equidispersión). El modelo Binomial Negativo,
en cambio, añade el parámetro \(\alpha\) para corregir la
sobredispersión, volviendo los errores estándar ligeramente más
flexibles y reales.

\textbf{2. Modelo más adecuado:} El modelo idóneo es el \textbf{Binomial
Negativo}. En la Pregunta 7 se demostró de manera estadísticamente
significativa (valor-P = 0.000) la existencia de un parámetro
\(\alpha > 0\) (\(\alpha = 0.0274\)). Usar el modelo Poisson en una
muestra con sobredispersión habría infravalorado la varianza e inflado
la significancia de las variables, generando conclusiones erróneas.

\textbf{3. Variables robustas a la especificación:} Aquellas variables
que conservaron consistentemente su significancia estadística
(\(P < 0.05\)) a lo largo de todos los modelos (desde MCO hasta Binomial
Negativo) demostraron ser estrictamente robustas. Estas fueron: *
\textbf{Positivas:} Horas de sueño (\texttt{sleep\_hours}), y los
índices cognitivos (\texttt{focus\_index},
\texttt{productivity\_score}). * \textbf{Negativas:} Tiempo de pantalla
(\texttt{screen\_time\_hours}), trabajar a tiempo parcial
(\texttt{part\_time\_job}) y tener entregas inminentes
(\texttt{upcoming\_deadline}).


    % Add a bibliography block to the postdoc
    
    
    
\end{document}
